\documentclass[12pt, openany]{book}

% ──────────────────────────────────────────────
% Packages
% ──────────────────────────────────────────────
\usepackage[margin=1in]{geometry}
\usepackage{amsmath, amssymb, amsthm}
\usepackage{mathtools}
\usepackage{bm}
\usepackage{tikz}
\usepackage{pgfplots}
\pgfplotsset{compat=1.18}
\usetikzlibrary{arrows.meta, decorations.markings, calc, patterns, positioning}
\usepackage{tcolorbox}
\tcbuselibrary{theorems, breakable, skins}
\usepackage{hyperref}
\usepackage{imakeidx}
\makeindex[intoc, title=Index]
\usepackage{enumitem}
\usepackage{fancyhdr}
\usepackage{titlesec}
\usepackage{epigraph}
\usepackage{xcolor}
\usepackage{booktabs}

% ──────────────────────────────────────────────
% Colors
% ──────────────────────────────────────────────
\definecolor{chapblue}{RGB}{0, 70, 130}
\definecolor{accentred}{RGB}{180, 40, 40}
\definecolor{softgray}{RGB}{245, 245, 245}
\definecolor{trajgreen}{RGB}{30, 130, 60}
\definecolor{darkgold}{RGB}{160, 120, 20}
\definecolor{purplehaze}{RGB}{100, 40, 140}

% ──────────────────────────────────────────────
% Theorem environments
% ──────────────────────────────────────────────
\newtcbtheorem[number within=chapter]{principle}{Principle}{
  colback=chapblue!5, colframe=chapblue!80,
  fonttitle=\bfseries, separator sign={.~},
  breakable
}{princ}

\newtcbtheorem[number within=chapter]{keyidea}{Key Idea}{
  colback=accentred!5, colframe=accentred!70,
  fonttitle=\bfseries, separator sign={.~},
  breakable
}{idea}

\newtcbtheorem[number within=chapter]{example}{Example}{
  colback=softgray, colframe=black!40,
  fonttitle=\bfseries, separator sign={.~},
  breakable
}{ex}

\newtcbtheorem[number within=chapter]{definition}{Definition}{
  colback=darkgold!5, colframe=darkgold!80,
  fonttitle=\bfseries, separator sign={.~},
  breakable
}{def}

\newtcbtheorem[number within=chapter]{remark}{Remark}{
  colback=purplehaze!5, colframe=purplehaze!60,
  fonttitle=\bfseries, separator sign={.~},
  breakable
}{rem}

\newtcbtheorem[number within=chapter]{warning}{Warning}{
  colback=accentred!5, colframe=accentred!80,
  fonttitle=\bfseries, separator sign={.~},
  breakable
}{warn}

\theoremstyle{plain}
\newtheorem{proposition}{Proposition}[chapter]
\newtheorem{theorem}{Theorem}[chapter]
\newtheorem{lemma}[theorem]{Lemma}
\newtheorem{corollary}[theorem]{Corollary}

% ──────────────────────────────────────────────
% Chapter formatting
% ──────────────────────────────────────────────
\titleformat{\chapter}[display]
  {\normalfont\huge\bfseries\color{chapblue}}
  {\chaptertitlename\ \thechapter}{20pt}{\Huge}
\titlespacing*{\chapter}{0pt}{-20pt}{30pt}

% ──────────────────────────────────────────────
% Header/Footer
% ──────────────────────────────────────────────
\pagestyle{fancy}
\fancyhf{}
\fancyhead[LE,RO]{\thepage}
\fancyhead[RE]{\textit{Control Is Motion}}
\fancyhead[LO]{\textit{\nouppercase{\leftmark}}}
\renewcommand{\headrulewidth}{0.4pt}
\setlength{\headheight}{14.5pt}

% ──────────────────────────────────────────────
% Custom commands
% ──────────────────────────────────────────────
\newcommand{\state}{\bm{x}}
\newcommand{\control}{\bm{u}}
\newcommand{\traj}{\bm{\gamma}}
\newcommand{\manifold}{\mathcal{M}}
\newcommand{\configspace}{\mathcal{Q}}
\newcommand{\statespace}{\mathcal{X}}
\newcommand{\controlspace}{\mathcal{U}}
\newcommand{\tube}{\mathcal{T}}
\newcommand{\Reals}{\mathbb{R}}
\newcommand{\dd}{\mathrm{d}}
\newcommand{\dt}{\,\dd t}
\newcommand{\norm}[1]{\left\lVert #1 \right\rVert}
\newcommand{\inner}[2]{\left\langle #1,\, #2 \right\rangle}

% Additional commands for specific chapters
\newcommand{\tangent}{\bm{T}}
\newcommand{\normal}{\bm{N}}
\newcommand{\binormal}{\bm{B}}
\newcommand{\frenet}{\bm{e}}
\newcommand{\ds}{\dd s}
\newcommand{\Sone}{S^1}
\newcommand{\Stwo}{S^2}
\newcommand{\Sthree}{S^3}
\newcommand{\SOthree}{\mathrm{SO}(3)}
\newcommand{\SEthree}{\mathrm{SE}(3)}
\newcommand{\SOtwo}{\mathrm{SO}(2)}
\newcommand{\SEtwo}{\mathrm{SE}(2)}
\newcommand{\Amat}{\bm{A}}
\newcommand{\Bmat}{\bm{B}}
\newcommand{\Cmat}{\bm{C}}
\newcommand{\Mmat}{\bm{M}}
\newcommand{\PhiMat}{\bm{\Phi}}
\newcommand{\gvec}{\bm{g}}
\newcommand{\Torus}{\mathbb{T}}
\newcommand{\monodromy}{\mathcal{M}}
\newcommand{\orbit}{\mathcal{O}}
\newcommand{\poincare}{\mathcal{P}}
\newcommand{\skewmat}[1]{\widehat{#1}}
\newcommand{\sothree}{\mathfrak{so}(3)}
\newcommand{\tangentbundle}{T\mathcal{M}}
\newcommand{\tangentspace}[1]{T_{#1}\mathcal{M}}
\newcommand{\transverse}{\perp}
% ──────────────────────────────────────────────
% Title Page
% ──────────────────────────────────────────────
\title{%
  {\Huge\bfseries\color{chapblue} Control Is Motion}\\[0.4cm]
  {\Large A New Framework for Nonlinear Systems\\
  That Move Through the World}\\[0.6cm]
  {\color{chapblue}\rule{\textwidth}{2pt}}
}
\author{{\Large D.\ Kraft}}
\date{}


\hypersetup{colorlinks=true, linkcolor=chapblue, filecolor=magenta, urlcolor=chapblue, citecolor=accentred}

\begin{document}

\maketitle
\thispagestyle{empty}
\cleardoublepage

% ──────────────────────────────────────────────
% Table of Contents
% ──────────────────────────────────────────────
\tableofcontents
\cleardoublepage

\chapter{Throwing Away the Target}

\epigraph{%
  A river does not flow toward the ocean because it knows where the ocean is.
  It flows because flowing is what rivers do.
}{--- \textup{Adapted from Lao Tzu}}

\vspace{0.5cm}

% ──────────────────────────────────────────────
\section{The Wrong Question}
% ──────────────────────────────────────────────

Open any classical control textbook and you will find, within the first few pages, a
picture like this: a system, an error signal, a reference point, and a controller
whose entire purpose is to drive that error to zero. The implicit promise is
seductive---\emph{give me a destination and I will take you there}.

But consider the golf swing.

A skilled golfer does not steer the clubhead toward the ball the way a thermostat
drives a room toward $72^\circ$F. There is no ``setpoint'' floating in space that
the club is chasing. Instead, the golfer's body executes a \emph{motion}---a
coordinated, whole-body trajectory through a high-dimensional configuration space.
The clubhead arrives at the ball not because it was pointed at the ball, but because
the entire motion was shaped so that, somewhere along its arc, the club happens to
pass through the ball's location at enormous speed.

\begin{keyidea}{The Central Thesis}{central}
  \textbf{Control is motion, not destination.}
  The fundamental object of nonlinear control is not a setpoint, not an equilibrium,
  and not a fixed target. It is a \emph{trajectory}---a curve through state space
  that the system is asked to follow, and around which disturbances are rejected.
  The quality of control is measured not by how close the system gets to a point,
  but by how faithfully it reproduces a desired motion.
\end{keyidea}

This is not merely a philosophical reframing. It changes the mathematics. It changes
what we optimize. It changes how we define stability. And it changes what ``success''
means in a controlled system.

The purpose of this book is to develop control theory from this trajectory-first
perspective, to show that the classical setpoint framework is a special (and often
misleading) case of a much richer theory, and to build the mathematical tools needed
to design controllers for systems whose purpose is to \emph{move}.

\vspace{0.3cm}

\begin{center}
\begin{tikzpicture}[scale=1.0]
  % Classical picture
  \begin{scope}[shift={(-4.5,0)}]
    \node[font=\bfseries\color{chapblue}] at (0, 2.8) {Classical View};
    \draw[-{Stealth[length=3mm]}, thick, gray] (-1.8, 0) -- (1.8, 0)
      node[right, font=\small] {$x_1$};
    \draw[-{Stealth[length=3mm]}, thick, gray] (0, -1.8) -- (0, 1.8)
      node[above, font=\small] {$x_2$};
    % Target
    \filldraw[accentred] (0.8, 0.6) circle (3pt)
      node[right, font=\small\bfseries, accentred] {~target $\state^*$};
    % Trajectories converging
    \foreach \a/\b in {-1.3/1.4, -1.5/-0.8, 1.4/-1.2, 1.0/1.5} {
      \draw[-{Stealth[length=2mm]}, thick, chapblue!70, opacity=0.7]
        (\a, \b) .. controls (0.2, 0.1) .. (0.75, 0.57);
    }
    \node[font=\small, text width=3.5cm, align=center] at (0, -2.5)
      {All trajectories converge\\to a single point.};
  \end{scope}

  % Our picture
  \begin{scope}[shift={(4.5,0)}]
    \node[font=\bfseries\color{trajgreen!80!black}] at (0, 2.8) {Trajectory View};
    \draw[-{Stealth[length=3mm]}, thick, gray] (-1.8, 0) -- (1.8, 0)
      node[right, font=\small] {$x_1$};
    \draw[-{Stealth[length=3mm]}, thick, gray] (0, -1.8) -- (0, 1.8)
      node[above, font=\small] {$x_2$};
    % Reference trajectory
    \draw[very thick, trajgreen, decoration={markings,
        mark=at position 0.3 with {\arrow{Stealth[length=3mm]}},
        mark=at position 0.7 with {\arrow{Stealth[length=3mm]}}
      }, postaction={decorate}]
      (-1.4, -1.2) .. controls (-0.5, 0.8) and (0.5, 1.6) .. (1.5, 0.3);
    \node[trajgreen!80!black, font=\small\bfseries] at (1.2, 1.3)
      {$\traj^*(t)$};
    % Tube around trajectory
    \draw[dashed, trajgreen!50, thick]
      (-1.2, -1.5) .. controls (-0.3, 0.5) and (0.7, 1.3) .. (1.7, 0.0);
    \draw[dashed, trajgreen!50, thick]
      (-1.6, -0.9) .. controls (-0.7, 1.1) and (0.3, 1.9) .. (1.3, 0.6);
    % Nearby trajectory
    \draw[thick, accentred!70, decoration={markings,
        mark=at position 0.5 with {\arrow{Stealth[length=2mm]}}
      }, postaction={decorate}]
      (-1.3, -1.0) .. controls (-0.4, 0.9) and (0.6, 1.5) .. (1.4, 0.4);
    \node[font=\small, text width=3.5cm, align=center] at (0, -2.5)
      {System tracks a moving\\reference through state space.};
  \end{scope}

  % Dividing line
  \draw[thick, gray!40] (0, -2.0) -- (0, 2.5);
\end{tikzpicture}
\end{center}

\vspace{0.3cm}

% ──────────────────────────────────────────────
\section{Why Setpoints Are a Special Case}
% ──────────────────────────────────────────────

Let us be precise about what we mean. Consider a nonlinear system
\begin{equation}\label{eq:nonlinear_sys}
  \dot{\state} = \bm{f}(\state, \control), \qquad
  \state \in \statespace \subseteq \Reals^n, \quad
  \control \in \controlspace \subseteq \Reals^m,
\end{equation}
where $\bm{f}$ is smooth and the state space $\statespace$ may carry additional
geometric structure (it may be a manifold, as we will see in Chapter~3).

In the classical framework, the control objective is:
\begin{quote}
  \emph{Find $\control(t)$ such that $\state(t) \to \state^*$ as $t \to \infty$.}
\end{quote}
The entire apparatus of Lyapunov stability, LQR, pole placement, and PID control
is built around this question. But notice what this objective assumes:

\begin{enumerate}[leftmargin=2cm, label=\textbf{A\arabic*.}]
  \item There exists a meaningful equilibrium $\state^*$ where $\bm{f}(\state^*, \control^*) = \bm{0}$.
  \item The system's purpose is to \emph{reach and remain at} this equilibrium.
  \item Performance is measured by how quickly and smoothly the error
        $\bm{e}(t) = \state(t) - \state^*$ decays.
\end{enumerate}

Now consider the systems this book is about:

\begin{itemize}[leftmargin=1.5cm]
  \item A golfer executing a swing---the body passes through a continuum of
        configurations and never ``arrives'' anywhere.
  \item A gymnast performing a floor routine---the motion \emph{is} the performance.
  \item A biped walking---there is no single posture that constitutes ``walking'';
        walking is a \emph{limit cycle}, a periodic orbit through state space.
  \item A robotic arm tracing a weld seam---the task is defined by the path, not the endpoint.
\end{itemize}

For all of these systems, assumptions \textbf{A1}--\textbf{A3} fail. The correct
control objective is fundamentally different:

\begin{principle}{The Trajectory Tracking Objective}{tracking}
  Given a reference trajectory $\traj^*: [0, T] \to \statespace$, find a feedback
  control law $\control = \bm{k}(\state, t)$ such that
  \begin{equation}\label{eq:tracking_obj}
    \norm{\state(t) - \traj^*(t)} \leq \varepsilon(t)
    \qquad \text{for all } t \in [0, T],
  \end{equation}
  where $\varepsilon(t)$ is a prescribed tracking tolerance (the ``tube width''),
  and the system is robust to disturbances $\bm{d}(t)$ satisfying
  $\norm{\bm{d}(t)} \leq \delta$.
\end{principle}

A setpoint is simply the degenerate case where $\traj^*(t) = \state^*$ for all $t$
and $T \to \infty$. The trajectory-tracking framework \emph{contains} the setpoint
framework, but not vice versa.

\begin{definition}{Trajectory Tube}{tube}
  Given a reference trajectory $\traj^*: [0,T] \to \statespace$ and a tube radius
  function $\varepsilon: [0,T] \to \Reals_{>0}$, the \textbf{trajectory tube} is the
  time-varying set
  \begin{equation}
    \tube(t) = \bigl\{\, \state \in \statespace : \norm{\state - \traj^*(t)} \leq \varepsilon(t) \,\bigr\}.
  \end{equation}
  A controller achieves \textbf{tube-invariance} if every solution starting in
  $\tube(0)$ remains in $\tube(t)$ for all $t \in [0, T]$.
\end{definition}

This definition is the trajectory-centric replacement for Lyapunov stability. Instead
of asking ``does the system converge to a point?'' we ask ``does the system stay
inside a moving tube?'' The tube can be tight where precision matters (e.g., at
ball impact in the golf swing) and loose where it does not (e.g., during the
backswing).

% ──────────────────────────────────────────────
\section{Anatomy of a Motion: The Golf Swing as Prototype}
% ──────────────────────────────────────────────

We will use the golf swing throughout this book as a running example---not because
golf is inherently more interesting than other motions, but because it compresses
an extraordinary number of control-theoretic challenges into roughly 1.2~seconds:

\begin{enumerate}[label=(\roman*)]
  \item \textbf{High dimensionality.} The human body has over 200 degrees of freedom.
        Even a simplified musculoskeletal model of the golf swing uses 20--40
        generalized coordinates.
  \item \textbf{Extreme nonlinearity.} Joint torques couple through the full rigid-body
        dynamics, including Coriolis and centrifugal terms that dominate at high
        angular velocities.
  \item \textbf{Underactuation.} Once the club is released in the downswing, the
        wrist hinge ``freewheels''---the golfer cannot directly control the club angle.
        The clubhead speed at impact comes from \emph{passive dynamics}, not
        active torque.
  \item \textbf{Time-criticality with flexible timing.} The swing must produce the
        right state (clubhead position, velocity, and orientation) at impact,
        but the \emph{exact} time of impact is not prescribed. The golfer
        is free to take the backswing slowly or quickly.
  \item \textbf{Repeatability over optimality.} A touring professional does not
        execute the theoretically optimal swing. They execute a swing they
        can \emph{repeat} 300 times in a tournament, under pressure, with
        minimal variance.
\end{enumerate}

\begin{example}{From Setpoint Thinking to Trajectory Thinking}{golf_shift}
  Suppose we model a simplified golf swing with generalized coordinates
  $\bm{q} = (\theta_{\text{torso}},\, \theta_{\text{shoulder}},\,
  \theta_{\text{arm}},\, \theta_{\text{wrist}})^\top$ and their velocities
  $\dot{\bm{q}}$, giving state $\state = (\bm{q}, \dot{\bm{q}}) \in \Reals^8$.

  \textbf{Setpoint approach} (wrong): Define a target state $\state^* =
  (\bm{q}_{\text{impact}}, \dot{\bm{q}}_{\text{impact}})$ and design a controller
  to drive $\state(t) \to \state^*$. This fails immediately: the system should not
  \emph{stop} at $\state^*$. Impact is a waypoint, not a destination. The club must
  be accelerating \emph{through} the ball, not decelerating toward it.

  \textbf{Trajectory approach} (right): Define a reference motion
  $\traj^*(s): [0, 1] \to \Reals^8$ parameterized by a phase variable $s$ (not
  clock time $t$). The controller's job is to ensure that the system traverses
  $\traj^*$ with the right sequencing and coordination, with a tight tube near the
  impact phase ($s \approx 0.7$) and a loose tube during setup ($s \approx 0$).
\end{example}

Point~(v) deserves special emphasis. Classical optimal control asks: ``What is
the minimum-time or minimum-energy trajectory?'' But this is the wrong question
for biological and athletic systems. The right question is:

\begin{principle}{The Repeatability Principle}{repeatability}
  The best trajectory is not the one that maximizes performance on a single trial,
  but the one that maximizes \emph{expected performance over many trials} in the
  presence of noise, fatigue, and uncertainty. Formally, we seek
  \begin{equation}\label{eq:repeatability}
    \traj^* = \arg\min_{\traj} \;
    \mathbb{E}\!\left[\, \int_0^T
      \ell\bigl(\state(t), \control(t)\bigr) \dt \,\right]
    + \lambda \, \mathrm{Var}\!\left[\, J(\traj, \bm{w}) \,\right],
  \end{equation}
  where $J$ is the performance cost, $\bm{w}$ represents stochastic disturbances
  (motor noise, perceptual error), and $\lambda > 0$ penalizes variance.
\end{principle}

This is why elite golfers do not all swing the same way. Each golfer finds a
trajectory that is locally optimal \emph{for their body and their noise characteristics}.
Robustness is not an afterthought bolted onto an optimal controller---it is the
primary design objective.

% ──────────────────────────────────────────────
\section{The Mathematical Landscape}
% ──────────────────────────────────────────────

To make the trajectory-first perspective rigorous, we need mathematical machinery
that most introductory control courses skip entirely. This section provides a roadmap
of the key concepts we will develop in subsequent chapters.

\subsection{Phase Variables and Time Reparameterization}

A crucial distinction separates \emph{what} the system does from \emph{when} it
does it. Consider two executions of the same golf swing: one with a slow backswing
and one with a fast backswing. The \emph{spatial path} through configuration space
may be identical, but the \emph{time parameterization} differs.

We formalize this by introducing a \textbf{phase variable} $s(t) \in [0, 1]$ that
monotonically increases from $0$ (start of motion) to $1$ (end of motion), with
dynamics
\begin{equation}\label{eq:phase}
  \dot{s} = \alpha(s, \state),
\end{equation}
where $\alpha > 0$ is the \emph{phase velocity}. The reference trajectory is defined
as a function of $s$, not $t$:
\begin{equation}
  \traj^* = \traj^*(s), \qquad s \in [0, 1].
\end{equation}
This decouples the \emph{shape} of the motion from its \emph{temporal execution}.
The controller has two separate tasks:
\begin{enumerate}[label=(\alph*)]
  \item \textbf{Path tracking}: keep $\state(t)$ close to the reference path
        $\traj^*(s(t))$, and
  \item \textbf{Timing control}: regulate $\dot{s}$ to achieve the desired
        speed profile along the path.
\end{enumerate}
This decomposition is impossible in the setpoint framework, where time and space
are fused together.

\subsection{Transverse Dynamics and Orbital Stability}

When we shift from point stability to trajectory stability, the relevant notion of
``closeness'' changes. We do not care about the distance from $\state(t)$ to
$\traj^*(t)$ at the same clock time---we care about the distance from $\state(t)$
to the \emph{nearest point on the trajectory}.

\begin{definition}{Transverse Distance}{transverse}
  Let $\traj^*: [0,1] \to \statespace$ be a smooth reference trajectory. The
  \textbf{transverse distance} from a state $\state$ to $\traj^*$ is
  \begin{equation}
    d_\perp(\state, \traj^*) = \min_{s \in [0,1]} \norm{\state - \traj^*(s)}.
  \end{equation}
  The minimizing argument $s^*(\state) = \arg\min_s \norm{\state - \traj^*(s)}$
  is the \textbf{projection phase}, giving the ``closest point on the reference.''
\end{definition}

This leads to a decomposition of the state into two components:
\begin{itemize}[leftmargin=1.5cm]
  \item The \textbf{tangential component}: where the system is \emph{along} the
        trajectory (measured by $s$).
  \item The \textbf{transverse component}: how far the system is \emph{from} the
        trajectory (measured by $d_\perp$).
\end{itemize}

\begin{center}
\begin{tikzpicture}[scale=1.2]
  % Reference trajectory
  \draw[very thick, trajgreen, decoration={markings,
      mark=at position 0.15 with {\arrow{Stealth[length=3mm]}},
      mark=at position 0.55 with {\arrow{Stealth[length=3mm]}},
      mark=at position 0.85 with {\arrow{Stealth[length=3mm]}}
    }, postaction={decorate}]
    (-3, -0.5) .. controls (-1, 1.5) and (1, 1.5) .. (3, -0.3);
  \node[trajgreen!70!black, font=\small\bfseries] at (3.3, 0.1) {$\traj^*(s)$};

  % Point on trajectory
  \coordinate (P) at (0.3, 1.45);
  \filldraw[trajgreen!70!black] (P) circle (2pt);
  \node[above left, font=\small, trajgreen!70!black] at (P) {$\traj^*(s^*)$};

  % Actual state
  \coordinate (X) at (0.9, 0.2);
  \filldraw[accentred] (X) circle (2pt);
  \node[below right, font=\small\bfseries, accentred] at (X) {$\state(t)$};

  % Transverse distance
  \draw[thick, accentred, dashed, {Stealth[length=2.5mm]}-{Stealth[length=2.5mm]}]
    (P) -- (X) node[midway, right, font=\small] {$d_\perp$};

  % Tangent vector
  \draw[-{Stealth[length=3mm]}, thick, chapblue]
    (P) -- ++(0.9, -0.15) node[above, font=\small] {tangential};

  % Phase annotation
  \draw[thin, gray, |->] (-3, -1.3) -- (3, -1.3)
    node[right, font=\small] {$s$};
  \node[font=\small, gray] at (-3, -1.6) {$0$};
  \node[font=\small, gray] at (3, -1.6) {$1$};
  \draw[thin, gray, dashed] (0.3, -1.3) -- (P);
  \filldraw[gray] (0.3, -1.3) circle (1.5pt);
  \node[font=\small, gray, below] at (0.3, -1.3) {$s^*$};
\end{tikzpicture}
\end{center}

\vspace{0.2cm}

A system is \textbf{orbitally stable} around $\traj^*$ if the transverse distance
$d_\perp(\state(t), \traj^*)$ remains small (or decays) even if the tangential
phase $s(t)$ drifts relative to the nominal timing. This is exactly the right notion
for the golf swing: we want the motion to follow the right \emph{path} even if the
golfer is slightly faster or slower than planned.

\subsection{Funnel Control and Time-Varying Tubes}

Combining the trajectory tube (Definition~\ref{def:tube}) with the
transverse/tangential decomposition leads to what we call \textbf{funnel control}:
the design of controllers that keep the system inside a prescribed time-varying tube
around the reference.

The tube width $\varepsilon(s)$ is a design parameter. For the golf swing:
\begin{equation}\label{eq:funnel_profile}
  \varepsilon(s) = \begin{cases}
    \varepsilon_{\text{loose}} & s \in [0,\, 0.3] \quad \text{(address \& backswing)}, \\[4pt]
    \varepsilon_{\text{loose}} - (\varepsilon_{\text{loose}} - \varepsilon_{\text{tight}})
      \dfrac{s - 0.3}{0.4} & s \in [0.3,\, 0.7] \quad \text{(transition \& downswing)}, \\[6pt]
    \varepsilon_{\text{tight}} & s \in [0.7,\, 0.8] \quad \text{(impact zone)}, \\[4pt]
    \text{relaxed} & s \in [0.8,\, 1.0] \quad \text{(follow-through)}.
  \end{cases}
\end{equation}

\begin{center}
\begin{tikzpicture}[xscale=8, yscale=2.5]
  % Axes
  \draw[-{Stealth[length=3mm]}, thick] (0, 0) -- (1.08, 0)
    node[right, font=\small] {phase $s$};
  \draw[-{Stealth[length=3mm]}, thick] (0, 0) -- (0, 1.3)
    node[above, font=\small] {tube width $\varepsilon(s)$};

  % Funnel profile
  \draw[very thick, accentred]
    (0, 1.0) -- (0.3, 1.0)
    -- (0.7, 0.15)
    -- (0.8, 0.15)
    (0.8, 0.15) .. controls (0.85, 0.3) and (0.9, 0.7) .. (1.0, 0.9);

  % Phase labels
  \node[font=\scriptsize, below, text width=1.3cm, align=center] at (0.15, -0.05)
    {backswing\\(loose)};
  \node[font=\scriptsize, below, text width=1.5cm, align=center] at (0.5, -0.05)
    {downswing\\(narrowing)};
  \node[font=\scriptsize, below, text width=1.2cm, align=center, accentred] at (0.75, -0.05)
    {\textbf{impact}\\(tight)};
  \node[font=\scriptsize, below, text width=1.5cm, align=center] at (0.92, -0.05)
    {follow-through\\(relaxed)};

  % Dashed guidelines
  \foreach \x in {0.3, 0.7, 0.8} {
    \draw[thin, gray, dashed] (\x, 0) -- (\x, 1.1);
  }
\end{tikzpicture}
\end{center}

This funnel profile encodes the physical insight that precision matters most at
impact and least during setup. The controller ``spends'' its control authority
where it matters, rather than trying to maintain uniform precision everywhere---a
lesson that classical setpoint control cannot teach.

% ──────────────────────────────────────────────
\section{What Changes When Control Is Motion}
% ──────────────────────────────────────────────

Adopting the trajectory-first perspective changes nearly every aspect of control
system design. Here we preview the major consequences, each of which will be
developed in detail in subsequent chapters.

\subsection{Stability Is Not Convergence---It Is Confinement}

In the setpoint framework, stability means convergence: the state approaches the
equilibrium and stays near it. In the trajectory framework, stability means
\emph{confinement}: the state remains inside the moving tube.

\begin{definition}{Trajectory Stability}{traj_stability}
  A reference trajectory $\traj^*$ is \textbf{stable} under the closed-loop
  dynamics $\dot{\state} = \bm{f}(\state, \bm{k}(\state, t))$ if for every
  $\varepsilon > 0$ there exists $\delta > 0$ such that
  \begin{equation}
    d_\perp(\state(0), \traj^*) < \delta \implies
    d_\perp(\state(t), \traj^*) < \varepsilon \quad \forall\, t \in [0, T].
  \end{equation}
  It is \textbf{asymptotically trajectory-stable} if, additionally,
  $d_\perp(\state(t), \traj^*) \to 0$ as the system evolves.
\end{definition}

Note the subtle but critical difference: we use the transverse distance $d_\perp$,
not the time-indexed error $\norm{\state(t) - \traj^*(t)}$. This allows the system
to be ahead of or behind schedule without being considered ``unstable.''

\subsection{The Cost Function Is a Functional, Not a Function}

In setpoint control, the cost is typically a function of the error $\bm{e}(t)$:
\begin{equation}
  J_{\text{setpoint}} = \int_0^\infty \bigl(\bm{e}^\top \bm{Q}\, \bm{e}
  + \control^\top \bm{R}\, \control\bigr) \dt.
\end{equation}
In trajectory control, the cost is a \textbf{functional} that operates on the
entire motion:
\begin{equation}\label{eq:traj_cost}
  J_{\text{traj}}[\state(\cdot), \control(\cdot)] =
  \underbrace{\int_0^T w_\perp(s)\, d_\perp^2\bigl(\state(t), \traj^*\bigr) \dt}_{%
    \text{tracking fidelity (transverse)}}
  + \underbrace{\int_0^T w_\parallel(s)\, (\dot{s} - \dot{s}_{\text{ref}})^2 \dt}_{%
    \text{timing fidelity (tangential)}}
  + \underbrace{\int_0^T \control^\top \bm{R}(s)\, \control \dt}_{%
    \text{effort}},
\end{equation}
where the weighting matrices $w_\perp(s)$, $w_\parallel(s)$, and $\bm{R}(s)$ are
\emph{phase-dependent}. This is what allows us to demand high precision at impact
and low effort during the backswing.

\subsection{Feedforward Is Not Optional---It Is Primary}

In the setpoint framework, feedforward is a ``nice to have'' that reduces transient
error. In the trajectory framework, feedforward is the \emph{primary} control
action. The reference trajectory $\traj^*(t)$ implicitly defines a feedforward
control signal
\begin{equation}
  \control_{\text{ff}}(t) = \bm{f}^{-1}\!\left(\dot{\traj}^*(t),\, \traj^*(t)\right),
\end{equation}
where $\bm{f}^{-1}$ denotes the inverse dynamics. The feedback controller handles
only the \emph{deviations} from this feedforward signal:
\begin{equation}\label{eq:ff_fb_decomp}
  \control(t) = \underbrace{\control_{\text{ff}}(t)}_{\text{executes the motion}}
  + \underbrace{\control_{\text{fb}}(\state(t), t)}_{\text{corrects deviations}}.
\end{equation}
For the golf swing, the feedforward component accounts for approximately 90\% of
the total joint torques. The feedback controller provides the remaining 10\% of
correction. A control design that treats everything as feedback is asking the
feedback loop to do ten times more work than necessary.

\subsection{Underactuation Creates Geometry, Not Just Constraints}

Many moving systems are underactuated: they have fewer control inputs than degrees
of freedom. In the setpoint framework, underactuation is an obstacle---it means we
cannot independently control all states. In the trajectory framework, underactuation
creates \emph{geometric structure} that we can exploit.

The set of states reachable from a given configuration through passive dynamics
forms a submanifold of the state space. A well-designed trajectory \emph{rides along}
this manifold, using the passive dynamics to generate motion that active torques
alone could not produce. This is how the golf swing works: the wrist hinge releases
passively, and the centrifugal acceleration of the downswing \emph{automatically}
squares the clubface through the double pendulum effect.

\begin{proposition}
  For a mechanical system with $n$ degrees of freedom and $m < n$ actuators, the
  set of dynamically feasible trajectories lies on a submanifold of dimension
  at most $2m$ in the $2n$-dimensional state space. Effective trajectory design
  requires understanding the geometry of this submanifold.
\end{proposition}

We will make this precise in Chapter~5 using the language of constraint manifolds
and their null-space structures.

% ──────────────────────────────────────────────
\section{A Roadmap for the Book}
% ──────────────────────────────────────────────

The remainder of this book develops the trajectory-first theory in a sequence of
increasingly sophisticated layers:

\begin{description}[leftmargin=1.5cm, labelwidth=1.2cm, font=\bfseries\color{chapblue}]
  \item[Ch.~2] \textbf{Curves in State Space.}
    The geometry of trajectories: parameterization, arc length, curvature, and torsion
    in high-dimensional state spaces. Frenet--Serret frames for nonlinear systems.

  \item[Ch.~3] \textbf{Configuration Manifolds.}
    Why state space is often not $\Reals^n$: Lie groups, joint spaces, and the
    geometry of articulated bodies. The golf swing lives on $SO(3)^k$, not $\Reals^k$.

  \item[Ch.~4] \textbf{Orbital Stability and Transverse Linearization.}
    Rigorous definitions of trajectory stability. Linearization transverse to the
    orbit. Floquet theory for periodic motions. Moving Poincar\'e sections.

  \item[Ch.~5] \textbf{Underactuation and Passive Dynamics.}
    Constraint Jacobians, null-space control, and the exploitation of passive
    mechanical coupling. The double pendulum and the physics of the wrist release.

  \item[Ch.~6] \textbf{Trajectory Optimization.}
    Direct collocation, shooting methods, and pseudospectral methods for finding
    optimal reference trajectories. The interplay between optimality and
    repeatability.

  \item[Ch.~7] \textbf{Funnel Synthesis.}
    Designing time-varying tubes with guaranteed invariance. Sums-of-squares
    programming for funnel computation. Robust funnels under model uncertainty.

  \item[Ch.~8] \textbf{Phase-Variable Control}~\cite{WesterveltGrizzle2007}.
    Decoupling path shape from timing. Virtual constraints and hybrid zero dynamics
    for walking and running. Central pattern generators as phase oscillators.

  \item[Ch.~9] \textbf{Stochastic Trajectories and Motor Variability.}
    Signal-dependent noise in biological actuators. The minimum-variance theory of human movement~\cite{TodorovJordan2002}. Why optimal trajectories are smooth~\cite{FlashHogan1985}.

  \item[Ch.~10] \textbf{Learning to Move.}
    Iterative learning control, policy search, and trajectory libraries. How to
    improve a motion over repeated trials without losing stability.

  \item[Ch.~11] \textbf{Case Study: The Complete Golf Swing.}
    Full application of the entire theory to a 15-DOF musculoskeletal model of the
    golf swing, from trajectory optimization through funnel synthesis to
    stochastic robustness analysis.
\end{description}

% ──────────────────────────────────────────────
\section{A New Language for Control}
% ──────────────────────────────────────────────

Before proceeding, let us establish the conceptual vocabulary that will replace
the classical language throughout this book:

\begin{center}
\renewcommand{\arraystretch}{1.4}
\begin{tabular}{@{} l @{\qquad$\longrightarrow$\qquad} l @{}}
  \textbf{\color{gray}Classical term} & \textbf{\color{chapblue}Trajectory term} \\
  \hline
  Setpoint / reference & Reference trajectory $\traj^*(s)$ \\
  Error $\bm{e}(t) = \state - \state^*$ & Transverse deviation $d_\perp(\state, \traj^*)$ \\
  Lyapunov stability & Orbital / trajectory stability \\
  Convergence to equilibrium & Confinement to tube \\
  Regulation & Path tracking + timing control \\
  Step response & Trajectory tracking response \\
  Steady-state error & Persistent transverse offset \\
  Gain margin & Funnel robustness margin \\
  Pole placement & Transverse eigenvalue assignment \\
\end{tabular}
\end{center}

\vspace{0.5cm}

% ──────────────────────────────────────────────
\section{Exercises}
% ──────────────────────────────────────────────

\begin{enumerate}[label=\textbf{1.\arabic*.}, leftmargin=1.5cm]

\item \textbf{(Conceptual.)} Consider a simple pendulum $\ddot{\theta} + \sin\theta = u$.
  \begin{enumerate}[label=(\alph*)]
    \item Identify the equilibria and classify their stability.
    \item Now consider the task of making the pendulum execute a full revolution
          (pumping up from rest to continuous spinning). Explain why this task
          \emph{cannot} be formulated as setpoint regulation.
    \item Sketch the reference trajectory in the $(\theta, \dot{\theta})$ phase
          plane for one complete revolution. What is the natural phase variable?
  \end{enumerate}

\item \textbf{(Mathematical.)} For the system $\dot{x}_1 = x_2$,\;
  $\dot{x}_2 = -x_1 + u$, consider the reference trajectory
  $\traj^*(t) = (\cos t,\, -\sin t)^\top$ (a unit circle in state space).
  \begin{enumerate}[label=(\alph*)]
    \item Verify that $u_{\text{ff}}(t) = 0$ produces exactly this trajectory.
    \item Compute the transverse linearization about $\traj^*$.
    \item Design a feedback law $u_{\text{fb}}$ that achieves asymptotic orbital
          stability.
  \end{enumerate}

\item \textbf{(Design.)} A simplified 2-DOF golf swing model has coordinates
  $\bm{q} = (\theta_{\text{arm}}, \theta_{\text{club}})$ with dynamics
  \[
    \bm{M}(\bm{q})\ddot{\bm{q}} + \bm{C}(\bm{q}, \dot{\bm{q}})\dot{\bm{q}}
    + \bm{g}(\bm{q}) = \bm{B}\,\control,
    \qquad \bm{B} = \begin{pmatrix} 1 \\ 0 \end{pmatrix}.
  \]
  This system is underactuated: only $\theta_{\text{arm}}$ is directly actuated.
  \begin{enumerate}[label=(\alph*)]
    \item Explain qualitatively how the club can accelerate through passive dynamics
          even though no torque is applied at the wrist.
    \item Define a phase variable $s$ and sketch a funnel profile $\varepsilon(s)$
          appropriate for this system.
    \item Discuss why a setpoint controller applied at impact ($s = 0.7$) would
          produce the wrong motion.
  \end{enumerate}

\item \textbf{(Philosophical.)} A thermostat, an autopilot, and a dancer's body
  are all ``controlled systems.'' Classify each according to whether its primary
  control objective is (i)~setpoint regulation, (ii)~trajectory tracking, or
  (iii)~something else entirely. Justify your classification.

\end{enumerate}

\vspace{1cm}
\begin{center}
  \rule{0.5\textwidth}{0.4pt}\\[0.5cm]
  {\small\itshape
  The river does not arrive at the ocean and stop.\\
  It becomes the ocean.\\
  The swing does not arrive at impact and stop.\\
  It becomes the follow-through.\\[0.3cm]
  Control is motion.}
\end{center}

\chapter{Curves in State Space}

\epigraph{%
  The shortest distance between two points is a straight line,
  but no interesting motion has ever been straight.
}{}

\vspace{0.5cm}

% ══════════════════════════════════════════════
\section{Why Geometry Comes First}
% ══════════════════════════════════════════════

In Chapter~1 we argued that control is fundamentally about \emph{motion}---the
faithful execution of a trajectory through state space. Before we can stabilize
a trajectory, optimize a trajectory, or confine the system to a tube around a
trajectory, we need to understand what a trajectory \emph{is} as a mathematical
object.

This chapter develops the differential geometry of curves in $\Reals^n$, adapted
to the needs of control theory. The central insight is that a trajectory has
intrinsic geometric properties---curvature, torsion, and a natural moving
frame---that exist independently of how fast the system moves along the curve.
These intrinsic properties determine how difficult the trajectory is to track,
how much control authority is needed to stay on course, and where the system
is most vulnerable to disturbances.

The reader may wonder: \emph{why not simply work in coordinates?} The answer is
that coordinate-based descriptions mix up properties of the \emph{curve} with
properties of the \emph{parameterization}. A golf swing executed slowly and the
same swing executed quickly trace the same curve in configuration space, but
their coordinate descriptions $\bm{q}(t)$ are entirely different functions. The
geometric tools we develop here allow us to separate what is intrinsic to the
motion from what is merely a choice of clock.

\begin{keyidea}{Intrinsic vs.\ Extrinsic}{intrinsic}
  A trajectory has two kinds of properties:
  \begin{itemize}[leftmargin=1cm]
    \item \textbf{Intrinsic}: properties of the curve itself---its shape,
      curvature, how it bends and twists through state space. These are
      independent of parameterization.
    \item \textbf{Extrinsic}: properties that depend on how the curve is
      traversed---speed, acceleration, timing. These depend on the
      parameterization.
  \end{itemize}
  Good control design separates these cleanly. The \emph{trajectory planner}
  designs the intrinsic shape. The \emph{timing law} specifies the extrinsic
  speed profile. The \emph{tracking controller} maintains both.
\end{keyidea}

% ══════════════════════════════════════════════
\section{Parameterized Curves}
% ══════════════════════════════════════════════

\begin{definition}{Parameterized Curve}{param_curve}
  A \textbf{parameterized curve} in $\Reals^n$ is a smooth map
  \begin{equation}
    \traj: I \to \Reals^n, \qquad \lambda \mapsto \traj(\lambda),
  \end{equation}
  where $I \subseteq \Reals$ is an interval. The parameter $\lambda$ may be time $t$,
  arc length $s$, a phase variable $\sigma$, or any other monotone scalar.
  The curve is \textbf{regular} if $\traj'(\lambda) \neq \bm{0}$ for all
  $\lambda \in I$, i.e., it never ``stops'' (though the system traversing it may
  slow down).
\end{definition}

The critical point is that many different parameterizations describe the same
geometric curve. If $\traj(\lambda)$ is a parameterized curve and
$\lambda = \phi(\mu)$ is a smooth, monotonically increasing reparameterization,
then $\tilde{\traj}(\mu) = \traj(\phi(\mu))$ traces the same set of points in
$\Reals^n$, but at different ``speeds.''

\begin{example}{The Circle: Three Parameterizations}{circle_param}
  Consider the unit circle in $\Reals^2$. Three natural parameterizations:
  \begin{align}
    \traj_1(t) &= (\cos t,\; \sin t), \quad t \in [0, 2\pi]
    && \text{(constant angular speed)}, \\
    \traj_2(t) &= (\cos t^2,\; \sin t^2), \quad t \in [0, \sqrt{2\pi}]
    && \text{(accelerating)}, \\
    \traj_3(s) &= (\cos s,\; \sin s), \quad s \in [0, 2\pi]
    && \text{(arc-length, same as $\traj_1$ here)}.
  \end{align}
  All three trace the same circle. The \emph{curve} is the same; the
  \emph{parameterizations} differ. Quantities like curvature are properties of
  the curve, not the parameterization.
\end{example}

\subsection{Velocity, Speed, and the Tangent Vector}

Given a parameterized curve $\traj(\lambda)$, the \textbf{velocity vector} is
\begin{equation}
  \bm{v}(\lambda) = \frac{\dd\traj}{\dd\lambda} = \traj'(\lambda).
\end{equation}
The \textbf{speed} is its magnitude:
\begin{equation}
  v(\lambda) = \norm{\traj'(\lambda)}.
\end{equation}
The \textbf{unit tangent vector} is the velocity normalized to unit length:
\begin{equation}\label{eq:unit_tangent}
  \tangent(\lambda) = \frac{\traj'(\lambda)}{\norm{\traj'(\lambda)}}
  = \frac{\traj'(\lambda)}{v(\lambda)}.
\end{equation}
The tangent vector $\tangent$ tells us the \emph{direction} of the curve at each
point---a purely geometric quantity. The speed $v$ tells us how \emph{fast} we
traverse it---a parameterization-dependent quantity. This separation is the first
instance of the intrinsic/extrinsic decomposition.

% ══════════════════════════════════════════════
\section{Arc Length: The Natural Parameter}
% ══════════════════════════════════════════════

Among all possible parameterizations, one is geometrically privileged: the
\textbf{arc-length parameterization}, in which the parameter measures distance
traveled along the curve.

\begin{definition}{Arc Length}{arc_length}
  The \textbf{arc length} of a curve $\traj(\lambda)$ from $\lambda = a$ to
  $\lambda = b$ is
  \begin{equation}\label{eq:arc_length}
    L = \int_a^b \norm{\traj'(\lambda)} \, \dd\lambda
    = \int_a^b v(\lambda) \, \dd\lambda.
  \end{equation}
  The \textbf{arc-length function} $s(\lambda)$ measures the distance from the
  starting point:
  \begin{equation}
    s(\lambda) = \int_a^{\lambda} \norm{\traj'(\mu)} \, \dd\mu,
    \qquad \text{so that } \frac{\dd s}{\dd\lambda} = v(\lambda).
  \end{equation}
\end{definition}

When a curve is parameterized by arc length, the speed is identically 1:
$\norm{\traj'(s)} = 1$ for all $s$. This means the tangent vector and the velocity
vector coincide: $\tangent(s) = \traj'(s)$. All geometric information is in the
direction of the derivative; none is hidden in its magnitude.

\begin{principle}{Arc Length as the Canonical Clock}{arc_clock}
  Arc-length parameterization strips away all timing information and reveals
  the pure geometry of the curve. In control terms, it answers the question:
  \emph{what does the motion look like if we forget how fast it goes?}

  For a golf swing, the arc-length parameterized trajectory $\traj(s)$
  describes the spatial \emph{path} of the body through configuration space.
  The timing law $s(t)$---how the golfer moves along this path in real
  time---is a separate design choice. A slow practice swing and a full-speed
  competition swing follow the same $\traj(s)$ with different $s(t)$.
\end{principle}

\begin{remark}{Arc Length in High Dimensions}{high_dim_arc}
  In $\Reals^n$ with the Euclidean metric, arc length is defined exactly as
  above. When the state space carries a non-Euclidean metric (as it will when
  we work on configuration manifolds in Chapter~3), the formula
  \eqref{eq:arc_length} generalizes to
  \begin{equation}
    L = \int_a^b \sqrt{\traj'(\lambda)^\top \bm{G}(\traj(\lambda))\,
    \traj'(\lambda)} \, \dd\lambda,
  \end{equation}
  where $\bm{G}(\state)$ is the Riemannian metric tensor (for mechanical systems,
  this is the mass-inertia matrix). This means ``distance in configuration
  space'' naturally accounts for how massive different parts of the body are.
  A degree of freedom connected to a heavy segment contributes more to arc
  length than one connected to a light segment---exactly as physical intuition
  suggests.
\end{remark}

% ══════════════════════════════════════════════
\section{Curvature: How Curves Bend}
% ══════════════════════════════════════════════

The most important intrinsic property of a curve is its \textbf{curvature}---a
measure of how rapidly the tangent direction changes as we move along the curve.

\begin{definition}{Curvature}{curvature}
  Let $\traj(s)$ be a curve parameterized by arc length in $\Reals^n$. The
  \textbf{curvature} at each point is
  \begin{equation}\label{eq:curvature_def}
    \kappa(s) = \norm{\traj''(s)} = \norm{\frac{\dd\tangent}{\dd s}}.
  \end{equation}
  Since $\norm{\tangent} = 1$, the derivative $\tangent'(s)$ is orthogonal to
  $\tangent(s)$, and $\kappa$ measures the rate of turning per unit distance.
  The curvature is always non-negative and is zero if and only if the curve is
  locally a straight line.
\end{definition}

The reciprocal $R(s) = 1/\kappa(s)$ is the \textbf{radius of curvature}---the
radius of the circle that best approximates the curve at that point. High curvature
means a tight bend; low curvature means a gentle sweep.

\subsection{Curvature in Arbitrary Parameterization}

In practice, we rarely work in arc-length parameterization. Given a curve
$\traj(t)$ parameterized by time (or any other parameter), the curvature can
be computed directly:

\begin{theorem}[Curvature Formula]\label{thm:curvature_formula}
  For a regular curve $\traj(t)$ in $\Reals^n$, the curvature is
  \begin{equation}\label{eq:curvature_general}
    \kappa(t) = \frac{\norm{\traj'(t) \times_n \traj''(t)}}{\norm{\traj'(t)}^3},
  \end{equation}
  where for $n = 2$, $\norm{\traj' \times_2 \traj''} = |x'y'' - y'x''|$,
  and for $n = 3$, $\times_3$ is the usual cross product. For general $n$,
  this generalizes via
  \begin{equation}\label{eq:curvature_general_n}
    \kappa(t) = \frac{\sqrt{\norm{\traj'}^2 \norm{\traj''}^2
    - (\traj' \cdot \traj'')^2}}{\norm{\traj'}^3}.
  \end{equation}
\end{theorem}

\begin{proof}
  Write $\traj(t) = \traj(s(t))$ where $s$ is arc length. Then
  $\traj' = \dot{s}\, \tangent$ and
  $\traj'' = \ddot{s}\, \tangent + \dot{s}^2\, \kappa\, \normal$,
  where $\normal$ is the unit normal (defined below). Since
  $\tangent \perp \normal$, we have
  $\norm{\traj'}^2\norm{\traj''}^2 - (\traj'\cdot\traj'')^2
  = \dot{s}^2(\ddot{s}^2 + \dot{s}^4\kappa^2) - \dot{s}^2\ddot{s}^2
  = \dot{s}^6\kappa^2$.
  Dividing by $\norm{\traj'}^6 = \dot{s}^6$ and taking the square root yields
  the result.
\end{proof}

\begin{example}{Curvature of an Elliptical Trajectory}{ellipse_curv}
  Consider the ellipse $\traj(t) = (a\cos t,\; b\sin t)$ with $a > b > 0$. Then
  \begin{align}
    \traj'(t) &= (-a\sin t,\; b\cos t), \\
    \traj''(t) &= (-a\cos t,\; -b\sin t).
  \end{align}
  Applying the 2D curvature formula:
  \begin{equation}
    \kappa(t) = \frac{|(-a\sin t)(-b\sin t) - (b\cos t)(-a\cos t)|}
    {(a^2\sin^2 t + b^2\cos^2 t)^{3/2}}
    = \frac{ab}{(a^2\sin^2 t + b^2\cos^2 t)^{3/2}}.
  \end{equation}
  The curvature is \emph{highest at the ends of the minor axis}
  ($t = 0, \pi$, where $\kappa = a/b^2$) and \emph{lowest at the ends of the
  major axis} ($t = \pi/2, 3\pi/2$, where $\kappa = b/a^2$). The tighter the
  bend, the higher the curvature---and, as we will see, the more control
  authority is needed to track it.
\end{example}

\subsection{Why Curvature Matters for Control}

Curvature is not merely a geometric curiosity. It has direct implications for
tracking difficulty.

\begin{principle}{The Curvature--Authority Principle}{curv_authority}
  Consider a system tracking a reference trajectory $\traj^*(s)$ at speed $v$.
  The centripetal acceleration required to follow the curve is
  \begin{equation}\label{eq:centripetal}
    a_\perp = \kappa\, v^2.
  \end{equation}
  This means:
  \begin{enumerate}[label=(\roman*)]
    \item High curvature demands high lateral control authority.
    \item The required authority grows as the \emph{square} of the speed.
    \item A trajectory that is easy to track slowly may be impossible to track
          at high speed.
  \end{enumerate}
  For the golf swing, this explains why the transition from backswing to
  downswing is so critical: the clubhead path has high curvature precisely
  where the speed is increasing rapidly.
\end{principle}

\begin{center}
\begin{tikzpicture}[scale=1.0]
  % Low curvature region
  \begin{scope}[shift={(-4,0)}]
    \node[font=\bfseries\small, chapblue] at (0, 2.2) {Low curvature};
    \draw[very thick, trajgreen, decoration={markings,
        mark=at position 0.5 with {\arrow{Stealth[length=3mm]}}
      }, postaction={decorate}]
      (-2, 0) .. controls (-0.5, 0.3) and (0.5, 0.3) .. (2, 0);
    \draw[<->, thin, gray] (-0.2, 0.2) -- (-0.2, 1.5)
      node[midway, left, font=\scriptsize] {$R$ large};
    % Osculating circle (large)
    \draw[thin, dashed, gray] (0, 3.5) circle (3.3);
    \node[font=\scriptsize, below] at (0, -0.5) {Gentle bend};
    \node[font=\scriptsize, below] at (0, -1.0) {Low $a_\perp$ needed};
  \end{scope}

  % High curvature region
  \begin{scope}[shift={(4,0)}]
    \node[font=\bfseries\small, accentred] at (0, 2.2) {High curvature};
    \draw[very thick, accentred!80, decoration={markings,
        mark=at position 0.5 with {\arrow{Stealth[length=3mm]}}
      }, postaction={decorate}]
      (-2, -0.8) .. controls (-0.3, 1.8) and (0.3, 1.8) .. (2, -0.8);
    % Osculating circle (small)
    \draw[thin, dashed, gray] (0, 0.9) circle (0.9);
    \draw[<->, thin, gray] (0, 0.9) -- (0.85, 0.75)
      node[midway, above, font=\scriptsize] {$R$ small};
    \node[font=\scriptsize, below] at (0, -1.3) {Tight bend};
    \node[font=\scriptsize, below] at (0, -1.8) {High $a_\perp$ needed};
  \end{scope}
\end{tikzpicture}
\end{center}

% ══════════════════════════════════════════════
\section{The Frenet--Serret Frame}
% ══════════════════════════════════════════════

The tangent vector $\tangent$ is only the first element of a natural coordinate
frame that moves with the curve. This \textbf{Frenet--Serret frame} (or simply
\textbf{moving frame}) provides a complete local coordinate system at every point
of the trajectory.

\subsection{Construction in $\Reals^n$}

\begin{definition}{Frenet--Serret Frame}{frenet}
  Let $\traj(s)$ be a curve parameterized by arc length in $\Reals^n$ with
  $\traj^{(k)}(s)$ denoting the $k$-th derivative with respect to $s$. Assume the
  first $p \leq n$ derivatives $\{\traj'(s), \traj''(s), \ldots, \traj^{(p)}(s)\}$
  are linearly independent at each point. The \textbf{Frenet--Serret frame}
  $\{\frenet_1(s), \frenet_2(s), \ldots, \frenet_p(s)\}$ is defined by
  applying the Gram--Schmidt process to these derivatives:
  \begin{align}
    \frenet_1 &= \traj' = \tangent, \label{eq:frenet_e1} \\[4pt]
    \tilde{\frenet}_k &= \traj^{(k)} - \sum_{j=1}^{k-1}
      \inner{\traj^{(k)}}{\frenet_j}\, \frenet_j, \qquad
    \frenet_k = \frac{\tilde{\frenet}_k}{\norm{\tilde{\frenet}_k}},
    \quad k = 2, \ldots, p. \label{eq:frenet_ek}
  \end{align}
  The frame vectors $\frenet_1, \ldots, \frenet_p$ form an orthonormal set at
  each point. If $p < n$, the frame can be extended to a full orthonormal basis
  by appending $n - p$ vectors, but these additional vectors have no intrinsic
  geometric meaning for the curve.
\end{definition}

For the familiar case $n = 3$:

\begin{center}
\renewcommand{\arraystretch}{1.5}
\begin{tabular}{@{} c l l @{}}
  \toprule
  \textbf{Vector} & \textbf{Name} & \textbf{Interpretation} \\
  \midrule
  $\frenet_1 = \tangent$ & Unit tangent & Direction of motion \\
  $\frenet_2 = \normal$ & Principal normal & Direction the curve is bending toward \\
  $\frenet_3 = \binormal$ & Binormal & $\tangent \times \normal$;
    normal to the osculating plane \\
  \bottomrule
\end{tabular}
\end{center}

\subsection{The Frenet--Serret Equations}

The power of the moving frame lies in how it evolves along the curve. The
derivatives of the frame vectors with respect to arc length are expressed
entirely in terms of the frame itself and a set of scalar functions called
\textbf{generalized curvatures}.

\begin{theorem}[Frenet--Serret Equations]\label{thm:frenet_serret}
  The Frenet--Serret frame satisfies
  \begin{equation}\label{eq:frenet_serret}
    \frac{\dd}{\ds}
    \begin{pmatrix} \frenet_1 \\ \frenet_2 \\ \frenet_3 \\ \vdots \\ \frenet_p \end{pmatrix}
    =
    \begin{pmatrix}
      0 & \kappa_1 & 0 & \cdots & 0 \\
      -\kappa_1 & 0 & \kappa_2 & \cdots & 0 \\
      0 & -\kappa_2 & 0 & \cdots & 0 \\
      \vdots & & \ddots & \ddots & \kappa_{p-1} \\
      0 & \cdots & 0 & -\kappa_{p-1} & 0
    \end{pmatrix}
    \begin{pmatrix} \frenet_1 \\ \frenet_2 \\ \frenet_3 \\ \vdots \\ \frenet_p \end{pmatrix},
  \end{equation}
  where $\kappa_1(s) = \kappa(s) > 0$ is the curvature,
  $\kappa_2(s)$ is the \textbf{torsion} (often denoted $\tau$),
  and $\kappa_3, \ldots, \kappa_{p-1}$ are the
  \textbf{higher-order curvatures}. The matrix is skew-symmetric, which
  guarantees that the frame remains orthonormal as it evolves.
\end{theorem}

The skew-symmetry of this matrix is not a coincidence---it is the hallmark of
a \emph{rotation}. As the frame moves along the curve, it rotates without
stretching. The curvatures $\kappa_1, \ldots, \kappa_{p-1}$ are the
``angular velocities'' of this rotation projected onto the various planes of
the frame.

\begin{example}{Frenet--Serret for a 3D Helix}{helix}
  The helix $\traj(t) = (a\cos t,\; a\sin t,\; bt)$ with speed
  $v = \sqrt{a^2 + b^2}$ has arc-length parameter $s = vt$. The Frenet--Serret
  quantities are:
  \begin{align}
    \tangent &= \frac{1}{v}(-a\sin t,\; a\cos t,\; b), \\
    \normal &= (-\cos t,\; -\sin t,\; 0), \\
    \binormal &= \frac{1}{v}(b\sin t,\; -b\cos t,\; a), \\
    \kappa &= \frac{a}{a^2 + b^2}, \qquad
    \tau = \frac{b}{a^2 + b^2}.
  \end{align}
  Both curvature and torsion are \emph{constant}, which makes the helix a
  natural reference trajectory for testing tracking controllers---analogous to
  the role of step inputs in classical control. The helix is to trajectory
  control what the step function is to setpoint control.
\end{example}

\subsection{The Moving Frame in State Space}

For control applications, the state space typically has dimension $n = 2N$ where
$N$ is the number of generalized coordinates (positions and velocities). A golf
swing model with 4~DOF has $n = 8$. The Frenet--Serret frame in $\Reals^8$
has up to 7 generalized curvatures, $\kappa_1$ through $\kappa_7$.

Not all of these are physically meaningful for every system. In practice, the
effective dimensionality of the curve---the number of linearly independent
derivatives---is determined by the system dynamics:

\begin{principle}{Dynamic Rank of a Trajectory}{dynamic_rank}
  For a controlled mechanical system with $N$ DOF and $m$ independent actuators,
  a generic trajectory has at most $p = 2m$ linearly independent arc-length
  derivatives. The Frenet--Serret frame therefore has at most $2m$ intrinsically
  meaningful vectors, and the curve has at most $2m - 1$ independent curvatures.

  For a fully actuated system ($m = N$), the trajectory can bend freely in all
  $2N$ state-space directions. For an underactuated system ($m < N$), the
  trajectory is geometrically constrained: it lives on a submanifold, and the
  curvatures in the ``forbidden'' directions are determined by the passive
  dynamics.
\end{principle}

% ══════════════════════════════════════════════
\section{Curvature and Torsion in the Language of Control}
% ══════════════════════════════════════════════

We now connect the geometric quantities to the dynamics and control of
nonlinear systems.

\subsection{Curvature as a Demand on the Controller}

Consider the system $\dot{\state} = \bm{f}(\state, \control)$ tracking a
reference trajectory $\traj^*(s)$ at instantaneous speed $v(t) = \dot{s}(t)$.
Using the chain rule:
\begin{equation}
  \dot{\state} = v\, \tangent, \qquad
  \ddot{\state} = \dot{v}\, \tangent + v^2\, \kappa\, \normal.
\end{equation}
The acceleration decomposes into:
\begin{itemize}[leftmargin=1.5cm]
  \item A \textbf{tangential component} $\dot{v}\, \tangent$ that changes the
    speed along the curve.
  \item A \textbf{normal component} $v^2\kappa\, \normal$ that ``steers'' the
    system around the bend.
\end{itemize}

\begin{definition}{Control Demand Decomposition}{control_demand}
  Along a reference trajectory $\traj^*(s)$ traversed at speed $v$, the
  instantaneous control demand decomposes as
  \begin{equation}\label{eq:control_demand}
    \control_{\text{demand}} = \underbrace{\control_\parallel}_{\text{tangential}}
    + \underbrace{\control_\perp}_{\text{normal}},
  \end{equation}
  where:
  \begin{align}
    \control_\parallel &\propto \dot{v}
      && \text{(accelerate/decelerate along path)}, \\
    \control_\perp &\propto v^2\, \kappa
      && \text{(centripetal force to follow curvature)}.
  \end{align}
  The normal demand $\control_\perp$ grows quadratically with speed and linearly
  with curvature. This is the \textbf{curvature cost} of the trajectory.
\end{definition}

\begin{example}{Curvature Cost in the Golf Downswing}{golf_curv_cost}
  In the transition from backswing to downswing, a professional golfer reverses
  the direction of the club in roughly 0.15 seconds. Let us estimate the
  curvature cost.

  At the top of the backswing, the clubhead speed is approximately zero. At
  impact (0.3\,s later), it reaches $\sim\!50$\,m/s. The path through
  configuration space has a tight reversal---a region of high curvature
  $\kappa$.

  Using $a_\perp = \kappa v^2$, and noting that mid-downswing
  $v \approx 25$\,m/s with the reversal curvature $\kappa \approx 2$\,rad/m
  (estimated from motion capture data), the centripetal acceleration demand is
  \begin{equation}
    a_\perp \approx 2 \times 25^2 = 1250\;\text{m/s}^2 \approx 127\;g.
  \end{equation}
  This enormous demand is met not by muscular torque alone, but by the
  \emph{passive coupling} in the kinematic chain---the very underactuation
  that the setpoint framework treats as a constraint. The geometry of the
  trajectory creates the forces needed to execute it.
\end{example}

\subsection{Torsion as Out-of-Plane Complexity}

If curvature measures how a trajectory bends \emph{within} its osculating
plane, \textbf{torsion} measures how the osculating plane itself rotates along
the curve. A planar curve has zero torsion everywhere. A helix has constant
nonzero torsion. A golf swing---which involves simultaneous rotation in the
transverse, sagittal, and frontal planes---has torsion that varies dramatically
through the motion.

\begin{definition}{Torsion}{torsion}
  For a curve in $\Reals^n$ with $n \geq 3$, the \textbf{torsion}
  $\tau(s) = \kappa_2(s)$ is the second generalized curvature in the
  Frenet--Serret equations:
  \begin{equation}
    \frac{\dd\normal}{\ds} = -\kappa\, \tangent + \tau\, \binormal.
  \end{equation}
  Torsion measures the rate at which the osculating plane rotates about the
  tangent direction, per unit arc length.
\end{definition}

In high-dimensional state spaces ($n \gg 3$), the higher curvatures
$\kappa_3, \ldots, \kappa_{p-1}$ measure increasingly subtle geometric
twisting. For control design, the first two curvatures ($\kappa$ and $\tau$)
capture the dominant geometric effects, with higher curvatures becoming
relevant primarily in highly articulated systems.

\begin{remark}{The Curvature Signature}{curv_signature}
  The ordered set of generalized curvatures $(\kappa_1(s), \kappa_2(s),
  \ldots, \kappa_{p-1}(s))$ as functions of arc length constitutes the
  \textbf{curvature signature} of the trajectory. By the fundamental theorem
  of curves, this signature (up to rigid motions) \emph{uniquely determines}
  the trajectory. Two motions with identical curvature signatures are
  geometrically congruent---they are the same shape, possibly translated
  and rotated.

  This has a striking implication for motor control: if a golfer can reproduce
  the curvature signature of their swing, the spatial path is automatically
  reproduced, regardless of timing.
\end{remark}

% ══════════════════════════════════════════════
\section{Reparameterization and Timing Laws}
% ══════════════════════════════════════════════

We now formalize the separation between the \emph{shape} of a trajectory and
the \emph{timing} of its execution.

\begin{definition}{Timing Law}{timing_law}
  Given a reference trajectory $\traj^*(s)$ parameterized by arc length, a
  \textbf{timing law} is a smooth, monotonically increasing function
  $s: [0, T] \to [0, L]$ with $s(0) = 0$ and $s(T) = L$. The
  time-parameterized trajectory is
  \begin{equation}
    \state^*(t) = \traj^*(s(t)).
  \end{equation}
  The timing law is equivalently specified by the speed profile $v(t) = \dot{s}(t)$
  or the phase velocity $v(s) = \dd s / \dd t$ expressed as a function of arc
  length.
\end{definition}

Different timing laws produce different motions through the same geometric
path:

\begin{center}
\begin{tikzpicture}[scale=1.0]
  % s(t) plots
  \begin{scope}[shift={(-4.5, 0)}]
    \draw[-{Stealth[length=2.5mm]}, thick] (0, 0) -- (3.5, 0)
      node[right, font=\small] {$t$};
    \draw[-{Stealth[length=2.5mm]}, thick] (0, 0) -- (0, 3)
      node[above, font=\small] {$s(t)$};
    % Constant speed
    \draw[very thick, chapblue] (0, 0) -- (3, 2.5);
    \node[chapblue, font=\scriptsize, right] at (2.5, 2.7) {constant speed};
    % Accelerating
    \draw[very thick, accentred] (0, 0) .. controls (2.5, 0.3) .. (3, 2.5);
    \node[accentred, font=\scriptsize, right] at (2.5, 1.2) {accelerating};
    % Decelerating
    \draw[very thick, trajgreen] (0, 0) .. controls (0.5, 2.2) .. (3, 2.5);
    \node[trajgreen, font=\scriptsize, left] at (0.7, 2.3) {decelerating};
    \node[font=\small\bfseries] at (1.5, -0.7) {Three timing laws};
  \end{scope}

  % Same path
  \begin{scope}[shift={(4, 0)}]
    \draw[-{Stealth[length=2.5mm]}, thick, gray] (-1.5, 0) -- (2, 0)
      node[right, font=\small] {$x_1$};
    \draw[-{Stealth[length=2.5mm]}, thick, gray] (0, -0.5) -- (0, 3)
      node[above, font=\small] {$x_2$};
    \draw[very thick, black!50, decoration={markings,
        mark=at position 0.5 with {\arrow{Stealth[length=3mm]}}
      }, postaction={decorate}]
      (-1, 0) .. controls (0, 2.5) and (1, 2.5) .. (1.5, 0.5);
    \node[font=\small\bfseries] at (0.3, -0.8) {Same spatial path};
    % Speed markers
    \foreach \p/\c in {0.2/chapblue, 0.4/chapblue, 0.6/chapblue, 0.8/chapblue} {
      \filldraw[\c, opacity=0.5]
        ({-1 + 2.5*\p + 0.1*sin(360*\p)}, {2.5*sin(180*\p) - 0.2*\p})
        circle (2pt);
    }
  \end{scope}
\end{tikzpicture}
\end{center}

\subsection{The Path--Timing Decomposition of Control}

The reparameterization framework leads to a clean decomposition of the
control problem:

\begin{principle}{Path--Timing Separation}{path_timing}
  The trajectory control problem decomposes into two subproblems:
  \begin{enumerate}[label=\textbf{(\Alph*)}]
    \item \textbf{Path design}: Choose the geometric curve $\traj^*(s)$ in state
      space that achieves the task (e.g., delivers the clubhead to the ball with
      the correct velocity vector).
    \item \textbf{Timing design}: Choose the timing law $s(t)$ that balances
      speed against tracking difficulty, control effort, and robustness.
  \end{enumerate}
  These two problems have different objective functions and constraints:
  \begin{center}
  \renewcommand{\arraystretch}{1.3}
  \begin{tabular}{@{} l l l @{}}
    \toprule
    & \textbf{Path design} & \textbf{Timing design} \\
    \midrule
    Objective & Task performance & Trackability \& effort \\
    Key quantity & Curvature $\kappa(s)$ & Speed profile $v(s)$ \\
    Constraint & Dynamic feasibility & Actuator limits \\
    \bottomrule
  \end{tabular}
  \end{center}
  The interaction between the two is captured by the curvature cost
  $\kappa(s)\, v(s)^2$, which the timing law must keep within the controller's
  authority at every point.
\end{principle}

\begin{example}{Timing the Downswing}{timing_downswing}
  Suppose the downswing path has a curvature profile $\kappa(s)$ with a peak
  of $\kappa_{\max}$ at the transition point. The maximum centripetal
  acceleration the controller can provide is $a_{\max}$. Then the speed at
  the transition is bounded by
  \begin{equation}
    v_{\text{transition}} \leq \sqrt{\frac{a_{\max}}{\kappa_{\max}}}.
  \end{equation}
  This is a \emph{geometric speed limit}: no matter how powerful the actuators,
  the system cannot traverse a tight bend faster than its control authority
  allows. The timing law must respect this limit, typically by slowing down
  through high-curvature regions and accelerating through straight sections.

  This is precisely what golfers do instinctively. The ``pause at the top''---the
  brief deceleration between backswing and downswing---is not aesthetic
  affectation. It is the timing law respecting the curvature speed limit at
  the point of maximum path curvature.
\end{example}

% ══════════════════════════════════════════════
\section{The Osculating Objects}
% ══════════════════════════════════════════════

The Frenet--Serret frame defines a sequence of ``best-fit'' geometric objects at
each point of the curve. These osculating objects provide increasingly refined
local approximations that are useful for controller design.

\begin{definition}{Osculating Objects}{osculating}
  At a point $\traj^*(s_0)$ on a curve in $\Reals^n$:
  \begin{enumerate}[label=(\roman*)]
    \item The \textbf{osculating line} is the tangent line: the best linear
      approximation.
    \item The \textbf{osculating circle} lies in the $(\tangent, \normal)$ plane,
      has radius $R = 1/\kappa$, and center at $\traj^*(s_0) + R\,\normal$.
      It is the best second-order (circular) approximation.
    \item The \textbf{osculating plane} is the plane spanned by $\tangent$ and
      $\normal$. Torsion measures how fast the curve leaves this plane.
    \item The \textbf{osculating sphere} is the best third-order approximation
      and accounts for both curvature and torsion.
  \end{enumerate}
\end{definition}

For trajectory control, the osculating circle is the most immediately useful.
It tells us: \emph{if the system were constrained to follow a circle at this
point, what circle would it be?} The radius of that circle is the radius of
curvature $R = 1/\kappa$, and the center indicates the direction toward which
the trajectory bends.

% ══════════════════════════════════════════════
\section{The Geometry of Trajectory Deviations}
% ══════════════════════════════════════════════

We now use the Frenet--Serret frame to analyze what happens when the system
deviates from the reference trajectory. This provides the geometric foundation
for the transverse linearization of Chapter~4.

\subsection{Frenet--Serret Coordinates}

Given a reference trajectory $\traj^*(s)$ and the associated Frenet--Serret
frame $\{\frenet_1(s), \ldots, \frenet_p(s)\}$, any state $\state$ near the
trajectory can be decomposed as:
\begin{equation}\label{eq:frenet_coords}
  \state = \traj^*(s^*) + \sum_{k=2}^{n} \xi_k\, \frenet_k(s^*),
\end{equation}
where $s^* = s^*(\state)$ is the projection phase (the arc-length parameter of
the nearest point on the reference), and $\xi_2, \ldots, \xi_n$ are the
\textbf{transverse coordinates} measuring deviation in each normal direction.

\begin{definition}{Frenet--Serret Coordinates}{frenet_coords}
  The \textbf{Frenet--Serret coordinate representation} of a state $\state$ relative
  to a reference trajectory $\traj^*$ is the tuple
  \begin{equation}
    (s^*,\; \xi_2,\; \xi_3,\; \ldots,\; \xi_n),
  \end{equation}
  where $s^*$ is the tangential (along-trajectory) coordinate and
  $\bm{\xi} = (\xi_2, \ldots, \xi_n)$ are the transverse (cross-trajectory)
  coordinates. The transverse distance from Chapter~1 is
  \begin{equation}
    d_\perp = \norm{\bm{\xi}} = \sqrt{\xi_2^2 + \cdots + \xi_n^2}.
  \end{equation}
\end{definition}

These coordinates have a powerful property: the reference trajectory itself
corresponds to $\bm{\xi} = \bm{0}$, and the trajectory-tracking problem
becomes the problem of \emph{regulating $\bm{\xi}$ to zero while $s^*$
advances}. This converts the tracking problem into a regulation problem, but in
the moving frame rather than a fixed frame.

\subsection{Curvature-Induced Coupling}

There is a subtlety: the Frenet--Serret coordinates are not inertial. Because
the frame rotates as it moves along the curve (governed by the Frenet--Serret
equations \eqref{eq:frenet_serret}), fictitious forces appear in the transverse
dynamics. The dominant effect is a \textbf{curvature-induced coupling}:

\begin{proposition}\label{prop:curv_coupling}
  In Frenet--Serret coordinates, the linearized transverse dynamics about the
  reference trajectory have the form
  \begin{equation}\label{eq:transverse_lin}
    \dot{\bm{\xi}} = \bigl[\bm{A}(s) - v^2 \bm{K}(s)\bigr]\, \bm{\xi}
    + \bm{B}(s)\, \control_{\text{fb}},
  \end{equation}
  where $\bm{A}(s)$ comes from linearizing the original dynamics, and
  $\bm{K}(s)$ is a matrix that depends on the generalized curvatures
  $\kappa_1(s), \ldots, \kappa_{p-1}(s)$. The term $v^2\bm{K}(s)$ represents
  the geometric coupling between the curve's shape and the transverse stability.
\end{proposition}

This result has a profound implication: \emph{the stability of trajectory
tracking depends on the geometry of the trajectory.} A trajectory with low
curvature is inherently easier to stabilize than one with high curvature,
even if the underlying system dynamics are identical. The curvature modifies
the effective dynamics transverse to the path.

% ══════════════════════════════════════════════
\section{Fundamental Theorem and Uniqueness}
% ══════════════════════════════════════════════

We close this chapter with a theorem that is central to the trajectory-first
philosophy.

\begin{theorem}[Fundamental Theorem of Curves]\label{thm:fundamental}
  Let $\kappa_1(s), \ldots, \kappa_{p-1}(s)$ be smooth functions on
  $[0, L]$ with $\kappa_k(s) > 0$ for $k = 1, \ldots, p-2$. Then there
  exists a curve $\traj: [0, L] \to \Reals^n$, parameterized by arc length,
  whose generalized curvatures are exactly $\kappa_1, \ldots, \kappa_{p-1}$.
  Moreover, this curve is unique up to rigid motions (translations and
  rotations) of $\Reals^n$.
\end{theorem}

This theorem says: \textbf{the curvature signature is the trajectory.} Two
trajectories with the same curvature functions are identical in shape. This has
a practical corollary for control:

\begin{corollary}
  To specify a reference trajectory, it suffices to specify the generalized
  curvatures as functions of arc length, together with an initial position
  and orientation. This representation is:
  \begin{enumerate}[label=(\roman*)]
    \item \emph{Compact}: $p - 1$ scalar functions instead of $n$ coordinate
      functions.
    \item \emph{Intrinsic}: independent of coordinate choice.
    \item \emph{Timing-independent}: the trajectory shape is fully determined
      without specifying when each point is reached.
  \end{enumerate}
\end{corollary}

\begin{example}{Curvature Signature of a Golf Swing}{golf_signature}
  A simplified 2-DOF model of the golf swing (arm + club) has a 4-dimensional
  state space ($\theta_{\text{arm}}, \theta_{\text{club}}, \dot{\theta}_{\text{arm}},
  \dot{\theta}_{\text{club}}$). The trajectory through this space has up to 3
  generalized curvatures: $\kappa_1(s)$, $\kappa_2(s)$, and $\kappa_3(s)$.

  The curvature signature reveals the structure of the swing:
  \begin{itemize}[leftmargin=1cm]
    \item $\kappa_1(s)$ peaks at the transition (the ``pause at the top'').
    \item $\kappa_2(s)$ peaks during the release, where the swing plane rotates.
    \item $\kappa_3(s)$ is small for a well-executed swing (the motion is
      approximately confined to a 3D subspace of $\Reals^4$).
  \end{itemize}

  A coach's instruction to ``flatten your swing plane'' is, in geometric terms,
  a request to reduce $\kappa_2(s)$ during a specific phase of the motion---a
  precise modification to the curvature signature.
\end{example}

% ══════════════════════════════════════════════
\section{Summary}
% ══════════════════════════════════════════════

This chapter has established the geometric language we will use throughout the
book:

\begin{center}
\renewcommand{\arraystretch}{1.4}
\begin{tabular}{@{} l p{9cm} @{}}
  \toprule
  \textbf{Concept} & \textbf{Control significance} \\
  \midrule
  Arc length $s$ & The natural ``distance along the path''---separates shape
    from timing \\
  Curvature $\kappa(s)$ & Determines the centripetal control demand
    $a_\perp = \kappa v^2$ \\
  Torsion $\tau(s)$ & Measures out-of-plane complexity; relates to
    multi-joint coordination \\
  Frenet--Serret frame & Provides the natural moving coordinate system for
    decomposing deviations into tangential and transverse components \\
  Curvature signature & Uniquely specifies the trajectory shape, independent
    of timing and coordinates \\
  Path--timing separation & Decomposes control into geometric path design and
    temporal speed planning \\
  \bottomrule
\end{tabular}
\end{center}

\vspace{0.5cm}

The key takeaway is that a trajectory is not just ``a function of time.'' It is a
geometric object with intrinsic structure, and this structure directly determines
the difficulty and requirements of controlling it. In Chapter~3, we will see that
the state spaces of mechanical systems are not flat $\Reals^n$---they are curved
manifolds with their own geometry. The interplay between the geometry of the
\emph{curve} and the geometry of the \emph{space} will produce the richest and
most practically important results of this book.

% ══════════════════════════════════════════════
\section{Exercises}
% ══════════════════════════════════════════════

\begin{enumerate}[label=\textbf{2.\arabic*.}, leftmargin=1.5cm]

\item \textbf{(Computation.)} Compute the curvature and torsion of the curve
  $\traj(t) = (t, t^2, t^3)$ at $t = 0$ and $t = 1$. At which point is the
  curve ``harder to track'' at constant speed?

\item \textbf{(Arc length.)} The logarithmic spiral $\traj(t) = e^{at}(\cos t,\;
  \sin t)$ with $a > 0$ is a model for certain biological growth patterns.
  \begin{enumerate}[label=(\alph*)]
    \item Find the arc-length function $s(t)$.
    \item Show that the curvature is $\kappa = 1/(e^{at}\sqrt{1 + a^2})$ and
      interpret what this means about tracking difficulty as the spiral grows.
    \item What timing law $s(t)$ would maintain constant centripetal acceleration?
  \end{enumerate}

\item \textbf{(Frenet--Serret.)} For the curve $\traj(s) = (\cos(s/\sqrt{2}),\;
  \sin(s/\sqrt{2}),\; s/\sqrt{2})$ (a helix parameterized by arc length):
  \begin{enumerate}[label=(\alph*)]
    \item Verify that $\norm{\traj'(s)} = 1$.
    \item Compute the Frenet--Serret frame $\{\tangent, \normal, \binormal\}$.
    \item Compute $\kappa$ and $\tau$ and verify the Frenet--Serret equations.
  \end{enumerate}

\item \textbf{(Path--timing separation.)} A robot end-effector must trace the
  path $\bm{p}(s) = (s,\; \sin(\pi s),\; 0)$ for $s \in [0, 2]$
  (one full sine wave).
  \begin{enumerate}[label=(\alph*)]
    \item Compute $\kappa(s)$ and identify the points of maximum curvature.
    \item If the maximum centripetal acceleration is $a_{\max} = 10\;\text{m/s}^2$,
      find the speed limit $v_{\max}(s)$ along the path.
    \item Design a timing law $s(t)$ that traverses the path in minimum time
      while respecting the speed limit. Sketch $v(s)$ and $s(t)$.
  \end{enumerate}

\item \textbf{(Conceptual.)} Consider two trajectories in $\Reals^4$ with
  identical curvature signatures $(\kappa_1(s), \kappa_2(s), \kappa_3(s))$
  but different initial conditions.
  \begin{enumerate}[label=(\alph*)]
    \item By the fundamental theorem, how are these trajectories related?
    \item If you designed a tracking controller for one, would it work for the
      other? Under what conditions?
    \item Relate this to the idea that a well-learned golf swing can be started
      from slightly different address positions.
  \end{enumerate}

\item \textbf{(Numerical project.)} Using motion capture data (or simulated data
  from a double pendulum), compute the curvature signature of a swing-like motion:
  \begin{enumerate}[label=(\alph*)]
    \item Numerically estimate $s(t)$ by integrating $\norm{\dot{\state}(t)}$.
    \item Compute $\kappa_1(s)$ using Eq.~\eqref{eq:curvature_general_n}.
    \item Identify the phases of the motion (backswing, transition, downswing,
      impact) from features of the curvature profile.
    \item Plot the curvature cost $\kappa(s)\, v(s)^2$ and identify the most
      demanding phase of the motion.
  \end{enumerate}

\end{enumerate}

\vspace{1cm}
\begin{center}
  \rule{0.5\textwidth}{0.4pt}\\[0.5cm]
  {\small\itshape
  The shape of a river is written in its curvature.\\
  The story of a motion is written the same way.\\[0.3cm]
  Learn to read the curvature, and you read the motion.}
\end{center}

\chapter{Configuration Manifolds}

\epigraph{%
  The map is not the territory. The coordinate chart is not the manifold.
  Every time you write $\theta$ for an angle, you are lying---just a little.
}{}

\vspace{0.5cm}

% ══════════════════════════════════════════════
\section{The Lie That Coordinates Tell}
% ══════════════════════════════════════════════

In Chapter~2 we developed the geometry of curves in $\Reals^n$. This was a
necessary starting point, but it was also a simplification. The state spaces
of real mechanical systems are \emph{not} $\Reals^n$.

Consider the simplest possible rotating joint: a hinge. Its configuration is
an angle $\theta$, which we habitually write as a real number. But an angle
is not a real number. The configuration $\theta = 0$ and the configuration
$\theta = 2\pi$ are \emph{the same physical state}. No point in $\Reals$
has this property. The true configuration space of a hinge is the
\emph{circle} $\Sone$---a one-dimensional manifold that is locally like
$\Reals$ but globally has a completely different topology.

This is not pedantry. It is the source of real engineering failures.

\begin{warning}{The Gimbal Lock Catastrophe}{gimbal_lock}
  If we represent orientation using three Euler angles
  $(\phi, \theta, \psi) \in \Reals^3$, we encounter \textbf{gimbal lock}:
  configurations where the representation degenerates and loses a degree of
  freedom. At $\theta = \pm\pi/2$, the axes for $\phi$ and $\psi$ align,
  making the system appear to have only 2~DOF when it actually has~3.

  This is not a hardware problem---it is a \emph{coordinate singularity},
  an artifact of forcing a non-Euclidean space ($\SOthree$) into Euclidean
  coordinates ($\Reals^3$). Apollo~13's inertial measurement unit
  experienced gimbal lock operationally. Every modern spacecraft avoids
  Euler angles for precisely this reason.

  The lesson for control: \textbf{if you use the wrong coordinates, your
  controller will have singularities that the physical system does not.}
\end{warning}

\begin{keyidea}{The Manifold Imperative}{manifold_imperative}
  The configuration space of a mechanical system is a \emph{smooth manifold}
  $\configspace$---a space that is locally like $\Reals^n$ but may have
  nontrivial global topology. Control theory built on $\Reals^n$ must be
  replaced by control theory built on $\configspace$ whenever:
  \begin{enumerate}[label=(\roman*)]
    \item Joints allow full rotation (wrapping around $\Sone$ or $\SOthree$).
    \item The trajectory traverses a significant fraction of the configuration space.
    \item Singularity-free representation is required for robustness.
  \end{enumerate}
  For the golf swing---which involves large rotations in multiple joints
  simultaneously---all three conditions are met.
\end{keyidea}

% ══════════════════════════════════════════════
\section{Smooth Manifolds: The Minimum You Need}
% ══════════════════════════════════════════════

We develop only as much manifold theory as the control engineer needs. The
reader seeking a comprehensive treatment should consult Lee~(2012) or
Abraham \& Marsden~(1978).

\begin{definition}{Smooth Manifold}{manifold}
  A \textbf{smooth manifold} $\manifold$ of dimension $n$ is a set together
  with a collection of \textbf{coordinate charts}
  $\{(U_\alpha, \varphi_\alpha)\}$ such that:
  \begin{enumerate}[label=(\roman*)]
    \item Each $U_\alpha \subseteq \manifold$ is an open set, and the
      $U_\alpha$ cover $\manifold$.
    \item Each $\varphi_\alpha: U_\alpha \to \Reals^n$ is a homeomorphism
      onto an open subset of $\Reals^n$.
    \item Where two charts overlap, the \textbf{transition maps}
      $\varphi_\beta \circ \varphi_\alpha^{-1}$ are smooth
      ($C^\infty$) functions from $\Reals^n$ to $\Reals^n$.
  \end{enumerate}
\end{definition}

The key intuition is: a manifold is a space where you can \emph{locally} use
coordinates, but no single coordinate system works \emph{globally}. You need
an atlas of overlapping charts, like a road atlas needs multiple pages to cover
a country.

\begin{example}{The Circle $\Sone$}{circle_manifold}
  The circle $\Sone = \{(\cos\theta, \sin\theta) : \theta \in \Reals\}$ is a
  1-dimensional manifold. It cannot be covered by a single coordinate chart:
  any angle parameterization $\theta \in [0, 2\pi)$ has a discontinuity
  at $\theta = 0 = 2\pi$. We need at least two charts:
  \begin{align}
    U_1 &= \Sone \setminus \{(-1, 0)\}, &
    \varphi_1(p) &= \theta \in (-\pi, \pi), \\
    U_2 &= \Sone \setminus \{(1, 0)\}, &
    \varphi_2(p) &= \theta \in (0, 2\pi).
  \end{align}
  On the overlap $U_1 \cap U_2$, the transition map is either the identity
  or a shift by $2\pi$---both smooth.

  For a hinge joint, this means: no single angle variable $\theta$ can
  represent all configurations without a discontinuity. If the joint rotates
  past the discontinuity, any controller using that angle will experience
  a ``jump''---a coordinate artifact, not a physical event.
\end{example}

\subsection{Tangent Spaces and Tangent Vectors}

On a manifold, there is no global notion of ``direction.'' Velocity is a
\emph{local} concept, defined at each point.

\begin{definition}{Tangent Space}{tangent_space}
  The \textbf{tangent space} $\tangentspace{p}$ at a point $p \in \manifold$
  is the vector space of all possible velocity vectors at $p$. If $\manifold$
  has dimension $n$, then $\tangentspace{p} \cong \Reals^n$ as a vector
  space, but each point has its \emph{own} tangent space.

  Given a curve $\gamma: (-\varepsilon, \varepsilon) \to \manifold$ with
  $\gamma(0) = p$, the tangent vector $\dot{\gamma}(0)$ is an element of
  $\tangentspace{p}$.
\end{definition}

\begin{definition}{Tangent Bundle}{tangent_bundle}
  The \textbf{tangent bundle} $T\!\manifold = \bigsqcup_{p \in \manifold}
  \tangentspace{p}$ is the collection of all tangent spaces, one for each
  point. It is itself a smooth manifold of dimension $2n$. A point in
  $T\!\manifold$ is a pair $(p, v)$ where $p \in \manifold$ and
  $v \in \tangentspace{p}$.

  For a mechanical system with configuration manifold $\configspace$, the
  \textbf{state space is the tangent bundle} $\statespace = \tangentbundle$.
  A state is a pair $(\bm{q}, \dot{\bm{q}})$---a configuration and a
  velocity.
\end{definition}

\begin{center}
\begin{tikzpicture}[scale=1.1]
  % Manifold surface
  \draw[thick, fill=chapblue!8]
    (-3, 0) .. controls (-2, 1.2) and (-0.5, 1.5) .. (0, 1.2)
    .. controls (0.5, 0.9) and (2, 1.5) .. (3, 0.8)
    .. controls (3.3, 0.2) and (2.5, -0.5) .. (1.5, -0.3)
    .. controls (0.5, -0.1) and (-0.5, -0.8) .. (-1.5, -0.5)
    .. controls (-2.5, -0.2) and (-3.3, 0.2) .. (-3, 0);
  \node[font=\large\bfseries, chapblue] at (-2.2, 0.5) {$\manifold$};

  % Point p
  \coordinate (P) at (0.5, 0.7);
  \filldraw[accentred] (P) circle (2.5pt);
  \node[below left, font=\small\bfseries, accentred] at (P) {$p$};

  % Tangent plane
  \draw[fill=darkgold!10, opacity=0.7, thick, darkgold!60]
    (-0.8, 1.8) -- (1.8, 2.5) -- (2.2, 0.5) -- (-0.4, -0.2) -- cycle;
  \node[font=\small, darkgold!80!black] at (1.8, 1.0) {$\tangentspace{p}$};

  % Tangent vector
  \draw[-{Stealth[length=3.5mm]}, very thick, accentred]
    (P) -- ++(1.0, 0.8) node[above right, font=\small] {$v \in \tangentspace{p}$};

  % Curve on manifold
  \draw[thick, trajgreen, decoration={markings,
      mark=at position 0.6 with {\arrow{Stealth[length=2.5mm]}}
    }, postaction={decorate}]
    (-1.5, 0.1) .. controls (-0.5, 0.3) and (0, 0.5) .. (P)
    .. controls (1.0, 0.9) and (1.8, 1.0) .. (2.3, 0.7);
  \node[trajgreen!80!black, font=\small] at (-1.7, 0.4) {$\gamma(t)$};
\end{tikzpicture}
\end{center}

% ══════════════════════════════════════════════
\section{The Configuration Manifolds of Joints}
% ══════════════════════════════════════════════

Every mechanical joint constrains relative motion between two rigid bodies.
The set of allowed relative configurations forms a manifold whose topology
depends on the joint type.

\begin{center}
\renewcommand{\arraystretch}{1.5}
\begin{tabular}{@{} l c c l @{}}
  \toprule
  \textbf{Joint type} & \textbf{DOF} & \textbf{Manifold} & \textbf{Topology} \\
  \midrule
  Revolute (hinge) & 1 & $\Sone$ & Circle \\
  Prismatic (slider) & 1 & $\Reals$ & Line (or interval) \\
  Universal & 2 & $\Stwo$ & 2-sphere \\
  Ball-and-socket & 3 & $\SOthree$ & Rotation group \\
  Planar & 3 & $\SEthree|_{\text{2D}} = \Reals^2 \times \Sone$ & Plane + rotation \\
  Free body & 6 & $\SEthree = \Reals^3 \rtimes \SOthree$ & Rigid motions of 3-space \\
  \bottomrule
\end{tabular}
\end{center}

The configuration space of a multi-body system is the \emph{product} of
individual joint manifolds (modulo kinematic constraints):

\begin{definition}{Product Configuration Space}{product_config}
  For a serial chain of $k$ joints with individual configuration manifolds
  $\manifold_1, \ldots, \manifold_k$, the total configuration space is
  \begin{equation}\label{eq:product_config}
    \configspace = \manifold_1 \times \manifold_2 \times \cdots \times \manifold_k.
  \end{equation}
  For example, a robot arm with three revolute joints has
  $\configspace = \Sone \times \Sone \times \Sone = \Torus^3$, the
  3-dimensional torus.
\end{definition}

\begin{example}{The Human Shoulder}{shoulder}
  The glenohumeral (shoulder) joint is approximately a ball-and-socket joint
  with configuration manifold $\SOthree$. Combined with the 1-DOF hinge at
  the elbow ($\Sone$) and the 1-DOF pronation/supination of the forearm
  ($\Sone$), the arm has configuration space
  \begin{equation}
    \configspace_{\text{arm}} = \SOthree \times \Sone \times \Sone.
  \end{equation}
  This is a 5-dimensional manifold with the topology of $\SOthree \times \Torus^2$.
  No set of 5 angle variables can parameterize this space without singularities.
\end{example}

% ══════════════════════════════════════════════
\section{The Rotation Group $\SOthree$}
% ══════════════════════════════════════════════

The rotation group $\SOthree$ appears wherever a joint allows full 3D rotation.
It is the single most important non-Euclidean manifold in mechanical control,
and its properties drive many of the challenges we will encounter.

\begin{definition}{The Special Orthogonal Group}{SO3}
  $\SOthree$ is the group of $3 \times 3$ real orthogonal matrices with
  determinant $+1$:
  \begin{equation}
    \SOthree = \{\bm{R} \in \Reals^{3 \times 3} :
    \bm{R}^\top\!\bm{R} = \bm{I},\; \det(\bm{R}) = 1\}.
  \end{equation}
  It is a 3-dimensional compact Lie group. As a manifold, it is diffeomorphic
  to the real projective space $\Reals P^3$.
\end{definition}

\subsection{Why Euler Angles Fail}

The most common parameterization of $\SOthree$ uses three Euler angles
$(\phi, \theta, \psi)$, defining the rotation matrix as a product of three
elemental rotations. This parameterization has two fundamental problems:

\begin{enumerate}[label=\textbf{P\arabic*.}, leftmargin=1.5cm]
  \item \textbf{Singularities.} Every 3-parameter representation of $\SOthree$
    has at least one singularity. For the ZYX convention, gimbal lock occurs
    at $\theta = \pm\pi/2$. This is a topological necessity, not a flaw
    in the choice of convention: $\SOthree$ is not homeomorphic to any open
    subset of $\Reals^3$.

  \item \textbf{Non-uniqueness.} Euler angles are periodic: $(\phi, \theta, \psi)$
    and $(\phi + 2\pi, \theta, \psi)$ represent the same rotation. Near the
    singularity, drastically different angle triples can represent nearly
    identical orientations, creating numerical chaos in controllers.
\end{enumerate}

\begin{principle}{Representation Principle for Rotations}{rotation_rep}
  For trajectory control on $\SOthree$, use one of:
  \begin{enumerate}[label=(\alph*)]
    \item \textbf{Rotation matrices} $\bm{R} \in \Reals^{3 \times 3}$: 9
      parameters with 6 orthogonality constraints. Singularity-free but
      redundant.
    \item \textbf{Unit quaternions} $\bm{q} \in S^3 \subset \Reals^4$: 4
      parameters with 1 unit-norm constraint. Singularity-free, minimal
      redundancy, excellent for interpolation and computation.
    \item \textbf{Exponential coordinates} (axis-angle or Lie algebra):
      3 parameters with a singularity at $\norm{\bm{\omega}} = \pi$, but
      natural for small rotations and linearization.
  \end{enumerate}
  Never use Euler angles for control of large-rotation trajectories.
\end{principle}

\subsection{The Lie Algebra $\sothree$}

The tangent space of $\SOthree$ at the identity is the \textbf{Lie algebra}
$\sothree$---the space of $3 \times 3$ skew-symmetric matrices:
\begin{equation}
  \sothree = \{\bm{\Omega} \in \Reals^{3\times 3} :
  \bm{\Omega}^\top = -\bm{\Omega}\}
  = \left\{ \skewmat{\bm{\omega}} :=
  \begin{pmatrix} 0 & -\omega_3 & \omega_2 \\
  \omega_3 & 0 & -\omega_1 \\ -\omega_2 & \omega_1 & 0 \end{pmatrix}
  : \bm{\omega} \in \Reals^3 \right\}.
\end{equation}
The ``hat map'' $(\cdot)^\wedge: \Reals^3 \to \sothree$ and its inverse
``vee map'' $(\cdot)^\vee: \sothree \to \Reals^3$ identify the Lie algebra
with $\Reals^3$ (the space of angular velocity vectors).

\begin{definition}{Exponential Map}{exp_map}
  The \textbf{exponential map} $\exp: \sothree \to \SOthree$ connects the
  Lie algebra to the Lie group:
  \begin{equation}\label{eq:exp_map}
    \bm{R} = \exp(\skewmat{\bm{\omega}})
    = \bm{I} + \frac{\sin\theta}{\theta}\skewmat{\bm{\omega}}
    + \frac{1 - \cos\theta}{\theta^2}\skewmat{\bm{\omega}}^2,
  \end{equation}
  where $\theta = \norm{\bm{\omega}}$. This is the \textbf{Rodrigues formula}.
  Geometrically, $\bm{R}$ is a rotation by angle $\theta$ about the axis
  $\bm{\omega}/\theta$.

  The inverse map $\log: \SOthree \to \sothree$ is well-defined for
  rotations with angle $\theta < \pi$ and gives the \textbf{rotation vector}
  representation.
\end{definition}

The exponential map is the bridge between the linear world (the Lie algebra,
where we can add and scale) and the nonlinear world (the Lie group, where
the actual dynamics live). This bridge is the central tool for linearization
on manifolds.

% ══════════════════════════════════════════════
\section{The Golf Swing Configuration Space}
% ══════════════════════════════════════════════

We now apply the manifold framework to our running example. The result will
show exactly why Euclidean control theory is insufficient for the golf swing.

\subsection{A Simplified Kinematic Model}

Consider a 4-segment model of the golf swing:
\begin{enumerate}[label=(\roman*)]
  \item \textbf{Torso}: rotates relative to ground via the spine. Modeled as a
    3-DOF ball joint: $\manifold_1 = \SOthree$.
  \item \textbf{Lead arm}: rotates relative to the torso via the shoulder.
    Modeled as a 3-DOF ball joint: $\manifold_2 = \SOthree$.
  \item \textbf{Forearm/hands}: rotates relative to the upper arm via the
    elbow (1~DOF) and wrist flexion (1~DOF): $\manifold_3 = \Sone \times \Sone$.
  \item \textbf{Club}: rotates relative to the hands via wrist
    cock/uncock (1~DOF): $\manifold_4 = \Sone$.
\end{enumerate}

The total configuration space is:
\begin{equation}\label{eq:golf_config}
  \configspace_{\text{golf}} = \SOthree \times \SOthree \times \Sone \times \Sone
  \times \Sone.
\end{equation}
This is a \textbf{9-dimensional manifold}. Its topology is
$\SOthree \times \SOthree \times \Torus^3$---decidedly not $\Reals^9$.

\begin{center}
\begin{tikzpicture}[scale=0.9, every node/.style={font=\small}]
  % Stick figure
  % Ground
  \draw[thick, gray] (-1, 0) -- (4, 0);

  % Torso
  \draw[very thick, chapblue] (1.5, 0.3) -- (1.5, 2.5);
  \filldraw[chapblue] (1.5, 2.5) circle (2pt);
  \node[right, chapblue, font=\footnotesize] at (1.8, 1.4) {torso};

  % Shoulder joint
  \draw[fill=accentred, accentred] (1.5, 2.5) circle (5pt);
  \node[above right, accentred, font=\footnotesize\bfseries] at (1.7, 2.6)
    {$\SOthree$};

  % Upper arm
  \draw[very thick, trajgreen] (1.5, 2.5) -- (0.3, 1.5);
  \filldraw[trajgreen] (0.3, 1.5) circle (2pt);
  \node[left, trajgreen, font=\footnotesize] at (0.1, 2.0) {arm};

  % Elbow joint
  \draw[fill=darkgold, darkgold] (0.3, 1.5) circle (4pt);
  \node[left, darkgold, font=\footnotesize\bfseries] at (0.0, 1.3) {$\Sone$};

  % Forearm
  \draw[very thick, purplehaze] (0.3, 1.5) -- (-0.3, 0.3);
  \node[left, purplehaze, font=\footnotesize] at (-0.3, 0.9) {forearm};

  % Wrist joint
  \draw[fill=darkgold, darkgold] (-0.3, 0.3) circle (4pt);
  \node[left, darkgold, font=\footnotesize\bfseries] at (-0.6, 0.1) {$\Sone \times \Sone$};

  % Club
  \draw[very thick, black!70] (-0.3, 0.3) -- (-1.2, -0.8);
  \filldraw[black] (-1.2, -0.8) circle (3pt);
  \node[below, font=\footnotesize] at (-1.2, -1.0) {clubhead};

  % Spine joint
  \draw[fill=accentred, accentred] (1.5, 0.3) circle (5pt);
  \node[right, accentred, font=\footnotesize\bfseries] at (1.8, 0.3) {$\SOthree$};

  % Configuration space annotation
  \draw[thick, gray, rounded corners, dashed]
    (3.2, -0.5) rectangle (7.2, 3.2);
  \node[font=\bfseries, chapblue] at (5.2, 2.8) {Configuration Space};
  \node[font=\normalsize, align=left] at (5.2, 1.8)
    {$\configspace = \SOthree \times \SOthree$};
  \node[font=\normalsize, align=left] at (5.2, 1.1)
    {$\phantom{\configspace = {}} \times\;\Sone \times \Sone \times \Sone$};
  \node[font=\small, align=left] at (5.2, 0.2)
    {dim $= 3 + 3 + 1 + 1 + 1 = 9$};
  \node[font=\small, accentred, align=center] at (5.2, -0.4)
    {\textbf{Not $\Reals^9$!}};
\end{tikzpicture}
\end{center}

\subsection{Why This Matters for Control}

The state space (tangent bundle) has dimension $2 \times 9 = 18$. A trajectory
through this space is a curve on an 18-dimensional manifold. The consequences
for control are immediate:

\begin{enumerate}[label=\textbf{C\arabic*.}, leftmargin=1.5cm]
  \item \textbf{No global linearization.} You cannot write
    $\state = \state^* + \delta\state$ with $\delta\state \in \Reals^{18}$,
    because the ``$+$'' operation is not defined globally on the manifold.
    Linearization must be done in the \emph{tangent space}, using the
    exponential map.

  \item \textbf{Error is not a vector.} The ``error'' between two configurations
    $\bm{q}_1, \bm{q}_2 \in \configspace$ is not $\bm{q}_2 - \bm{q}_1$
    (subtraction is undefined on manifolds). The correct notion of error
    is the \emph{logarithmic map}: $\bm{e} = \log(\bm{q}_1^{-1} \bm{q}_2)$,
    which lives in the Lie algebra.

  \item \textbf{Geodesics replace straight lines.} The ``shortest path'' between
    two configurations is not a straight line in any coordinate system---it
    is a \emph{geodesic} on the manifold, determined by the Riemannian metric
    (mass-inertia matrix).

  \item \textbf{Parallel transport replaces vector addition.} To compare
    tangent vectors (velocities, errors) at different points on the trajectory,
    we must \emph{transport} them along the manifold using the connection.
    Simply subtracting velocity vectors at different times is mathematically
    meaningless.
\end{enumerate}

\begin{warning}{Euler Angle Catastrophe in Golf Modeling}{euler_golf}
  Many biomechanics studies parameterize the golf swing using Euler angles for
  each joint. For the shoulder joint ($\SOthree$), this means three angles
  $(\phi_s, \theta_s, \psi_s)$. But during the backswing, the shoulder
  passes through configurations near gimbal lock ($\theta_s \approx \pm\pi/2$).

  In these regions:
  \begin{itemize}[leftmargin=1cm]
    \item The Jacobian mapping angular velocities to Euler angle rates becomes
      singular.
    \item Small physical rotations produce large jumps in angle coordinates.
    \item Any controller computing torques from angle errors will produce
      \emph{infinite} commanded torques at the singularity.
  \end{itemize}
  The fix is not to ``choose a better Euler angle convention'' (all conventions
  have singularities somewhere). The fix is to work on $\SOthree$ directly,
  using rotation matrices or quaternions.
\end{warning}

% ══════════════════════════════════════════════
\section{Riemannian Metrics and the Mass Matrix}
% ══════════════════════════════════════════════

A manifold by itself has no notion of distance, angle, or length. To measure
these, we need a \textbf{Riemannian metric}---a smooth choice of inner product
on each tangent space.

\begin{definition}{Riemannian Metric}{riemannian}
  A \textbf{Riemannian metric} on a manifold $\manifold$ is a smooth assignment
  of an inner product $\inner{\cdot}{\cdot}_p$ on each tangent space
  $\tangentspace{p}$. In local coordinates $\bm{q} = (q_1, \ldots, q_n)$, the
  metric is represented by a positive-definite matrix $\bm{G}(\bm{q})$:
  \begin{equation}
    \inner{\bm{v}}{\bm{w}}_{\bm{q}} = \bm{v}^\top \bm{G}(\bm{q})\, \bm{w},
    \qquad \bm{v}, \bm{w} \in T_{\bm{q}}\manifold.
  \end{equation}
\end{definition}

For mechanical systems, nature provides a canonical Riemannian metric: the
\textbf{kinetic energy metric}.

\begin{principle}{The Mass Matrix as Riemannian Metric}{mass_metric}
  For a mechanical system with generalized coordinates $\bm{q}$ and
  mass-inertia matrix $\bm{M}(\bm{q})$, the kinetic energy
  \begin{equation}
    T = \frac{1}{2}\dot{\bm{q}}^\top \bm{M}(\bm{q})\, \dot{\bm{q}}
  \end{equation}
  defines a Riemannian metric $\bm{G} = \bm{M}(\bm{q})$ on $\configspace$.
  This metric has deep physical significance:
  \begin{enumerate}[label=(\roman*)]
    \item \textbf{Distance} between configurations is measured in ``units of
      kinetic energy''---configurations connected by high-inertia motions
      are ``far apart.''
    \item \textbf{Geodesics} (shortest paths) under this metric are the
      trajectories of the \emph{unforced} system---free motions with no
      applied torques.
    \item \textbf{Arc length} from Chapter~2, computed with this metric,
      naturally weights each DOF by its effective inertia.
  \end{enumerate}
\end{principle}

\begin{example}{Inertia-Weighted Distance in the Swing}{inertia_distance}
  Consider two configurations that differ by either:
  \begin{enumerate}[label=(\alph*)]
    \item A $10^\circ$ rotation of the torso (high inertia, $I \approx 12\;\text{kg}\cdot\text{m}^2$), or
    \item A $10^\circ$ rotation of the wrist (low inertia, $I \approx 0.03\;\text{kg}\cdot\text{m}^2$).
  \end{enumerate}
  In Euclidean coordinates, both represent ``$10^\circ$ of change.'' In the
  kinetic energy metric, the torso rotation is $\sqrt{12/0.03} = 20$ times
  ``farther'' than the wrist rotation. This correctly reflects the physical
  reality: rotating the torso requires 400 times more energy than rotating
  the wrist by the same angle.

  The curvature of a trajectory measured in this metric therefore naturally
  captures the \emph{physical} difficulty of the motion, not merely the
  angular change.
\end{example}

\subsection{Geodesics as Natural Trajectories}

\begin{definition}{Geodesic}{geodesic}
  A \textbf{geodesic} on a Riemannian manifold $(\manifold, \bm{G})$ is a
  curve that locally minimizes arc length. Geodesics satisfy the
  \textbf{geodesic equation}:
  \begin{equation}\label{eq:geodesic}
    \ddot{q}^i + \sum_{j,k} \Gamma^i_{jk}(\bm{q})\, \dot{q}^j\, \dot{q}^k = 0,
  \end{equation}
  where $\Gamma^i_{jk}$ are the \textbf{Christoffel symbols} of the metric
  $\bm{G}$.
\end{definition}

The remarkable fact is that the geodesic equation \eqref{eq:geodesic} is
\emph{exactly} the equation of motion for the unforced mechanical system
$\bm{M}(\bm{q})\ddot{\bm{q}} + \bm{C}(\bm{q}, \dot{\bm{q}})\dot{\bm{q}} = \bm{0}$.
The Coriolis and centrifugal terms are precisely the Christoffel symbol terms.
Free motions of mechanical systems are geodesics on the configuration manifold
equipped with the kinetic energy metric.

\begin{principle}{Geodesics and Passive Dynamics}{geodesic_passive}
  For trajectory design, geodesics represent the ``easiest'' paths---the
  motions the system would naturally follow with no applied forces. A
  well-designed trajectory stays close to geodesics where possible,
  deviating only where task requirements (e.g., impact conditions) demand
  it.

  In the golf swing, the follow-through is nearly geodesic: after impact,
  the golfer relaxes and the body follows its natural, inertia-weighted
  path through configuration space. The downswing deviates significantly
  from the geodesic because the task (delivering maximum clubhead speed at
  impact) requires active torques to redirect the motion.
\end{principle}

% ══════════════════════════════════════════════
\section{Curves on Manifolds Revisited}
% ══════════════════════════════════════════════

With the manifold framework in place, we can extend every concept from
Chapter~2 to non-Euclidean configuration spaces.

\subsection{Covariant Differentiation}

On $\Reals^n$, we differentiate vector fields by taking component-wise
derivatives. On a manifold, this does not work---the tangent spaces at
different points are different vector spaces, so we cannot simply subtract
vectors at different points.

\begin{definition}{Covariant Derivative}{covariant}
  The \textbf{covariant derivative} (or \textbf{Levi-Civita connection})
  $\nabla$ is the unique way to differentiate vector fields along curves on
  a Riemannian manifold that:
  \begin{enumerate}[label=(\roman*)]
    \item Is compatible with the metric:
      $\frac{\dd}{\dd t}\inner{\bm{V}}{\bm{W}} =
      \inner{\nabla_{\dot{\gamma}} \bm{V}}{\bm{W}} +
      \inner{\bm{V}}{\nabla_{\dot{\gamma}} \bm{W}}$.
    \item Is torsion-free: $\nabla_{\bm{V}}\bm{W} - \nabla_{\bm{W}}\bm{V}
      = [\bm{V}, \bm{W}]$.
  \end{enumerate}
  In local coordinates, the covariant derivative of a vector field
  $\bm{V}$ along a curve $\gamma(t)$ is
  \begin{equation}
    (\nabla_{\dot{\gamma}} \bm{V})^i
    = \dot{V}^i + \sum_{j,k} \Gamma^i_{jk}\, \dot{\gamma}^j\, V^k.
  \end{equation}
\end{definition}

The covariant derivative replaces the ordinary derivative in all of our
Chapter~2 formulas. The generalized curvatures, the Frenet--Serret frame,
and the transverse decomposition all carry over, with $\dd/\dd s$ replaced
by $\nabla_{\dot{\gamma}}$.

\subsection{Manifold Curvature of Trajectories}

\begin{definition}{Geodesic Curvature}{geodesic_curvature}
  The \textbf{geodesic curvature} of a curve $\gamma(s)$ on a Riemannian
  manifold, parameterized by arc length, is
  \begin{equation}
    \kappa_g(s) = \norm{\nabla_{\dot{\gamma}} \dot{\gamma}}.
  \end{equation}
  This measures how much the curve deviates from a geodesic. A geodesic
  has $\kappa_g = 0$ everywhere.
\end{definition}

Geodesic curvature is the manifold generalization of the curvature from
Chapter~2, and it has the same control interpretation: \emph{geodesic
curvature measures the force per unit mass needed to follow the curve},
beyond what the natural dynamics would provide for free.

\begin{proposition}
  The control force required to track a trajectory $\gamma^*(s)$ at speed
  $v$ on a configuration manifold with metric $\bm{M}(\bm{q})$ has magnitude
  \begin{equation}
    \norm{\bm{\tau}_{\perp}} = \kappa_g(s)\, v^2,
  \end{equation}
  where $\norm{\cdot}$ is measured in the dual (force) metric. Comparing with
  the Euclidean formula $a_\perp = \kappa v^2$ from Chapter~2: the Euclidean
  curvature $\kappa$ is replaced by the geodesic curvature $\kappa_g$, which
  accounts for the curvature of the manifold itself.
\end{proposition}

This is a deeper version of the Curvature--Authority Principle. The control
demand is reduced along directions where the manifold's own curvature
``helps''---i.e., where the natural dynamics carry the system in the desired
direction. The control demand increases where the trajectory fights the
natural dynamics.

\begin{example}{Geodesic Curvature of the Downswing}{downswing_geod_curv}
  During the downswing, the trajectory has two sources of curvature:
  \begin{enumerate}[label=(\alph*)]
    \item \textbf{Euclidean curvature}: the path bends in coordinate space
      (the club changes direction).
    \item \textbf{Manifold curvature}: the Christoffel symbol terms
      ($\Gamma^i_{jk}\dot{q}^j\dot{q}^k$) represent Coriolis and centrifugal
      effects that curve the natural trajectories.
  \end{enumerate}
  The geodesic curvature $\kappa_g$ is the \emph{difference}. When the
  Coriolis/centrifugal terms push the system in the direction the trajectory
  wants to go (as happens during the wrist release, where centrifugal
  acceleration ``unhinges'' the club), $\kappa_g$ is \emph{less} than the
  Euclidean $\kappa$. The passive dynamics are doing work for free.

  This is the geometric explanation for the ``free speed'' of the golf swing:
  a well-designed trajectory aligns with the manifold's geodesics during the
  speed-critical phase, exploiting the natural dynamics rather than fighting
  them.
\end{example}

% ══════════════════════════════════════════════
\section{Error Metrics on Manifolds}
% ══════════════════════════════════════════════

We now address the critical question: how do we measure tracking error when
the state space is a manifold?

\subsection{The Logarithmic Error}

\begin{definition}{Configuration Error on a Lie Group}{log_error}
  For two configurations $\bm{q}_1, \bm{q}_2 \in \configspace$ where
  $\configspace$ is a Lie group, the \textbf{error} is
  \begin{equation}\label{eq:log_error}
    \bm{e} = \log(\bm{q}_1^{-1}\, \bm{q}_2) \in \text{Lie algebra}.
  \end{equation}
  For the rotation group $\SOthree$, this gives
  \begin{equation}
    \bm{e}_R = \log(\bm{R}_1^\top \bm{R}_2)^\vee \in \Reals^3,
  \end{equation}
  which is the rotation vector needed to go from $\bm{R}_1$ to $\bm{R}_2$.
  This error is zero if and only if $\bm{q}_1 = \bm{q}_2$, and it varies
  smoothly with both arguments (away from the cut locus at $\norm{\bm{e}} = \pi$).
\end{definition}

\begin{remark}{Why Not Just Subtract Quaternions?}{quat_subtract}
  Given two unit quaternions $\bm{q}_1$ and $\bm{q}_2$, a tempting error
  metric is $\bm{e} = \bm{q}_2 - \bm{q}_1$. But this lives in $\Reals^4$,
  not in the tangent space of $S^3$. Worse, it violates the double-cover
  property: $\bm{q}$ and $-\bm{q}$ represent the same rotation, so the
  error $\bm{q}_2 - \bm{q}_1$ and $-\bm{q}_2 - \bm{q}_1$ should be
  ``the same,'' but they point in opposite directions.

  The logarithmic error correctly handles the group structure and produces
  an error vector in the Lie algebra that has a clear physical interpretation
  (the angular velocity needed to correct the error in unit time).
\end{remark}

\subsection{The Trajectory Tube on a Manifold}

The trajectory tube from Chapter~1 now generalizes naturally:

\begin{definition}{Manifold Trajectory Tube}{manifold_tube}
  Let $\traj^*: [0, T] \to \configspace$ be a reference trajectory on a
  Riemannian manifold $(\configspace, \bm{G})$ and let
  $\varepsilon: [0, T] \to \Reals_{>0}$ be a tube radius function. The
  \textbf{manifold trajectory tube} is
  \begin{equation}
    \mathcal{T}(t) = \left\{\, \bm{q} \in \configspace :
    d_{\bm{G}}\!\left(\bm{q},\, \traj^*(t)\right) \leq \varepsilon(t) \,\right\},
  \end{equation}
  where $d_{\bm{G}}$ is the \textbf{geodesic distance}---the length of the
  shortest geodesic connecting $\bm{q}$ and $\traj^*(t)$.
\end{definition}

The geodesic distance correctly accounts for the topology of the configuration
space. On $\SOthree$, the maximum distance between any two orientations is $\pi$
(a half-rotation), and the tube wraps properly around the group.

% ══════════════════════════════════════════════
\section{Equations of Motion on Manifolds}
% ══════════════════════════════════════════════

The Euler--Lagrange equations on a manifold take a coordinate-free form that
reveals the geometric structure:

\begin{theorem}[Lagrangian Mechanics on Manifolds]\label{thm:lagrange_manifold}
  For a mechanical system on a configuration manifold $(\configspace, \bm{M})$
  with potential energy $V: \configspace \to \Reals$ and generalized forces
  $\bm{\tau}$, the equations of motion are
  \begin{equation}\label{eq:euler_lagrange_manifold}
    \bm{M}(\bm{q})\, \nabla_{\dot{\bm{q}}} \dot{\bm{q}} = -\mathrm{grad}\, V + \bm{B}\, \control,
  \end{equation}
  where $\nabla_{\dot{\bm{q}}} \dot{\bm{q}}$ is the covariant acceleration and
  $\mathrm{grad}\, V$ is the Riemannian gradient of the potential energy.
\end{theorem}

In local coordinates, this reduces to the familiar
$\bm{M}\ddot{\bm{q}} + \bm{C}\dot{\bm{q}} + \bm{g} = \bm{B}\control$,
but the coordinate-free form \eqref{eq:euler_lagrange_manifold} makes clear that:
\begin{itemize}[leftmargin=1.5cm]
  \item The Coriolis matrix $\bm{C}$ is not a separate ``term''---it is the
    Christoffel symbol contribution to the covariant derivative.
  \item The equation has the same form on any manifold, regardless of coordinate
    choice.
  \item The ``centrifugal and Coriolis forces'' are geometric artifacts of the
    manifold's curvature, not real forces. On a flat manifold ($\Reals^n$),
    they vanish.
\end{itemize}

% ══════════════════════════════════════════════
\section{Practical Implications for Controller Implementation}
% ══════════════════════════════════════════════

We conclude with concrete guidance for implementing trajectory controllers
on manifold-valued configuration spaces.

\begin{principle}{Manifold-Aware Controller Design}{manifold_control}
  A trajectory tracking controller for a system on $\configspace$ must:
  \begin{enumerate}[label=(\roman*)]
    \item \textbf{Compute errors using the logarithmic map}, not coordinate
      subtraction:
      \begin{equation}
        \bm{e}_q = \log\!\left((\bm{q}^*)^{-1}\, \bm{q}\right), \qquad
        \bm{e}_v = \dot{\bm{q}} - \mathrm{PT}_{q^* \to q}(\dot{\bm{q}}^*),
      \end{equation}
      where $\mathrm{PT}$ denotes parallel transport of the reference velocity
      to the current configuration.
    \item \textbf{Compute feedforward using the covariant derivative}, not
      ordinary differentiation of coordinates.
    \item \textbf{Design feedback gains in the tangent space}, then map the
      resulting control back to the manifold via the exponential map.
    \item \textbf{Monitor for cut-locus proximity}: when the error approaches
      $\pi$ (for $\SOthree$), the logarithmic map becomes ill-conditioned.
      Switch to a recovery strategy.
  \end{enumerate}
\end{principle}

\begin{example}{PD Control on $\SOthree$}{pd_so3}
  The manifold-correct PD controller for tracking a reference orientation
  $\bm{R}^*(t)$ with a rigid body (moment of inertia $\bm{J}$,
  angular velocity $\bm{\omega}$) is:
  \begin{equation}\label{eq:pd_so3}
    \bm{\tau} = -k_p\, \log(\bm{R}^{*\top}\bm{R})^\vee
    - k_d\, (\bm{\omega} - \bm{R}^{*\top}\bm{R}\,\bm{\omega}^*)
    + \bm{J}\dot{\bm{\omega}}^* + \bm{\omega} \times \bm{J}\bm{\omega},
  \end{equation}
  where the first term is the manifold-correct proportional error
  (logarithmic map), the second is the velocity error (parallel-transported),
  and the last two terms are the feedforward (covariant derivative of the
  reference).

  Compare with the na\"{\i}ve Euler-angle PD controller:
  $\tau = -k_p(\bm{\theta} - \bm{\theta}^*) - k_d(\dot{\bm{\theta}} - \dot{\bm{\theta}}^*)$.
  The na\"{\i}ve controller has singularities; Eq.~\eqref{eq:pd_so3} does not.
  The na\"{\i}ve controller produces torques proportional to angle differences
  (which can be misleading near $\pm\pi$); Eq.~\eqref{eq:pd_so3} produces
  torques proportional to the \emph{true rotation} needed to correct the error.
\end{example}

% ══════════════════════════════════════════════
\section{Summary}
% ══════════════════════════════════════════════

\begin{center}
\renewcommand{\arraystretch}{1.4}
\begin{tabular}{@{} l l @{}}
  \toprule
  \textbf{Euclidean assumption} & \textbf{Manifold replacement} \\
  \midrule
  $\configspace = \Reals^n$ & $\configspace = \SOthree^k \times \Torus^m \times \cdots$ \\
  Error: $\bm{e} = \bm{q} - \bm{q}^*$ & Error: $\bm{e} = \log((\bm{q}^*)^{-1}\bm{q})$ \\
  Derivative: $\ddot{\bm{q}}$ & Covariant derivative: $\nabla_{\dot{q}}\dot{\bm{q}}$ \\
  Shortest path: straight line & Shortest path: geodesic \\
  Distance: $\norm{\bm{q}_2 - \bm{q}_1}$ & Geodesic distance: $d_{\bm{G}}(\bm{q}_1, \bm{q}_2)$ \\
  Curvature: $\kappa$ (Ch.~2) & Geodesic curvature: $\kappa_g$ \\
  Trajectory tube: Euclidean ball & Geodesic ball \\
  Coriolis terms: ``extra forces'' & Christoffel symbols: manifold geometry \\
  PD: $-k_p\bm{e} - k_d\dot{\bm{e}}$ & $-k_p\log(\cdot)^\vee - k_d(\text{parallel transport})$ \\
  \bottomrule
\end{tabular}
\end{center}

\vspace{0.3cm}

The central message of this chapter is that the configuration spaces of
mechanical systems have geometry, and this geometry matters for control.
Working in $\Reals^n$ when the true space is $\SOthree \times \Torus^k$
introduces artificial singularities, incorrect error metrics, and
controllers that fail when the system moves far from a nominal configuration.

The trajectory-first philosophy of this book demands that we take the
geometry seriously from the start. The curves we designed in Chapter~2 live
on manifolds, and the stability theory we will develop in Chapter~4 must
respect the manifold structure. The payoff is controllers that work globally,
without singularities, and that correctly exploit the natural dynamics
encoded in the manifold's Riemannian geometry.

% ══════════════════════════════════════════════
\section{Exercises}
% ══════════════════════════════════════════════

\begin{enumerate}[label=\textbf{3.\arabic*.}, leftmargin=1.5cm]

\item \textbf{(Topology.)} Determine the configuration manifold and its
  dimension for each system:
  \begin{enumerate}[label=(\alph*)]
    \item A planar double pendulum (two revolute joints in 2D).
    \item A spherical pendulum (a point mass on a rigid rod, free to swing
      in all directions).
    \item A spacecraft with two reaction wheels (the body rotates freely in
      3D; each wheel adds one internal DOF).
  \end{enumerate}
  For each, identify whether $\Reals^n$ coordinates would have singularities.

\item \textbf{(Euler angle singularity.)} For the ZYX Euler angle representation
  $\bm{R} = R_z(\phi)\, R_y(\theta)\, R_x(\psi)$:
  \begin{enumerate}[label=(\alph*)]
    \item Compute the Jacobian $\bm{J}$ relating angular velocity
      $\bm{\omega}$ to Euler angle rates $(\dot{\phi}, \dot{\theta}, \dot{\psi})$.
    \item Show that $\bm{J}$ is singular at $\theta = \pm\pi/2$.
    \item Find an orientation trajectory $\bm{R}(t)$ that passes through the
      singularity and compute the Euler angle rates. What happens?
  \end{enumerate}

\item \textbf{(Logarithmic error.)} For two rotations
  $\bm{R}_1 = \bm{I}$ and $\bm{R}_2 = R_z(\alpha)$ (rotation by $\alpha$
  about $z$):
  \begin{enumerate}[label=(\alph*)]
    \item Compute $\log(\bm{R}_1^\top\bm{R}_2)^\vee$ for
      $\alpha = \pi/6, \pi/2, \pi, 5\pi/3$.
    \item Compare with the ``error'' $\alpha - 0 = \alpha$ as a function of
      $\alpha$. At what point do they differ significantly?
    \item What happens at $\alpha = \pi$? Explain geometrically.
  \end{enumerate}

\item \textbf{(Geodesics.)} For a double pendulum with masses $m_1, m_2$ and
  lengths $l_1, l_2$:
  \begin{enumerate}[label=(\alph*)]
    \item Write the mass matrix $\bm{M}(\theta_1, \theta_2)$.
    \item Compute the Christoffel symbols (or use the Coriolis matrix formula
      $C_{ijk} = \frac{1}{2}(\partial M_{ij}/\partial q_k +
      \partial M_{ik}/\partial q_j - \partial M_{jk}/\partial q_i)$).
    \item Numerically integrate the geodesic equation from initial condition
      $(\theta_1, \theta_2) = (0, 0)$ with $(\dot{\theta}_1, \dot{\theta}_2) = (1, 0)$.
      Plot the geodesic on $\Torus^2$.
  \end{enumerate}

\item \textbf{(Manifold PD control.)} Implement and compare:
  \begin{enumerate}[label=(\alph*)]
    \item An Euler-angle PD controller for a rigid body tracking a reference
      orientation $\bm{R}^*(t) = R_z(\omega t)$ (constant-speed rotation).
    \item The manifold PD controller from Eq.~\eqref{eq:pd_so3} for the
      same task.
    \item Run both controllers with the reference trajectory passing through
      $\theta = \pi/2$ (gimbal lock for ZYX convention). Compare the
      tracking errors and commanded torques.
  \end{enumerate}

\item \textbf{(Golf swing configuration space.)} For the 9-DOF golf swing
  model of Eq.~\eqref{eq:golf_config}:
  \begin{enumerate}[label=(\alph*)]
    \item How many Euler angles would be needed to parameterize the full
      configuration space? How many singularity surfaces exist?
    \item Propose a quaternion-based parameterization for the two $\SOthree$
      joints. What is the total number of parameters? What constraints
      must be maintained?
    \item Discuss the trade-offs between the Euler angle and quaternion
      representations for numerical simulation of the golf swing.
  \end{enumerate}

\end{enumerate}

\vspace{1cm}
\begin{center}
  \rule{0.5\textwidth}{0.4pt}\\[0.5cm]
  {\small\itshape
  The sphere cannot be flattened without tearing.\\
  The rotation group cannot be charted without singularity.\\
  Do not fight the topology---ride it.\\[0.3cm]
  The manifold is not an obstacle. It is the landscape\\
  through which every motion flows.}
\end{center}

\chapter{Orbital Stability and Transverse Linearization}

\epigraph{%
  Stability is not the absence of motion.\\
  It is the faithfulness of motion.
}{}

\vspace{0.5cm}

% ══════════════════════════════════════════════
\section{The Problem with Lyapunov}
% ══════════════════════════════════════════════

The classical Lyapunov framework asks a simple question: \emph{does the state
converge to an equilibrium?} We demonstrated in Chapter~1 that this is the
wrong question for systems whose purpose is to move. But the problem goes
deeper than philosophical framing---Lyapunov stability is \emph{structurally
incapable} of analyzing trajectory-tracking systems.

Consider a system perfectly tracking a reference trajectory $\traj^*(t)$.
The state $\state(t)$ matches $\traj^*(t)$ at every instant. Now apply a
small timing perturbation: the system executes the same motion but shifted
by $\varepsilon$ seconds. The perturbed state is $\tilde{\state}(t) =
\traj^*(t - \varepsilon)$, and the time-indexed error is
\begin{equation}\label{eq:timing_error}
  \bm{e}(t) = \tilde{\state}(t) - \traj^*(t)
  = \traj^*(t - \varepsilon) - \traj^*(t)
  \approx -\varepsilon\, \dot{\traj}^*(t).
\end{equation}
This error does \emph{not} decay to zero---it persists indefinitely,
oscillating with a magnitude proportional to $\varepsilon\, v(t)$ where $v$
is the speed. By Lyapunov's definition, the trajectory is \emph{unstable}.
But the system is executing the correct motion perfectly! The only
``error'' is that it is slightly ahead of or behind schedule.

\begin{keyidea}{The Fundamental Inadequacy of Point Stability}{lyap_failure}
  Lyapunov stability conflates two independent phenomena:
  \begin{enumerate}[label=(\roman*)]
    \item \textbf{Transverse deviations}: the system drifts \emph{away from}
      the trajectory path. This is a genuine control failure.
    \item \textbf{Tangential deviations}: the system is ahead of or behind
      schedule \emph{along} the trajectory path. This is a timing difference,
      not a tracking failure.
  \end{enumerate}
  Lyapunov stability rejects both. Orbital stability accepts tangential
  deviations and rejects only transverse ones. For systems whose purpose is
  motion, orbital stability is the correct notion.
\end{keyidea}

\begin{center}
\begin{tikzpicture}[scale=1.0]
  % --- Lyapunov view ---
  \begin{scope}[shift={(-4.5,0)}]
    \node[font=\bfseries\small, chapblue] at (0, 3.0) {Lyapunov View};
    \draw[-{Stealth[length=2.5mm]}, thick, gray] (-2, 0) -- (2, 0)
      node[right, font=\small] {$x_1$};
    \draw[-{Stealth[length=2.5mm]}, thick, gray] (0, -2) -- (0, 2.5)
      node[above, font=\small] {$x_2$};
    % Reference trajectory
    \draw[very thick, trajgreen, decoration={markings,
        mark=at position 0.3 with {\arrow{Stealth[length=2.5mm]}},
        mark=at position 0.7 with {\arrow{Stealth[length=2.5mm]}}
      }, postaction={decorate}]
      (-1.5, -1.5) .. controls (-0.5, 1) and (0.5, 1.5) .. (1.5, 0.3);
    % Reference point at time t
    \filldraw[trajgreen!80!black] (0.0, 1.15) circle (2.5pt);
    \node[above right, font=\scriptsize, trajgreen!80!black] at (0.1, 1.2) {$\traj^*(t)$};
    % Perturbed (timing shift)
    \filldraw[accentred] (-0.55, 0.5) circle (2.5pt);
    \node[below left, font=\scriptsize, accentred] at (-0.6, 0.4) {$\traj^*(t\!-\!\varepsilon)$};
    % Error arrow
    \draw[-{Stealth[length=2mm]}, thick, accentred, dashed]
      (0.0, 1.15) -- (-0.5, 0.55);
    \node[font=\scriptsize, accentred] at (0.5, 0.55) {$\bm{e}(t) \neq 0$};
    % Verdict
    \node[font=\small\bfseries, accentred] at (0, -2.5) {``Unstable''};
    \node[font=\scriptsize] at (0, -2.9) {(wrong answer)};
  \end{scope}

  % Dividing line
  \draw[thick, gray!30] (0, -2.2) -- (0, 2.8);

  % --- Orbital view ---
  \begin{scope}[shift={(4.5,0)}]
    \node[font=\bfseries\small, trajgreen!80!black] at (0, 3.0) {Orbital View};
    \draw[-{Stealth[length=2.5mm]}, thick, gray] (-2, 0) -- (2, 0)
      node[right, font=\small] {$x_1$};
    \draw[-{Stealth[length=2.5mm]}, thick, gray] (0, -2) -- (0, 2.5)
      node[above, font=\small] {$x_2$};
    % Reference trajectory
    \draw[very thick, trajgreen, decoration={markings,
        mark=at position 0.3 with {\arrow{Stealth[length=2.5mm]}},
        mark=at position 0.7 with {\arrow{Stealth[length=2.5mm]}}
      }, postaction={decorate}]
      (-1.5, -1.5) .. controls (-0.5, 1) and (0.5, 1.5) .. (1.5, 0.3);
    % Perturbed (timing shift) - ON the trajectory
    \filldraw[chapblue] (-0.55, 0.5) circle (2.5pt);
    \node[below left, font=\scriptsize, chapblue] at (-0.6, 0.4) {on $\orbit$};
    % Transverse distance = 0
    \node[font=\scriptsize, chapblue] at (1.0, 0.9) {$d_\transverse = 0$};
    % Verdict
    \node[font=\small\bfseries, trajgreen!80!black] at (0, -2.5) {``Stable''};
    \node[font=\scriptsize] at (0, -2.9) {(correct answer)};
  \end{scope}
\end{tikzpicture}
\end{center}

% ══════════════════════════════════════════════
\section{Orbital Stability: Formal Definitions}
% ══════════════════════════════════════════════

We now define orbital stability rigorously. The key object is the
\textbf{orbit}---the image of the trajectory as a set, stripped of all
timing information.

\begin{definition}{Orbit}{orbit}
  Given a trajectory $\traj^*: [0, T] \to \statespace$, the \textbf{orbit}
  is the set
  \begin{equation}
    \orbit = \{\traj^*(t) : t \in [0, T]\} \subset \statespace.
  \end{equation}
  The orbit is a one-dimensional submanifold of $\statespace$ (a curve
  without parameterization). For periodic trajectories with period $T$,
  $\traj^*(0) = \traj^*(T)$ and $\orbit$ is a closed curve (a loop).
\end{definition}

\begin{definition}{Orbital Stability}{orbital_stability}
  A trajectory $\traj^*(t)$ (or its orbit $\orbit$) is \textbf{orbitally
  stable} if for every $\varepsilon > 0$ there exists $\delta > 0$ such that
  \begin{equation}\label{eq:orbital_stability}
    d(\state(0), \orbit) < \delta \implies
    d(\state(t), \orbit) < \varepsilon \quad \forall\, t \geq 0,
  \end{equation}
  where $d(\state, \orbit) = \inf_{p \in \orbit} \norm{\state - p}$ is the
  distance from the state to the orbit.

  It is \textbf{asymptotically orbitally stable} if additionally
  $d(\state(t), \orbit) \to 0$ as $t \to \infty$.
\end{definition}

Critically, orbital stability does not require $\norm{\state(t) - \traj^*(t)}
\to 0$. The system may drift along the orbit (ahead or behind schedule) while
remaining perfectly stable. Only deviations \emph{away from} the orbit matter.

\begin{definition}{Orbital Stability with Asymptotic Phase}{orbital_phase}
  A trajectory is orbitally stable \textbf{with asymptotic phase} if it is
  asymptotically orbitally stable and additionally there exists a phase shift
  $\theta^*$ such that
  \begin{equation}
    \norm{\state(t) - \traj^*(t + \theta^*)} \to 0 \quad \text{as } t \to \infty.
  \end{equation}
  This means the system not only returns to the orbit but eventually
  ``locks in'' to a definite timing, shifted by $\theta^*$ from the
  nominal schedule. This is the strongest practical form of trajectory
  stability.
\end{definition}

\begin{example}{Limit Cycles: The Prototype}{limit_cycle}
  The Van der Pol oscillator $\ddot{x} - \mu(1 - x^2)\dot{x} + x = 0$
  possesses a \textbf{limit cycle}---a periodic orbit that is asymptotically
  orbitally stable with asymptotic phase. All nearby trajectories spiral
  toward the limit cycle and eventually synchronize to a definite phase.

  The limit cycle is the simplest example of a system whose purpose is to
  move: the ``target'' is not a point but an orbit, and stability means
  the orbit attracts nearby trajectories. Walking, running, and heartbeats
  are all limit cycles. The golf swing is a transient trajectory (not
  periodic), but the same stability concepts apply over the finite interval
  $[0, T]$.
\end{example}

% ══════════════════════════════════════════════
\section{The Transverse--Tangential Decomposition}
% ══════════════════════════════════════════════

To analyze orbital stability, we must decompose the dynamics into components
along and across the orbit. This is the \emph{transverse--tangential
decomposition}, the central technical tool of this chapter.

\subsection{Projection onto the Orbit}

\begin{definition}{Nearest-Point Projection}{projection}
  Given a state $\state$ sufficiently close to the orbit $\orbit$, the
  \textbf{nearest-point projection} is
  \begin{equation}\label{eq:projection}
    \pi(\state) = \arg\min_{p \in \orbit} \norm{\state - p},
  \end{equation}
  and the associated \textbf{projected phase} is the unique $s^*(\state)$
  satisfying $\pi(\state) = \traj^*(s^*)$. The projection is well-defined
  and smooth in a tubular neighborhood of $\orbit$ (as long as $\state$
  is closer to $\orbit$ than the minimum radius of curvature of the orbit).
\end{definition}

\subsection{The Moving Hyperplane Decomposition}

At each point $\traj^*(s)$ on the orbit, the state space $\Reals^n$
decomposes into a one-dimensional tangential direction and an
$(n{-}1)$-dimensional transverse hyperplane:

\begin{definition}{Transverse Hyperplane}{transverse_hyperplane}
  At the point $\traj^*(s)$ on the orbit, define:
  \begin{align}
    \text{Tangential direction:} &\quad
      \bm{t}(s) = \frac{\dot{\traj}^*(s)}{\norm{\dot{\traj}^*(s)}}, \\
    \text{Transverse hyperplane:} &\quad
      \Sigma(s) = \{\bm{v} \in \Reals^n : \bm{v}^\top \bm{t}(s) = 0\}.
  \end{align}
  A state $\state$ near $\traj^*(s)$ decomposes as
  \begin{equation}\label{eq:decomposition}
    \state = \traj^*(s^*) + \bm{\xi}, \qquad
    \bm{\xi} \in \Sigma(s^*),
  \end{equation}
  where $s^*$ is the projected phase and $\bm{\xi}$ is the \textbf{transverse
  deviation}---the component of the error perpendicular to the orbit. Orbital
  stability is equivalent to the stability of the origin $\bm{\xi} = \bm{0}$
  in the transverse dynamics.
\end{definition}

\begin{center}
\begin{tikzpicture}[scale=1.15]
  % Orbit
  \draw[very thick, trajgreen, decoration={markings,
      mark=at position 0.15 with {\arrow{Stealth[length=3mm]}},
      mark=at position 0.5 with {\arrow{Stealth[length=3mm]}},
      mark=at position 0.85 with {\arrow{Stealth[length=3mm]}}
    }, postaction={decorate}]
    (-3, -0.5) .. controls (-1.5, 2) and (1.5, 2) .. (3, -0.3);
  \node[trajgreen!70!black, font=\small\bfseries] at (3.4, 0.1) {$\orbit$};

  % Point on orbit
  \coordinate (P) at (0, 1.65);
  \filldraw[trajgreen!70!black] (P) circle (2.5pt);
  \node[above, font=\small, trajgreen!70!black] at (0, 1.85) {$\traj^*(s^*)$};

  % Tangent vector
  \draw[-{Stealth[length=3.5mm]}, very thick, chapblue]
    (P) -- ++(1.2, -0.05);
  \node[chapblue, font=\small\bfseries] at (1.5, 1.85) {$\bm{t}(s^*)$};

  % Transverse hyperplane (shown as a line in 2D)
  \draw[thick, darkgold, dashed] ($(P) + (-0.1, -1.8)$) -- ($(P) + (0.1, 1.8)$);
  \node[darkgold, font=\small\bfseries, rotate=85] at (-0.5, 2.8) {$\Sigma(s^*)$};

  % Actual state
  \coordinate (X) at (0.3, 0.2);
  \filldraw[accentred] (X) circle (2.5pt);
  \node[right, font=\small\bfseries, accentred] at (0.4, 0.2) {$\state(t)$};

  % Transverse deviation
  \draw[-{Stealth[length=2.5mm]}, thick, accentred]
    (P) -- ($(P) + (0.05, -1.45)$);
  \node[accentred, font=\small] at (-0.7, 0.5) {$\bm{\xi}$};

  % Phase annotation
  \draw[thin, gray, |->] (-3, -1.6) -- (3, -1.6) node[right, font=\small] {$s$};
  \draw[thin, gray, dashed] (0, -1.6) -- (P);
  \filldraw[gray] (0, -1.6) circle (1.5pt);
  \node[font=\small, gray, below] at (0, -1.6) {$s^*$};
\end{tikzpicture}
\end{center}

% ══════════════════════════════════════════════
\section{Transverse Linearization}
% ══════════════════════════════════════════════

The power of the transverse decomposition is that it reduces orbital
stability to the stability of a \emph{time-varying linear system}---the
transverse linearization.

\subsection{Deriving the Transverse Dynamics}

Consider the nonlinear system $\dot{\state} = \bm{f}(\state, \control)$
with a reference trajectory $\traj^*(t)$ satisfying
$\dot{\traj}^* = \bm{f}(\traj^*, \control^*(t))$. Using the decomposition
$\state = \traj^*(s^*) + \bm{\xi}$ and linearizing about $\bm{\xi} = \bm{0}$:

\begin{theorem}[Transverse Linearization~\cite{Shiriaev2010, ManchesterSlotine2017}]\label{thm:transverse_lin}
  The linearized transverse dynamics about the orbit $\orbit$ are
  \begin{equation}\label{eq:transverse_linear}
    \dot{\bm{\xi}} = \Amat_\transverse(s)\, \bm{\xi}
    + \Bmat_\transverse(s)\, \delta\control,
  \end{equation}
  where $\delta\control = \control - \control^*$ is the feedback perturbation, and
  \begin{align}
    \Amat_\transverse(s) &= \bm{P}(s) \left[
      \frac{\partial \bm{f}}{\partial \state}\bigg|_{\traj^*, \control^*}
      - \frac{\dot{\traj}^* \otimes (\nabla_{\state} s^*)^\top}
      {\norm{\dot{\traj}^*}^2} \cdot
      \frac{\partial \bm{f}}{\partial \state}\bigg|_{\traj^*, \control^*}
    \right] \bm{P}(s), \label{eq:A_transverse} \\[6pt]
    \Bmat_\transverse(s) &= \bm{P}(s)\,
      \frac{\partial \bm{f}}{\partial \control}\bigg|_{\traj^*, \control^*},
      \label{eq:B_transverse}
  \end{align}
  and $\bm{P}(s) = \bm{I} - \bm{t}(s)\bm{t}(s)^\top$ is the orthogonal
  projector onto $\Sigma(s)$.
\end{theorem}

\begin{remark}{Understanding the Transverse $\Amat_\transverse$}{A_trans_explain}
  Equation \eqref{eq:A_transverse} has a clear structure:
  \begin{itemize}[leftmargin=1cm]
    \item The first term $\bm{P}\, \frac{\partial \bm{f}}{\partial \state}\, \bm{P}$
      is the standard Jacobian projected onto the transverse space. This
      captures how the system dynamics affect cross-trajectory deviations.
    \item The second term subtracts the component of the Jacobian that acts along
      the tangent direction---this is exactly the tangential dynamics that
      orbital stability ``forgets.''
  \end{itemize}
  The resulting $\Amat_\transverse(s)$ has dimension $(n{-}1) \times (n{-}1)$,
  one less than the full state Jacobian. The removed eigenvalue corresponds
  to motion along the orbit---the tangential direction in which we are
  \emph{not} demanding stability.
\end{remark}

\begin{principle}{The Transverse Reduction Principle}{transverse_reduction}
  The reference trajectory $\traj^*(t)$ is asymptotically orbitally stable
  if and only if the origin of the transverse linearization
  \eqref{eq:transverse_linear} (with $\delta\control = \bm{0}$) is
  asymptotically stable. This reduces the orbital stability question from
  a nonlinear problem on the full $n$-dimensional state space to a
  \emph{linear} problem on the $(n{-}1)$-dimensional transverse space.

  For controller design: if we can stabilize the linear system
  \eqref{eq:transverse_linear}, we achieve orbital stability of the
  nonlinear system. All the tools of linear time-varying (LTV) control
  theory apply.
\end{principle}

\subsection{The Transverse Jacobian: A Practical Formula}

For implementation, the key quantities are computed as follows. Let
$\bm{J}(t) = \frac{\partial \bm{f}}{\partial \state}\big|_{\traj^*, \control^*}$
be the Jacobian along the reference. Choose any orthonormal basis
$\{\bm{n}_1(s), \ldots, \bm{n}_{n-1}(s)\}$ for $\Sigma(s)$---the Frenet--Serret
frame from Chapter~2 provides a natural choice. Collect these into the matrix
$\bm{N}(s) = [\bm{n}_1 \mid \cdots \mid \bm{n}_{n-1}] \in \Reals^{n \times (n-1)}$.
Then:

\begin{equation}\label{eq:transverse_practical}
  \Amat_\transverse(s) = \bm{N}(s)^\top \left[\bm{J}(s)
  - \frac{\bm{f}^* \cdot \bm{J}(s)^\top \bm{t}(s)}{\norm{\bm{f}^*}^2}\, \bm{f}^*\, \bm{t}(s)^\top
  - \dot{\bm{N}}(s)\, \bm{N}(s)^\top \right] \bm{N}(s),
\end{equation}
where $\bm{f}^* = \bm{f}(\traj^*, \control^*)$ and $\dot{\bm{N}}$ accounts
for the rotation of the transverse frame along the orbit (the Frenet--Serret
equations).

\begin{example}{Transverse Linearization of a Planar System}{planar_transverse}
  Consider the system $\dot{x}_1 = x_2$, $\dot{x}_2 = -x_1 + u$ tracking
  the circular orbit $\traj^*(t) = (\cos t,\, -\sin t)$ with $u^* = 0$.

  The Jacobian is $\bm{J} = \bigl(\begin{smallmatrix} 0 & 1 \\ -1 & 0
  \end{smallmatrix}\bigr)$, the tangent vector is
  $\bm{t} = (-\sin t,\, -\cos t)^\top / 1 = (-\sin t,\, -\cos t)^\top$,
  and the normal is $\bm{n} = (-\cos t,\, \sin t)^\top$.

  Computing $\Amat_\transverse = \bm{n}^\top[\bm{J} - \text{tangential terms}]\bm{n}$,
  we find:
  \begin{equation}
    \Amat_\transverse(t) = 0, \qquad \Bmat_\transverse(t) = \bm{n}^\top
    \begin{pmatrix} 0 \\ 1 \end{pmatrix} = \sin t.
  \end{equation}
  The open-loop transverse dynamics are \emph{neutrally stable}
  ($\Amat_\transverse = 0$). Feedback $\delta u = -k\, \xi$ gives
  $\dot{\xi} = -k\sin^2\!t\, \xi$ (after averaging), which is stable
  for $k > 0$. The orbit is stabilizable.
\end{example}

% ══════════════════════════════════════════════
\section{Floquet Theory: Stability of Periodic Orbits}
% ══════════════════════════════════════════════

For periodic trajectories (limit cycles, walking gaits, repetitive
manufacturing), the transverse linearization is a \emph{periodic} linear
system. Floquet theory provides the definitive stability analysis tool.

\subsection{The State Transition Matrix}

\begin{definition}{State Transition Matrix}{stm}
  For the linear time-varying system $\dot{\bm{\xi}} = \Amat_\transverse(t)\,
  \bm{\xi}$, the \textbf{state transition matrix}
  $\PhiMat(t, t_0)$ satisfies
  \begin{equation}
    \dot{\Phi}(t, t_0) = \Amat_\transverse(t)\, \PhiMat(t, t_0), \qquad
    \PhiMat(t_0, t_0) = \bm{I}.
  \end{equation}
  The solution is $\bm{\xi}(t) = \PhiMat(t, t_0)\, \bm{\xi}(t_0)$.
\end{definition}

\subsection{The Monodromy Matrix}

\begin{definition}{Monodromy Matrix}{monodromy}
  For a periodic orbit with period $T$, the \textbf{monodromy matrix}
  is the state transition matrix over one complete period:
  \begin{equation}
    \monodromy = \PhiMat(T, 0).
  \end{equation}
  The eigenvalues $\mu_1, \ldots, \mu_{n-1}$ of $\monodromy$ are the
  \textbf{Floquet multipliers} (also called \textbf{characteristic
  multipliers}).
\end{definition}

\begin{theorem}[Floquet Stability Criterion]\label{thm:floquet}
  A periodic orbit is:
  \begin{enumerate}[label=(\roman*)]
    \item \textbf{Orbitally stable} if all Floquet multipliers satisfy
      $|\mu_k| \leq 1$, with multipliers on the unit circle being
      semisimple.
    \item \textbf{Asymptotically orbitally stable} if all Floquet multipliers
      satisfy $|\mu_k| < 1$ (all strictly inside the unit circle).
    \item \textbf{Orbitally unstable} if any Floquet multiplier satisfies
      $|\mu_k| > 1$.
  \end{enumerate}
\end{theorem}

\begin{remark}{The Missing Multiplier}{missing_multiplier}
  If we had computed the monodromy matrix for the \emph{full}
  $n$-dimensional system (not the transverse reduction), there would be
  $n$ Floquet multipliers, with one of them exactly equal to $+1$. This
  ``trivial'' multiplier corresponds to the tangential direction---perturbations
  along the orbit neither grow nor decay. The transverse reduction removes
  this trivial multiplier, leaving only the $n - 1$ multipliers that
  matter for orbital stability.

  This is why the transverse reduction is essential: without it, the
  monodromy matrix \emph{always} has an eigenvalue at $+1$, making the
  orbit appear only marginally stable even when it is asymptotically
  orbitally stable.
\end{remark}

\begin{example}{Floquet Analysis of a Walking Gait}{walking_floquet}
  Consider a simplified compass-gait biped with 4-dimensional state space
  $(\theta_1, \theta_2, \dot{\theta}_1, \dot{\theta}_2)$. The walking gait
  is a periodic orbit with period $T$ (one stride). The transverse
  linearization has dimension $n - 1 = 3$.

  After numerically integrating $\PhiMat(T, 0)$ for the transverse dynamics,
  suppose we find the monodromy matrix has eigenvalues
  $\mu_1 = 0.7$, $\mu_2 = 0.4$, $\mu_3 = -0.3$.
  Since all $|\mu_k| < 1$, the gait is asymptotically orbitally stable.

  The physical interpretation: $\mu_1 = 0.7$ means that after one stride,
  transverse deviations in the first mode are reduced to 70\% of their
  initial value. After 10 strides, they are reduced to $0.7^{10} \approx 3\%$.
  The negative multiplier $\mu_3 = -0.3$ indicates an alternating (flip-flop)
  convergence pattern in the third mode.

  For the golf swing (a non-periodic motion), we cannot use the monodromy
  matrix directly, but the state transition matrix $\PhiMat(T, 0)$ plays an
  analogous role---its singular values measure how much transverse
  deviations amplify or attenuate over the course of the motion.
\end{example}

% ══════════════════════════════════════════════
\section{Moving Poincar\'e Sections}
% ══════════════════════════════════════════════

A complementary tool for analyzing orbital stability is the
\textbf{Poincar\'e section}---a codimension-1 surface that the orbit
crosses transversally. While Floquet theory gives continuous-time
information, Poincar\'e sections give \emph{discrete-time} (strobed)
information.

\begin{definition}{Poincar\'e Section}{poincare_section}
  A \textbf{Poincar\'e section} for an orbit $\orbit$ is a smooth
  $(n{-}1)$-dimensional surface $\poincare \subset \statespace$ such that:
  \begin{enumerate}[label=(\roman*)]
    \item $\orbit$ intersects $\poincare$ transversally (not tangentially).
    \item In a neighborhood of the intersection, every orbit near $\orbit$
      also crosses $\poincare$.
  \end{enumerate}
  The \textbf{Poincar\'e return map} $\bm{P}: \poincare \to \poincare$ sends
  each point on $\poincare$ to the next intersection with $\poincare$ under
  the flow.
\end{definition}

For periodic orbits, the orbit crosses $\poincare$ once per period, and
the Poincar\'e return map has a fixed point at the intersection
$\state^* = \orbit \cap \poincare$. Orbital stability of $\orbit$ is
equivalent to stability of this fixed point under the return map.

\begin{theorem}[Poincar\'e--Floquet Connection]\label{thm:poincare_floquet}
  The Jacobian of the Poincar\'e return map at the fixed point equals the
  transverse monodromy matrix:
  \begin{equation}
    D\bm{P}(\state^*) = \monodromy_\transverse.
  \end{equation}
  The Floquet multipliers and the eigenvalues of the Poincar\'e map are
  identical.
\end{theorem}

\subsection{Moving Poincar\'e Sections for Non-Periodic Trajectories}

The golf swing is not periodic---it is a \emph{finite-time} trajectory.
The classical Poincar\'e section (which relies on periodicity) does not
directly apply. However, we can generalize to \textbf{moving Poincar\'e
sections}: a continuous family of transverse hyperplanes, one at each point
of the orbit.

\begin{definition}{Moving Poincar\'e Section}{moving_poincare}
  A \textbf{moving Poincar\'e section} along a trajectory $\traj^*(s)$,
  $s \in [0, L]$, is the family of hyperplanes
  \begin{equation}
    \poincare(s) = \traj^*(s) + \Sigma(s),
  \end{equation}
  where $\Sigma(s)$ is the transverse hyperplane at $s$. The
  \textbf{section-to-section map} from $\poincare(s_1)$ to $\poincare(s_2)$
  is defined by following the flow from $s_1$ to $s_2$:
  \begin{equation}
    \bm{P}_{s_1 \to s_2}: \poincare(s_1) \to \poincare(s_2), \qquad
    \bm{\xi}(s_1) \mapsto \bm{\xi}(s_2).
  \end{equation}
  The linearized section-to-section map is the state transition matrix
  $\PhiMat_\transverse(s_2, s_1)$ of the transverse dynamics.
\end{definition}

For the golf swing, this gives us a powerful analysis tool:

\begin{principle}{Finite-Time Orbital Stability}{finite_orbital}
  For a finite-time trajectory $\traj^*: [0, T] \to \statespace$, define
  the \textbf{transverse amplification ratio} from phase $s_1$ to $s_2$:
  \begin{equation}\label{eq:amplification}
    \rho(s_1, s_2) = \norm{\PhiMat_\transverse(s_2, s_1)}_2
    = \sigma_{\max}\!\left(\PhiMat_\transverse(s_2, s_1)\right),
  \end{equation}
  where $\sigma_{\max}$ is the largest singular value. If
  $\rho(s_1, s_2) < 1$, then transverse deviations shrink between $s_1$
  and $s_2$. If $\rho(s_1, s_2) > 1$, they grow.

  For trajectory design, we require:
  \begin{equation}
    \rho(0, s_{\text{impact}}) \ll 1 \qquad \text{(deviations shrink
    toward impact)}.
  \end{equation}
  This is a precise formulation of the funnel concept from Chapter~1:
  the tube narrows because the transverse dynamics are contracting.
\end{principle}

\begin{example}{Amplification Profile of the Golf Swing}{golf_amplification}
  Consider the 8-dimensional state space of a 4-DOF golf swing model.
  The transverse dynamics have dimension 7. The transverse state transition
  matrix $\PhiMat_\transverse(s, 0)$ maps initial deviations to deviations at
  phase $s$.

  Numerical computation reveals the amplification ratio $\rho(0, s)$ as a
  function of $s$:
  \begin{itemize}[leftmargin=1cm]
    \item $s \in [0, 0.3]$ (backswing): $\rho$ is approximately 1---deviations
      neither grow nor shrink. The motion is kinematically slow and nearly
      linear.
    \item $s \in [0.3, 0.5]$ (transition): $\rho$ \emph{increases} past 1.
      This is the unstable phase---the high curvature and speed change
      amplify deviations. Without feedback, errors grow here.
    \item $s \in [0.5, 0.7]$ (late downswing): $\rho$ \emph{decreases}.
      The passive dynamics of the double-pendulum release are
      self-stabilizing in the transverse directions. The wrist release
      ``funnels'' the club toward the correct impact state.
    \item $s = 0.75$ (impact): $\rho(0, 0.75) < 1$. The net effect over
      the entire swing is transverse contraction---deviations are smaller
      at impact than at address.
  \end{itemize}

  This profile explains why skilled golfers are accurate despite the
  extreme nonlinearity of the swing: the trajectory is designed so that its
  \emph{passive transverse dynamics} are contracting near impact, creating
  a natural funnel that corrects errors without active feedback.
\end{example}

% ══════════════════════════════════════════════
\section{Controller Design via Transverse Stabilization}
% ══════════════════════════════════════════════

We now use the transverse linearization to design tracking controllers.
The approach is systematic: stabilize the transverse linear system
\eqref{eq:transverse_linear} using standard LTV techniques, then map
the result back to the nonlinear system.

\subsection{Time-Varying LQR for Transverse Stabilization}

\begin{principle}{Transverse LQR}{transverse_lqr}
  The trajectory tracking problem reduces to solving the \emph{differential
  Riccati equation} for the transverse dynamics:
  \begin{equation}\label{eq:riccati}
    -\dot{\bm{S}}(t) = \Amat_\transverse^\top \bm{S}
    + \bm{S}\, \Amat_\transverse
    - \bm{S}\, \Bmat_\transverse \bm{R}^{-1} \Bmat_\transverse^\top \bm{S}
    + \bm{Q}(t),
  \end{equation}
  with terminal condition $\bm{S}(T) = \bm{S}_T$. The optimal transverse
  feedback gain is
  \begin{equation}
    \bm{K}_\transverse(t) = \bm{R}^{-1}\, \Bmat_\transverse(t)^\top\, \bm{S}(t).
  \end{equation}
  The full control law is
  \begin{equation}\label{eq:full_control_law}
    \control(t) = \underbrace{\control^*(t)}_{\text{feedforward}}
    + \underbrace{\bm{N}(s^*)\, \bm{K}_\transverse(t)\,
    \bm{N}(s^*)^\top (\state - \traj^*(s^*))}_{\text{transverse feedback}},
  \end{equation}
  where $\bm{N}(s^*)$ maps between the transverse coordinates $\bm{\xi}$
  and the full state coordinates.
\end{principle}

The weighting matrices $\bm{Q}(t)$ and $\bm{R}$ have clear physical
interpretations:
\begin{itemize}[leftmargin=1.5cm]
  \item $\bm{Q}(t)$ penalizes transverse deviation. Making $\bm{Q}$ large
    near impact tightens the funnel where precision matters.
  \item $\bm{R}$ penalizes control effort. Making $\bm{R}$ small allows
    aggressive correction but may demand unrealistic actuator authority.
  \item The \emph{phase-dependent} nature of $\bm{Q}(t)$ is essential:
    we want tight tracking near impact and loose tracking during the backswing.
    This is the time-varying tube from Chapter~1, now realized through the
    Riccati equation.
\end{itemize}

\subsection{Phase-Varying Gains and the Funnel Connection}

The terminal-value Riccati equation \eqref{eq:riccati} naturally produces
gains that tighten as the terminal time approaches. This is exactly the
funnel behavior we designed geometrically in Chapter~1:

\begin{theorem}[Riccati Funnel]\label{thm:riccati_funnel}
  Let $\bm{S}(t)$ be the solution of \eqref{eq:riccati} with terminal
  penalty $\bm{S}_T \succ 0$. The ellipsoidal set
  \begin{equation}
    \mathcal{E}(t) = \left\{\bm{\xi} \in \Sigma(s) :
    \bm{\xi}^\top \bm{S}(t)\, \bm{\xi} \leq c \right\}
  \end{equation}
  defines a time-varying tube whose width (in the direction of eigenvector
  $\bm{v}_k$ of $\bm{S}$) is proportional to $1/\sqrt{\lambda_k(t)}$, where
  $\lambda_k(t)$ is the corresponding eigenvalue of $\bm{S}(t)$.

  As $t \to T$ and $\bm{S}(t) \to \bm{S}_T$, the tube tightens. The
  closed-loop transverse dynamics ensure that any trajectory starting inside
  the tube at $t = 0$ remains inside for all $t \in [0, T]$, provided the
  control authority is sufficient.
\end{theorem}

This is the formal connection between the geometric funnel of Chapter~1 and
the algebraic Riccati equation: \textbf{the funnel is the level set of the
Riccati cost-to-go.}

\subsection{The Complete Tracking Controller Architecture}

Assembling all the pieces, the trajectory tracking controller has three
layers:

\begin{center}
\begin{tikzpicture}[
  block/.style={draw, rounded corners, minimum height=1.2cm,
    minimum width=3.5cm, font=\small, align=center, thick},
  arrow/.style={-{Stealth[length=3mm]}, thick}
]
  % Blocks
  \node[block, fill=trajgreen!15, draw=trajgreen!70] (ff) at (0, 3)
    {\textbf{Feedforward}\\ $\control^*(t)$ from\\ inverse dynamics};
  \node[block, fill=chapblue!10, draw=chapblue!70] (proj) at (0, 1.2)
    {\textbf{Projection}\\ $s^* = \pi(\state)$\\ $\bm{\xi} = \bm{P}\,(\state - \traj^*(s^*))$};
  \node[block, fill=darkgold!15, draw=darkgold!70] (fb) at (0, -0.6)
    {\textbf{Transverse Feedback}\\ $\delta\control = \bm{K}_\transverse(s^*)\, \bm{\xi}$};
  \node[block, fill=accentred!10, draw=accentred!70] (plant) at (5.5, 1.2)
    {\textbf{Plant}\\ $\dot{\state} = \bm{f}(\state, \control)$};

  % Arrows
  \draw[arrow] (ff) -| (4, 3) |- node[above, font=\scriptsize, pos=0.25]
    {$\control^*$} (plant.170);
  \draw[arrow] (fb) -| (4, -0.6) |- node[below, font=\scriptsize, pos=0.25]
    {$\delta\control$} (plant.190);
  \draw[arrow] (plant) -- ++(2, 0) node[right, font=\small] {$\state(t)$}
    |- (7.5, -2) -| (-3, -2) |- (proj.west);
  \draw[arrow] (proj) -- (fb);

  % Sum
  \node[font=\large] at (4.3, 1.2) {$+$};
\end{tikzpicture}
\end{center}

\begin{enumerate}[label=\textbf{Layer \arabic*.}, leftmargin=2cm]
  \item \textbf{Feedforward} $\control^*(t)$: computed offline from inverse
    dynamics, provides $\sim$90\% of the total control effort.
  \item \textbf{Projection}: measures the current state $\state(t)$,
    projects onto the orbit to find $s^*$, and computes the transverse
    deviation $\bm{\xi}$.
  \item \textbf{Transverse feedback} $\delta\control = \bm{K}_\transverse(s^*)\, \bm{\xi}$:
    stabilizes the transverse dynamics using phase-varying gains from the
    Riccati equation.
\end{enumerate}

\begin{remark}{Why Not Just Use Standard LQR?}{why_not_lqr}
  One could linearize the \emph{full} dynamics about $\traj^*(t)$ and apply
  standard time-varying LQR to the $n$-dimensional error
  $\bm{e}(t) = \state(t) - \traj^*(t)$. This works, but:
  \begin{enumerate}[label=(\alph*)]
    \item It ``wastes'' control effort trying to correct tangential deviations
      (timing errors), which are benign.
    \item It fights the inherent marginal stability in the tangential direction,
      leading to unnecessarily high gains.
    \item It does not produce the correct funnel shape---it penalizes all
      errors equally, including timing errors.
  \end{enumerate}
  The transverse approach avoids all three problems. It is geometrically
  correct, computationally smaller (dimension $n - 1$ vs.\ $n$), and
  physically meaningful.
\end{remark}

% ══════════════════════════════════════════════
\section{Robustness of Orbital Stability}
% ══════════════════════════════════════════════

A practical controller must be robust to model uncertainty and external
disturbances. The transverse framework provides natural robustness margins.

\begin{definition}{Transverse Gain and Phase Margins}{transverse_margins}
  For the closed-loop transverse system
  $\dot{\bm{\xi}} = [\Amat_\transverse - \Bmat_\transverse \bm{K}_\transverse]\, \bm{\xi}$,
  the \textbf{transverse gain margin} and \textbf{transverse phase margin}
  are the gain and phase margins of the transverse loop transfer function.
  These are the trajectory-tracking analogues of the classical setpoint
  stability margins.
\end{definition}

\begin{proposition}[Robustness of Transverse LQR]\label{prop:lqr_robust}
  If the transverse LQR is designed with $\bm{R} = r\bm{I}$, then the
  closed-loop transverse system has guaranteed gain margin
  $[1/2, \infty)$ and phase margin $\pm 60^\circ$, inheriting the
  classical LQR robustness guarantees in each transverse channel.
\end{proposition}

For the golf swing, robustness translates directly to \emph{repeatability}:
a swing with high transverse robustness margins produces consistent results
even when the golfer's motor commands are noisy.

% ══════════════════════════════════════════════
\section{Summary}
% ══════════════════════════════════════════════

\begin{center}
\renewcommand{\arraystretch}{1.4}
\begin{tabular}{@{} l p{8.5cm} @{}}
  \toprule
  \textbf{Concept} & \textbf{Significance} \\
  \midrule
  Orbital stability & The correct stability notion for moving systems:
    only cross-orbit deviations matter \\
  Transverse linearization & Reduces orbital stability to $(n{-}1)$-dim
    LTV stability; removes the tangential ``nuisance'' direction \\
  Floquet multipliers & Eigenvalues of the transverse monodromy matrix;
    determine stability of periodic orbits \\
  Moving Poincar\'e sections & Extend Poincar\'e analysis to non-periodic
    trajectories; connect to the state transition matrix \\
  Amplification ratio $\rho$ & Maximum singular value of $\PhiMat_\transverse$;
    quantifies transverse error growth or contraction \\
  Transverse LQR & Phase-varying feedback from the Riccati equation; the
    optimal funnel controller \\
  Riccati funnel & The LQR cost-to-go defines the funnel boundary; connects
    geometric and algebraic views \\
  \bottomrule
\end{tabular}
\end{center}

\vspace{0.3cm}

This chapter has provided the mathematical backbone for everything that
follows. The transverse linearization converts the nonlinear trajectory-tracking
problem into a linear, lower-dimensional problem that inherits all the powerful
tools of LTV control theory. The Floquet/Poincar\'e framework provides both
analytical insight and numerical algorithms for assessing stability.

The key insight for the golf swing is that orbital stability---not Lyapunov
stability---is the correct framework. The swing is designed so that its
passive transverse dynamics create a natural funnel toward impact, and the
transverse LQR controller tightens this funnel using phase-varying gains
that match the demands of each phase of the motion.

In Chapter~5, we will see how underactuation constrains the transverse dynamics,
creating geometric structure that can be \emph{exploited} rather than merely
tolerated.

% ══════════════════════════════════════════════
\section{Exercises}
% ══════════════════════════════════════════════

\begin{enumerate}[label=\textbf{4.\arabic*.}, leftmargin=1.5cm]

\item \textbf{(Conceptual.)} Explain why a pendulum swinging on its limit
  cycle (pumped by a periodic torque) is orbitally stable but \emph{not}
  Lyapunov stable. Sketch the phase portrait and identify the tangential
  and transverse directions.

\item \textbf{(Computation.)} For the Van der Pol oscillator
  $\ddot{x} - \mu(1 - x^2)\dot{x} + x = 0$ with $\mu = 1$:
  \begin{enumerate}[label=(\alph*)]
    \item Numerically find the limit cycle by simulating from
      $(x, \dot{x}) = (3, 0)$ until transient effects decay.
    \item Compute the transverse linearization along the limit cycle.
    \item Compute the monodromy matrix and Floquet multipliers numerically.
    \item Verify that one multiplier of the full (2D) system is exactly 1.
  \end{enumerate}

\item \textbf{(Transverse linearization.)} For the system
  $\dot{x}_1 = x_2 + u$, $\dot{x}_2 = -x_1^3$, consider the trajectory
  $\traj^*(t) = (\cos t, -\sin t)$ with appropriate $u^*$.
  \begin{enumerate}[label=(\alph*)]
    \item Find $u^*(t)$ such that $\traj^*$ is a solution.
    \item Compute the full Jacobian $\bm{J}(t)$ along $\traj^*$.
    \item Compute the transverse Jacobian $\Amat_\transverse(t)$.
    \item Is the orbit stable open-loop? Design stabilizing transverse feedback.
  \end{enumerate}

\item \textbf{(Numerical project: Transverse LQR.)} For the double pendulum
  from Exercise~3.4:
  \begin{enumerate}[label=(\alph*)]
    \item Compute a reference swing trajectory using direct collocation
      (or load from a data file).
    \item Compute the transverse linearization along the trajectory.
    \item Solve the differential Riccati equation with $\bm{Q}(t)$ that
      increases 10-fold near the terminal phase (``impact'').
    \item Simulate the closed-loop system with random initial perturbations
      and plot the transverse deviation $\norm{\bm{\xi}(t)}$ over time.
    \item Compute the amplification ratio profile $\rho(0, s)$ and identify
      the ``natural funnel'' phase.
  \end{enumerate}

\item \textbf{(Design.)} A walking robot with period-1 gait has Floquet
  multipliers $\mu_1 = 0.9$, $\mu_2 = 1.1$, $\mu_3 = 0.5$.
  \begin{enumerate}[label=(\alph*)]
    \item Is the gait orbitally stable? Which mode is problematic?
    \item Design a transverse feedback controller that places the unstable
      multiplier $\mu_2$ at $0.6$. What gain is required?
    \item After stabilization, what is the worst-case convergence rate
      per stride?
    \item How many strides does it take to reduce a 10\% transverse
      deviation to 1\%?
  \end{enumerate}

\item \textbf{(Philosophical.)} Compare the transverse LQR approach
  (this chapter) with the standard trajectory-tracking LQR (linearize about
  $\traj^*(t)$, regulate $\bm{e} = \state - \traj^*(t)$ to zero).
  \begin{enumerate}[label=(\alph*)]
    \item Under what conditions do the two approaches give identical feedback
      gains? (Hint: consider straight-line trajectories.)
    \item Construct a numerical example where standard LQR uses significantly
      more control effort than transverse LQR.
    \item Discuss which approach is preferable for the golf swing and why.
  \end{enumerate}

\end{enumerate}

\vspace{1cm}
\begin{center}
  \rule{0.5\textwidth}{0.4pt}\\[0.5cm]
  {\small\itshape
  The pendulum does not care if you are watching at noon or midnight.\\
  Its orbit knows no clock.\\[0.2cm]
  Stability is not arrival. It is the promise\\
  that the motion, once begun, will keep its shape---\\
  even when the world pushes back.}
\end{center}

\chapter{Underactuation and Passive Dynamics}

\epigraph{%
  Do not force the motion. Let the physics do the work.
}{}

\vspace{0.5cm}

% ══════════════════════════════════════════════
\section{The Loss of Control Authority}
% ══════════════════════════════════════════════

In classical robotics, the goal is often full actuation: one motor for every degree of freedom. If a robot has six joints, it has six motors. This fully actuated paradigm allows the system to follow essentially any smooth trajectory in the configuration space, subject only to actuator saturation limits. The control problem is algebraically simple because the mapping from torques to joint accelerations is invertible.

But nature does not build fully actuated systems, and neither do athletes.

When a gymnast swings on a high bar, their center of mass rotates around the bar, but they have no motor at the bar---that ``joint'' is a passive pin constraint. When a human walks, the foot acting as the pivot is unactuated during the single-support phase; we cannot create torque from a point contact. And crucially for our running example, when a golfer releases the club during the downswing, the wrist joint acts as a free hinge. The club rapidly accelerates to over 100 miles per hour, \emph{driven entirely by the passive physics} of the arms pulling the grip, not by muscular torque applied at the wrist.

\begin{definition}{Underactuated System}{underactuated}
  Consider a mechanical system with $n$ degrees of freedom (DOF) and configuration coordinates $\bm{q} \in \configspace \subseteq \Reals^n$. The equations of motion take the standard manipulator form:
  \begin{equation}
    \Mmat(\bm{q})\ddot{\bm{q}} + \Cmat(\bm{q}, \dot{\bm{q}})\dot{\bm{q}} + \gvec(\bm{q}) = \Bmat(\bm{q}) \control
  \end{equation}
  where $\Mmat$ is the inertia matrix, $\Cmat$ contains Coriolis and centrifugal terms, $\gvec$ is the gravity vector, $\control \in \Reals^m$ is the vector of control inputs, and $\Bmat(\bm{q}) \in \Reals^{n \times m}$ maps inputs into generalized forces.
  
  The system is \textbf{underactuated} if $m < n$, meaning there are fewer actuators than degrees of freedom. In this case, the matrix $\Bmat$ is not full rank (it maps a lower-dimensional control space into the $n$-dimensional configuration space).
\end{definition}

If we lack the authority to command arbitrary accelerations, the classical setpoint perspective sees this as a profound failing. A setpoint controller wants to cancel the nonlinearities, feedback-linearize the system, and dictate the exact motion. Underactuation prevents this brute-force cancellation. 

However, from the trajectory-first perspective developed in Chapter~1, underactuation is not merely a restriction---it is a geometric feature. The \emph{passive dynamics} (the unactuated portions of the dynamics) define a manifold of feasible trajectories. We don't fight the physics; we ride it. 

% ══════════════════════════════════════════════
\section{The Double Pendulum: A Window into the Golf Swing}
% ══════════════════════════════════════════════

To understand how passive dynamics can be exploited, we study the simplest model that captures the physics of the golf downswing: the planar underactuated double pendulum.

Let the first link (the arms) have angle $\theta_1$, length $l_1$, mass $m_1$, and inertia $I_1$. Let the second link (the club) have angle $\theta_2$, length $l_2$, mass $m_2$, and inertia $I_2$. The configuration is $\bm{q} = (\theta_1, \theta_2)^\top$.

Crucially, we assume there is a motor at the shoulder (controlling torque $u_1 = \tau_{\text{shoulder}}$) but \emph{no motor} at the wrist ($\tau_{\text{wrist}} = 0$).

The input matrix is therefore:
\begin{equation}
  \Bmat = \begin{pmatrix} 1 \\ 0 \end{pmatrix}, \qquad \control = [u_1].
\end{equation}

The dynamics expand into two coupled equations:
\begin{align}
  M_{11}\ddot{\theta}_1 + M_{12}\ddot{\theta}_2 + C_1 + g_1 &= u_1, \label{eq:pend_actuated} \\
  M_{21}\ddot{\theta}_1 + M_{22}\ddot{\theta}_2 + C_2 + g_2 &= 0. \label{eq:pend_unactuated}
\end{align}

Equation \eqref{eq:pend_unactuated} is the \textbf{unactuated dynamics} or \textbf{passive constraint}. It dictates the motion of the club. Notice that we cannot directly command $\ddot{\theta}_2$. However, we \emph{can} influence it through the coupling terms:
\begin{enumerate}[label=(\roman*)]
  \item \textbf{Inertial coupling} ($M_{21}\ddot{\theta}_1$): Accelerating or decelerating the arms immediately exerts a force on the club.
  \item \textbf{Velocity coupling} ($C_2$): This term contains the centrifugal forces. As the arms swing rapidly, centrifugal force naturally pulls the club outward.
\end{enumerate}

\begin{keyidea}{The Physics of the Wrist Release}{wrist_release}
  In a golf swing, the huge clubhead speed at impact is generated by exploiting the $C_2$ and $M_{21}$ terms. Early in the downswing, the golfer maintains an acute angle between the arms and the club ($\theta_2 - \theta_1 \approx 90^\circ$). The arms accelerate ($\ddot{\theta}_1 > 0$). 
  
  As the hands approach the bottom of the swing arc, the arms decelerate ($\ddot{\theta}_1 < 0$). This negative acceleration, transferred through the $M_{21}$ coupling, causes the club to rapidly whip forward. Simultaneously, the massive velocity of the arms ($\dot{\theta}_1$) creates a large centrifugal force ($C_2$) that drives the club in-line with the arms ($\theta_2 \to \theta_1$). The golfer does not ``unhinge'' their wrists with muscular effort; the physics forces the release.
\end{keyidea}

Let's look at the unactuated row algebraically to see this whipping action:
\begin{equation}
  \ddot{\theta}_2 = -M_{22}^{-1} \bigl( M_{21}\ddot{\theta}_1 + C_2(\bm{q}, \dot{\bm{q}}) + g_2(\bm{q}) \bigr).
\end{equation}
Since $M_{22}$ is scalar and positive, a large negative $\ddot{\theta}_1$ (braking the arms) produces a large positive $\ddot{\theta}_2$ (accelerating the club). This is the parametric excitation that drives the moving target control problem. It is an optimal transfer of momentum.

\begin{center}
\begin{tikzpicture}[scale=1.2]
  % Arm (Link 1)
  \draw[thick, chapblue, line width=1.5mm] (0, 0) -- (2, -2);
  \filldraw[chapblue!50!black] (0, 0) circle (4pt) node[above left, black, font=\small] {Shoulder ($u_1$)};
  
  % Club (Link 2)
  \draw[thick, accentred, line width=1mm] (2, -2) -- (1, -4);
  \filldraw[accentred!50!black] (2, -2) circle (3pt) node[right, black, font=\small, xshift=2pt] {Wrist (passive)};
  \filldraw[gray!80!black] (1, -4) circle (4pt) node[below, black, font=\small] {Clubhead};

  % Angles
  \draw[dashed, gray] (0, 0) -- (0, -2);
  \draw[-{Stealth[length=2mm]}, gray] (0, -1) arc (270:315:1);
  \node[font=\small] at (0.3, -1.2) {$\theta_1$};

  \draw[dashed, gray] (2, -2) -- (2, -4);
  \draw[-{Stealth[length=2mm]}, gray] (2, -3) arc (270:245:1);
  \node[font=\small] at (1.8, -3.2) {$\theta_2$};
  
  % Centrifugal
  \draw[-{Stealth[length=3mm]}, thick, trajgreen] (2, -2) -- (3, -3) node[right, font=\scriptsize] {Centrifugal pull};
  
  % Arm deceleration
  \draw[-{Stealth[length=3mm]}, thick, chapblue!80] (1, -1) arc (135:100:1.4) node[above, font=\scriptsize, align=center] {Arm brakes\\($\ddot{\theta}_1 < 0$)};
  
  % Club acceleration
  \draw[-{Stealth[length=3mm]}, thick, accentred!80] (1.5, -3) arc (245:285:1) node[below right, font=\scriptsize, align=center] {Club whips\\($\ddot{\theta}_2 \gg 0$)};
\end{tikzpicture}
\end{center}

% ══════════════════════════════════════════════
\section{The Annihilator and the Null Space}
% ══════════════════════════════════════════════

The double pendulum example shows that the passive dynamics constrain the allowable states and accelerations. We generalize this using the mathematics of the \textbf{null space}.

Recall the dynamics:
\begin{equation}
  \Mmat(\bm{q})\ddot{\bm{q}} + \Cmat(\bm{q}, \dot{\bm{q}})\dot{\bm{q}} + \gvec(\bm{q}) = \Bmat(\bm{q}) \control.
\end{equation}

\begin{definition}{Left Annihilator}{annihilator}
  We define a full-rank matrix $\Bmat^\perp(\bm{q}) \in \Reals^{(n-m) \times n}$, called the \textbf{left annihilator} of $\Bmat$, such that:
  \begin{equation}
    \Bmat^\perp(\bm{q}) \Bmat(\bm{q}) = \bm{0} \quad \text{for all } \bm{q}.
  \end{equation}
  The rows of $\Bmat^\perp$ span the left null space of $\Bmat$.
\end{definition}

In our double pendulum where $\Bmat = [1, 0]^\top$, the annihilator is trivially $\Bmat^\perp = [0, 1]$.

By multiplying the equations of motion by the annihilator $\Bmat^\perp$, we eliminate the control input $\control$, revealing the intrinsic constraints on the system's acceleration:

\begin{equation}\label{eq:unactuated_dynamics}
  \Bmat^\perp \Mmat \ddot{\bm{q}} + \Bmat^\perp (\Cmat\dot{\bm{q}} + \gvec) = \bm{0}.
\end{equation}

\begin{principle}{The Theorem of Feasible Trajectories}{feasible_traj}
  A smooth curve $\traj^*(t)$ is a dynamically feasible trajectory for the underactuated system if and only if it satisfies the unactuated dynamics:
  \begin{equation}
    \Bmat^\perp(\traj^*) \Bigl[ \Mmat(\traj^*)\ddot{\traj}^* + \Cmat(\traj^*, \dot{\traj}^*)\dot{\traj}^* + \gvec(\traj^*) \Bigr] = \bm{0}
  \end{equation}
  for all $t \in [0, T]$. If a curve satisfies this condition, one can always find the necessary control input $\control^*(t)$ by projecting the dynamics onto the range of $\Bmat$.
\end{principle}

For control design, this means we cannot arbitrarily select a path in $\statespace$ and expect a controller to track it. The reference trajectory \emph{must} live in the restricted subspace defined by \eqref{eq:unactuated_dynamics}. The trajectory optimization algorithms in Chapter~6 will explicitly enforce this constraint.

% ══════════════════════════════════════════════
\section{Partial Feedback Linearization}
% ══════════════════════════════════════════════

While we cannot feedback-linearize the entire underactuated system, we can feedback-linearize the \emph{actuated} degrees of freedom. This technique is called \textbf{partial feedback linearization} or \textbf{collocated feedback linearization}.

By cleverly substituting coordinates, we can partition the system into a linear, fully actuated subsystem (say, the golfer's arms) and a nonlinear, passive subsystem (the club) that is driven by the state of the first.

Partition the coordinates as $\bm{q} = (\bm{q}_a, \bm{q}_u)$, representing actuated and unactuated DOFs. We can rewrite the dynamics and choose $\control$ to exactly cancel the nonlinearities of the actuated part:
\begin{equation}
  \ddot{\bm{q}}_a = \bm{v}_a
\end{equation}
where $\bm{v}_a$ is our new synthetic control input. The unactuated dynamics become:
\begin{equation}\label{eq:zero_dynamics}
  \Mmat_{uu}(\bm{q})\ddot{\bm{q}}_u + \Mmat_{ua}(\bm{q})\bm{v}_a + \text{nonlinear terms} = 0.
\end{equation}

If we force the actuated coordinates to track a reference trajectory (e.g., using a high-gain PD controller, $\bm{v}_a = \ddot{\bm{q}}_a^* - K_d \dot{\tilde{\bm{q}}}_a - K_p \tilde{\bm{q}}_a$), the unactuated coordinates are left to evolve according to \eqref{eq:zero_dynamics}. The stability of the resulting motion depends entirely on whether \eqref{eq:zero_dynamics} is a stable differential equation when forced by $\bm{q}_a^*$. This leads us to the concept of the \emph{zero dynamics}.

% ══════════════════════════════════════════════
\section{Zero Dynamics and the Golf Swing Funnel}
% ══════════════════════════════════════════════

If we perfectly stabilize the actuated tracking errors to zero ($\bm{q}_a \equiv \bm{q}_a^*$), the remaining dynamics governing $\bm{q}_u$ are called the \textbf{zero dynamics}.

For the golf swing, the zero dynamics describe how the club moves if the golfer's arms perfectly follow the nominal swing path. Are the zero dynamics of the golf swing stable?

Revisiting the concept of Transverse Linearization (Chapter~4), if we linearize the unactuated dynamics around the trajectory $\traj^*(t)$, we analyze the transverse divergence of the club from its optimal path.

We find a beautiful physical phenomenon:
\begin{itemize}
  \item During the early downswing, the club's zero dynamics are \emph{unstable}. If the club gets slightly out of position, the centrifugal force tends to pull it further off the optimal plane.
  \item During the late downswing (the release phase), the club's zero dynamics become \emph{strongly stable}. The massive acceleration $\ddot{\theta}_2 \gg 0$ acts like a stiff restoring spring. The centrifugal force aligns the club perfectly with the arms. 
\end{itemize}

This is the manifestation of the time-varying tube (the Funnel) from Chapter~1! The passive physics naturally form a wide, loose tube that tightens into an incredibly precise singularity exactly at the moment of impact. The golfer does not need to use wrist torque to steer the clubhead to the ball; if the shoulder inputs are grossly correct, the physics of the double pendulum \emph{pulls} the clubhead into the ball.

\begin{remark}{Why Good Swings Feel Effortless}{effortless}
  Novice golfers attempt to fully actuate the swing. They use their wrists to steer the club, treating it as a fully actuated robotic arm. This fights the passive dynamics, resulting in slower club speeds and higher variability. Pro golfers use their arms and torso to set up the boundary conditions for the passive wrist release. By acting primarily in the null space (letting the unactuated dynamics dominate), the swing maximizes speed and consistency. The physics does the work.
\end{remark}

% ══════════════════════════════════════════════
\section{Summary}
% ══════════════════════════════════════════════

\begin{center}
\renewcommand{\arraystretch}{1.4}
\begin{tabular}{@{} l p{8.5cm} @{}}
  \toprule
  \textbf{Concept} & \textbf{Significance} \\
  \midrule
  Underactuation & Fewer actuators ($m$) than coordinates ($n$). Cannot arbitrarily command accelerations in configuration space. \\
  Passive Dynamics & The natural, unforced motion of the unactuated DOFs. \\
  Coupling & Inertial ($M_{21}$) and velocity ($C_2$) terms transmit energy from actuated to unactuated DOFs. \\
  Left Annihilator & A matrix $\Bmat^\perp$ that strips away the control input to reveal the intrinsic geometric constraints of the motion. \\
  Partial Feedback Lin. & Cancelling nonlinearities only on the actuated joints to simplify control. \\
  Zero Dynamics & The residual dynamics of the unactuated subsystem when the actuated subsystem perfectly tracks its reference. \\
  \bottomrule
\end{tabular}
\end{center}

\vspace{0.3cm}

Underactuation forces a shift from ``forcing'' a system to an equilibrium, to ``orchestrating'' the natural physics of the system. The trajectory is a choreography between the motors and gravity/inertia. In Chapter~6, we will see how numerical trajectory optimization can synthesize these choreographies automatically by solving optimal control problems that explicitly respect the $\Bmat^\perp$ constraint.

% ══════════════════════════════════════════════
\section{Exercises}
% ══════════════════════════════════════════════

\begin{enumerate}[label=\textbf{5.\arabic*.}, leftmargin=1.5cm]

\item \textbf{(Annihilator construction.)} A cart-pole system has state $\bm{q} = (x, \theta)^\top$, where $x$ is the cart position and $\theta$ is the pole angle. The motor applies horizontal force to the cart.
  \begin{enumerate}
    \item What is the input matrix $\Bmat$?
    \item Construct a left annihilator $\Bmat^\perp$.
    \item Derive the unactuated dynamics constraint. What physical conservation law does this represent when gravity is ignored?
  \end{enumerate}

\item \textbf{(Passive pumping.)} In a playground swing, the person swinging cannot exert an external torque on the swing setup (there is no motor at the top pivot). By what mechanism (which dynamics term) does the swinger add energy to the system by leaning back and forth? Relate this to the golf wrist release.

\item \textbf{(Zero Dynamics Stability.)} Consider a generic system partitioned into $\ddot{q}_a = v_a$ and $\ddot{q}_u + c q_u = q_a$. Let the reference trajectory be $q_a^*(t) = \sin(t)$.
  \begin{enumerate}
    \item Assume we enforce $q_a = q_a^*$ perfectly. Write the zero dynamics for $q_u$.
    \item Under what conditions on the parameter $c$ are the zero dynamics stable? Is it possible to have a trajectory that passes through both stable and unstable regimes of zero dynamics?
  \end{enumerate}

\end{enumerate}

\vspace{1cm}
\begin{center}
  \rule{0.5\textwidth}{0.4pt}\\[0.5cm]
  {\small\itshape
  Control is not dominion over nature;\\
  it is an intelligent negotiation with it.}
\end{center}

\chapter{Trajectory Optimization}

\epigraph{%
  We do not merely seek a path that obeys the laws of physics.\\
  We seek the path that bends them to our advantage.
}{}

\vspace{0.5cm}

% ══════════════════════════════════════════════
\section{The Moving Target Problem}
% ══════════════════════════════════════════════

In Chapter 1, we fundamentally reframed our control objective from "reaching a destination" to "tracking a trajectory." Up to this point, we have assumed that this optimal reference trajectory $\traj^*(t)$ is given to us by an oracle. This chapter answers the obvious next question: \emph{Where does this reference come from?}

For systems whose purpose is to move---like our running example of the golf swing---the optimal trajectory is not arbitrary. We are not just trying to go from point A to point B while minimizing energy. We are trying to fulfill what we term the \textbf{Moving Target Objective}: optimizing the kinematics (velocity, pose, acceleration) of a system precisely at the moment it intersects a target manifold (e.g., the golf ball). 

Crucially, \emph{the exact clock time that impact occurs does not matter}. A golfer has the freedom to start their backswing slowly and pause at the top. The optimization is entirely concerned with the \emph{shape} of the motion and its terminal kinematics. Time is a resource, not a constraint.

\begin{definition}{The Moving Target Optimal Control Problem (OCP)}{moving_target_ocp}
  Find the state trajectory $\state(t)$, control input $\control(t)$, and terminal time $t_f$ that minimize the cost functional:
  \begin{equation}\label{eq:ocp_cost}
    J = \phi(\state(t_f)) + \int_0^{t_f} L(\state(t), \control(t)) \dt
  \end{equation}
  subject to the dynamics and constraints:
  \begin{align}
    \dot{\state}(t) &= \bm{f}(\state(t), \control(t)) \qquad &\text{(Dynamics)}\\
    \bm{h}(\state(t), \control(t)) &\leq \bm{0} \qquad &\text{(Path constraints)}\\
    \psi(\state(t_f)) &= \bm{0} \qquad &\text{(Terminal manifold/Moving goal)}\\
    \state(0) &= \state_0 \qquad &\text{(Initial state)}
  \end{align}
\end{definition}

In the golf downswing, the terminal cost $\phi(\state(t_f))$ might be negative clubhead speed (to maximize speed), the integral cost $L$ might penalize excessive joint torques to ensure repeatability, and the terminal manifold constraint $\psi(\state(t_f)) = \bm{0}$ forces the clubhead to be geometrically located at the ball with a squared face. Notice that $t_f$ is a \emph{free parameter} in the optimization.

% ══════════════════════════════════════════════
\section{Transcribing the Continuous Problem}
% ══════════════════════════════════════════════

The OCP described above occurs in infinite-dimensional function space. To solve it on a computer, we must transcribe it into a finite-dimensional Nonlinear Programming problem (NLP):
\begin{align*}
  \min_{\bm{w}} \quad & f(\bm{w}) \\
  \text{subject to} \quad & \bm{c}(\bm{w}) \leq \bm{0},
\end{align*}
where $\bm{w}$ is a large vector of decision variables. How we discretize the integrals and differential equations defines the transcription method. We will focus on two modern approaches: \textbf{Direct Shooting} and \textbf{Direct Collocation}.

\subsection{Direct Multi-Shooting}

In direct multi-shooting, we break the time interval $[0, t_f]$ into $N$ segments. The decision variables are the control parameters $\control_k$ on each interval and the state values $\state_k$ at the beginning of each interval.

The dynamic constraint $\dot{\state} = \bm{f}(\state, \control)$ is enforced by numerical integration (e.g., Runge-Kutta). We require that the integrated state at the end of interval $k$ exactly matches the assumed starting state of interval $k+1$:
\begin{equation}
  \state_{k+1} - \bm{\Phi}(\state_k, \control_k, \Delta t) = \bm{0} \qquad \text{(Defect constraint)}
\end{equation}
where $\bm{\Phi}$ represents the numerical integrator step. Multi-shooting is excellent for highly nonlinear dynamics because it keeps the integration intervals short, preventing the massive nonlinear sensitivities that plague single-shooting methods.

\subsection{Direct Collocation}

Direct Collocation takes this a step further. We represent both the state trajectory and the control signal as piecewise polynomials (e.g., cubic splines for state, linear splines for control). The decision variables are the coefficients of these polynomials, typically located at discrete \textbf{collocation points}.

Instead of using a numerical integrator, we enforce the differential equation algebraicaly at the collocation points:
\begin{equation}
  \dot{\tilde{\state}}(t_c) - \bm{f}(\tilde{\state}(t_c), \tilde{\control}(t_c)) = \bm{0} \qquad \text{(Collocation constraint)}
\end{equation}
where $\tilde{\state}$ and $\tilde{\control}$ are the polynomial approximations. 

\begin{keyidea}{Fidelity vs. Flexibility}{collocation_vs_shooting}
  Direct collocation results in a massive, highly sparse NLP. Because the solver has access to the state at every grid point simultaneously, it can easily manipulate the physical trajectory to satisfy difficult path constraints without waiting to "simulate" the result. Collocation is widely considered the state-of-the-art for generating complex dynamic motions like walking, jumping, and swinging.
\end{keyidea}

% ══════════════════════════════════════════════
\section{Optimizing Underactuated Geometries}
% ══════════════════════════════════════════════

Recall from Chapter 5 that the underactuated dynamics enforce a strict algebraic constraint on the system:
\begin{equation}
  \Bmat^\perp(\state) \Bigl[ \Mmat(\state)\ddot{\state} + \Cmat(\state, \dot{\state})\dot{\state} + \gvec(\state) \Bigr] = \bm{0}.
\end{equation}

When optimizing for a golf swing, direct collocation handles this underactuation organically. We do not need to explicitly partition the variables into actuated and unactuated sets for the solver. The collocation constraint $\dot{\state} - \bm{f} = \bm{0}$ already implicitly contains the passive dynamic constraints. 

\begin{example}{The Golf Swing OCP}{golf_ocp}
  Consider our planar double pendulum model from Chapter 5. The solver is given the task to maximize $\dot{\theta}_2(t_f)$ (club speed). 
  
  \textbf{Initial Guess:} If we seed the NLP solver with a naive guess (e.g., constant shoulder torque), the club speed will be pitifully low.
  
  \textbf{The Optimized Result:} The solver discovers the optimal physical orchestration entirely on its own. It will apply a massive positive $u_1$ to accelerate the arms, and then immediately reverse $u_1$ to negative maximum just prior to $t_f$. The optimizer mathematicaly "discovers" the necessary inertial coupling (the parametric excitation of $M_{21}\ddot{\theta}_1$) to whip the unactuated club link through the terminal manifold constraint! The passive dynamics are beautifully exploited to satisfy the objective.
\end{example}

% ══════════════════════════════════════════════
\section{Repeatability: The Stochastic OCP}
% ══════════════════════════════════════════════

The previous sections frame a deterministic problem: "Find the swing that maximizes terminal speed."

But as we argued in Chapter 1, biological systems and athletes do not optimize for peak deterministic performance. If you instruct an NLP solver to maximize golf club speed without regard for motor noise, the solver will ride the very edge of the constraint boundaries. It will produce a swing so delicately timed that a 5-millisecond execution error would result in a complete whiff.

\begin{principle}{The Repeatability Addendum}{repeatability_OCP}
  The theoretically optimal trajectory is useless if its local topology amplifies small execution errors into catastrophic terminal misses. The true objective must account for \emph{variance}.
\end{principle}

We modify our deterministic cost $J = \phi(\state(t_f)) + \int L$ into a stochastic expectation:
\begin{equation}
  \min_{\traj^*} \quad \mathbb{E}\left[ \phi(\state(t_f, \bm{w})) \right] + \lambda \cdot \mathrm{Var}\left[ \psi(\state(t_f, \bm{w})) \right]
\end{equation}
where $\bm{w} \sim \mathcal{N}(0, \Sigma_w)$ represents signal-dependent motor noise (a known property of human muscle actuation, where noise scales proportionally to the applied torque).

To solve this in a collocation framework without massive Monte Carlo simulations, we use \textbf{Covariance Steering} or \textbf{Linear Covariance Analysis}. We augment the state vector $\state(t)$ with its local covariance matrix $\bm{\Sigma}(t)$. The dynamics of $\bm{\Sigma}(t)$ are propagated alongside the nominal trajectory using the linearized Jacobian matrices along the path. 

The optimizer is now forced to "widen the funnel" as it approaches the moving goal. It sacrifices a small amount of peak clubhead speed to select a trajectory geometry where the passive dynamics are heavily stabilizing in the transverse direction. The result is a trajectory that a human can actually \emph{repeat}.

% ══════════════════════════════════════════════
\section{Summary}
% ══════════════════════════════════════════════

\begin{center}
\renewcommand{\arraystretch}{1.4}
\begin{tabular}{@{} l p{8.5cm} @{}}
  \toprule
  \textbf{Concept} & \textbf{Significance} \\
  \midrule
  Moving Target OCP & Finding a trajectory that optimizes terminal kinematics when intersecting a spatial manifold, without prescribing the clock time. \\
  Multi-Shooting & Discretizing the time domain into segments, using numerical integrators to enforce dynamics constraints between nodes. \\
  Direct Collocation & Representing the entire trajectory as splines and enforcing dynamics algebraically at discrete points arrayed along the trajectory. \\
  Passive Exploitation & Solvers naturally optimize underactuated topologies by weaponizing coupling matrices ($M_{21}, C_2$). \\
  Stochastic Robustness & Augmenting the OCP with expected variance to ensure the resulting trajectory can be consistently executed by noisy nonlinear plants. \\
  \bottomrule
\end{tabular}
\end{center}

\vspace{0.3cm}

Determining the shape of the reference trajectory is the cornerstone of Control Is Motion. With our passive geometry understood (Chap 5) and our optimal reference trajectory computed (Chap 6), we must now surround that reference trajectory with a guaranteed, robust control tube to reject disturbances online. We turn to Funnel Synthesis in Chapter 7.

% ══════════════════════════════════════════════
\section{Exercises}
% ══════════════════════════════════════════════

\begin{enumerate}[label=\textbf{6.\arabic*.}, leftmargin=1.5cm]

\item \textbf{(Free Terminal Time Formulation.)} Explain how one sets up an OCP where final time $t_f$ is a decision variable in a fixed-grid Direct Collocation solver. (Hint: consider scaling time $t = \tau \cdot t_f$, where $\tau \in [0, 1]$). Show how the dynamic constraint $\dot{\state} = \bm{f}$ changes under this substitution.

\item \textbf{(Kinematic Goals vs Setpoints.)} You are optimizing a baseball bat swing. Your state is $(\theta, \dot{\theta})$. Assume the ball arrives at the plate between $t=0.5$ and $t=0.6$ seconds. Formulate the path constraint and the terminal manifold constraint $\psi$ assuming you want maximum $\dot{\theta}$ when $\theta = \pi/2$, but the exact time of impact within that window does not matter.

\item \textbf{(Signal-Dependent Noise Penalty.)} Human muscle torque output $\tau$ has variance $\sigma^2 = k |\tau|^2$. Write a modified integral cost function $L(\state, \tau)$ that penalizes the introduction of variance into the system. Interpret what this does to the optimal trajectory shape.

\end{enumerate}

\vspace{1cm}
\begin{center}
  \rule{0.5\textwidth}{0.4pt}\\[0.5cm]
  {\small\itshape
  The perfect motion is rarely the fastest.\\
  It is the one that survives reality.}
\end{center}

\chapter{Funnel Synthesis}

\epigraph{%
  You cannot control the exact trajectory a stone will take down a mountain.\\
  But you can build a trough to ensure it reaches the bottom.
}{}

\vspace{0.5cm}

% ══════════════════════════════════════════════
\section{Beyond LQR: The Need for Guaranteed Invariance}
% ══════════════════════════════════════════════

In Chapter 4, we showed that applying time-varying Linear Quadratic Regulation (LQR) to the transverse dynamics creates a contracting envelope of stability. The Riccati equation naturally shapes this envelope into a "funnel" that narrows as the terminal constraint approaches. 

However, LQR guarantees are inherently \emph{local}. They rely on the linear approximation of the transverse dynamics $\dot{\bm{\xi}} = \bm{A}_\perp(t)\bm{\xi} + \bm{B}_\perp(t)\delta\control$. For actual physical systems traversing highly nonlinear state spaces—like a golf club accelerating to $50$ m/s while subject to massive centrifugal and Coriolis forces—this linearization quickly breaks down. If a disturbance pushes the state too far from the reference trajectory $\traj^*(t)$, the nonlinearities will dominate, the LQR restoring forces will compute the wrong corrective action, and the clubhead will completely miss the ball.

We need a mathematically rigorous way to define the \emph{exact boundary} of our funnel. We must prove that for the fully nonlinear equations of motion, any state starting within the mouth of the funnel at time $t=0$ will remain inside the funnel for all $t \in [0, t_f]$.

\begin{definition}{Dynamic Funnel}{dynamic_funnel}
  Given a reference trajectory $\traj^*(t)$, a feedback controller $\control = \bm{k}(\state, t)$, and a scalar boundary function $V(\state, t)$, the set:
  \begin{equation}
    \mathcal{F}(t) = \{ \state \in \statespace \mid V(\state, t) \leq 1 \}
  \end{equation}
  constitutes a valid \textbf{dynamic funnel} if the set is positively invariant under the closed-loop dynamics. That is, if $\state(0) \in \mathcal{F}(0)$, then $\state(t) \in \mathcal{F}(t)$ for all $t > 0$.
\end{definition}

% ══════════════════════════════════════════════
\section{Time-Varying Lyapunov Functions for Trajectories}
% ══════════════════════════════════════════════

To verify invariance, we use the machinery of Lyapunov theory, adapted for trajectories rather than equilibria. We require our boundary function $V(\state, t)$ to act as a time-varying Lyapunov function.

\begin{theorem}[Funnel Invariance]\label{thm:funnel_invariance}
  The set $\mathcal{F}(t) = \{ \state \mid V(\state, t) \leq 1 \}$ is invariant if, on the boundary where $V(\state, t) = 1$, the function is strictly decreasing along the system trajectories:
  \begin{equation}
    \dot{V}(\state, t) = \frac{\partial V}{\partial t} + \nabla_{\state} V(\state, t) \cdot \bm{f}(\state, \bm{k}(\state, t)) < 0 \quad \text{for all } \state \text{ s.t. } V(\state, t) = 1.
  \end{equation}
\end{theorem}

If the derivative of the boundary function points inward everywhere on the boundary, trajectories may circulate inside the funnel, but they can never escape it. For a moving target problem, we want the spatial volume of $\mathcal{F}(t)$ to shrink as $t \to t_f$, terminating in a tight bounded region at impact.

\begin{keyidea}{Funnels Define the Usable State Space}{usable_space}
  A computed funnel is not just an analysis tool; it is the ultimate judge of repeatability. It explicitly defines the exact set of initial conditions (e.g., the kinematic state at the top of the backswing) that are mathematically guaranteed to successfully strike the ball despite the presence of extreme nonlinearities.
\end{keyidea}

% ══════════════════════════════════════════════
\section{Sums-of-Squares (SOS) Programming}
% ══════════════════════════════════════════════

Proving that $\dot{V} < 0$ on the boundary $V=1$ for a general nonlinear function $\bm{f}$ is famously difficult. We have an infinite number of states on the boundary to check! We circumvent this by restricting our dynamics $\bm{f}$ algorithms to polynomials (via Taylor expansion or substitution) and leveraging \textbf{Sums-of-Squares (SOS) Programming}~\cite{Boyd1994, MajumdarTedrake2017}.

A polynomial $p(x)$ is a Sum of Squares if it can be written as $p(x) = \sum_{i} q_i(x)^2$ for some polynomials $q_i(x)$. If a polynomial is SOS, it is globally non-negative. Testing if a polynomial is SOS turns out to be equivalent to solving a Semidefinite Program (SDP)—a convex optimization problem that can be solved very efficiently.

We can reframe Theorem \ref{thm:funnel_invariance} using the S-procedure. To guarantee that $\dot{V}(\state, t) < 0$ whenever $V(\state, t) = 1$, we require:
\begin{equation}
  -\dot{V}(\state, t) - L(\state, t)(1 - V(\state, t)) \quad \text{is SOS}
\end{equation}
where $L(\state, t)$ is a strictly positive polynomial multiplier. If $V(\state, t) = 1$, the second term vanishes, and the SOS condition forces $-\dot{V} > 0$, achieving our invariance proof!

% ══════════════════════════════════════════════
\section{Computing the Funnel for the Golf Swing}
% ══════════════════════════════════════════════

Let us return to the double pendulum of the golf swing. Our reference trajectory $\traj^*(t)$ from Chapter 6 gives us the nominal kinematics. Our transverse LQR from Chapter 4 gives us a feedback matrix $\bm{K}(t)$ designed to stabilize deviations. We wish to compute the largest possible funnel around this sequence that the closed-loop system is mathematically guaranteed to stay inside.

We define our candidate Lyapunov function as a time-varying quadratic form:
\begin{equation}
  V(\state, t) = (\state - \traj^*(t))^\top \bm{P}(t) (\state - \traj^*(t))
\end{equation}
where $\bm{P}(t)$ is a positive definite matrix that changes continuously along the trajectory. The funnel boundary is the time-varying ellipsoid where $V=1$. By altering $\bm{P}(t)$, we alter the shape and orientation of the funnel.

\begin{example}{Volume Maximization via SDP}{vol_max}
  We transcribe the continuous time $t$ into discrete nodes $t_k$, just as we did in direct collocation. At each node, we set up an SDP that seeks to \emph{maximize} the volume of the ellipsoid defined by $\bm{P}(t_k)$, subject to the SOS invariance constraints connecting it to adjacent nodes.

  The optimizer actively computes the exact structural limits of the golf swing's passive dynamics. It discovers that early in the downswing, the funnel can be incredibly "wide" (forgiving), but as centrifugal forces build, the SOS constraints force $\bm{P}(t_k)$ to become highly elongated along the unactuated dimensions in order to maintain the invariance guarantee. The resulting geometry is a true numerical portrait of athletic consistency.
\end{example}

In Chapter 8, we will explore how to make these funnels robust to temporal variation by migrating out of the time domain entirely and into the phase domain.

\chapter{Phase-Variable Control}

\epigraph{%
  Time is a river that carries us forward. Phase is the boat we build to navigate it.
}{}

\vspace{0.5cm}

% ══════════════════════════════════════════════
\section{Spatial Paths vs. Temporal Execution}
% ══════════════════════════════════════════════

The most glaring flaw in classical control architecture, when applied to biomechanics or advanced robotics, is its rigid adherence to the clock. If you construct a time-varying trajectory $\traj^*(t)$ for a robotic arm and command a tracking controller to follow it, the controller will evaluate the error as $\bm{e}(t) = \state(t) - \traj^*(t)$. 

If the robot hits an unmodeled friction patch and slows down, the reference trajectory $\traj^*(t)$ continues advancing relentlessly. The controller calculates a massive position error not because the robot deviated physically from the path, but simply because it is *behind schedule*. In a desperate attempt to catch up, the actuators max out, applying violent torques that throw the system entirely off the stable path.

\begin{keyidea}{The Autonomy of Time}{autonomy_of_time}
  A system executing a motion should dictate its own progress. If the system is impeded, the reference should slow down. We must decouple the spatial geometry of the motion from its temporal execution. We achieve this by replacing chronological time $t$ with a monotonically increasing \textbf{phase variable} $s$.
\end{keyidea}

Instead of parameterizing our optimal trajectory as $\traj^*(t)$, we parameterize it as a function of the phase: $\traj^*(s)$. The phase $s \in [0, 1]$ represents the percentage of motion completed. For a walking robot, $s$ could be the forward position of the hip. For a golf swing, $s$ could be the angular position of the shoulder relative to the ball. The reference trajectory is now a purely spatial curve.

\begin{definition}{Phase Variable}{phase_variable}
  A phase variable $s(\state)$ is a continuously differentiable scalar function of the state $\state$ such that its time derivative is strictly positive along all feasible trajectories inside the nominal funnel:
  \begin{equation}
    \dot{s}(\state) > 0 \quad \text{for all } \state \in \mathcal{F}.
  \end{equation}
\end{definition}

% ══════════════════════════════════════════════
\section{Virtual Constraints}
% ══════════════════════════════════════════════

If the goal is no longer to match a time sequence but rather to conform to a spatial path defined by $s(\state)$, we can express the control objective as a set of holonomic constraints on the state. However, because these constraints are not mechanical linkages but rather enforced by software, they are termed \textbf{Virtual Constraints}.

Let our control objective be to drive the state $\state$ to the invariant manifold $\mathcal{Z}$ defined by:
\begin{equation}
  \mathcal{Z} = \{ \state \mid \state - \traj^*(s(\state)) = \bm{0} \}.
\end{equation}

By using partial feedback linearization (Chapter 5), we can design a synthetic input $\bm{v}_a$ that drives the actuated degrees of freedom $\bm{q}_a$ onto their nominal surfaces proportional to the phase:
\begin{equation}
  y = \bm{q}_a - \traj^*_a(s(\state)) \to 0.
\end{equation}
As $y \to 0$, the system converges toward the virtual constraints. Crucially, the rate at which it progresses along the path is entirely governed by $\dot{s}$, which is itself determined by the current momentum and state of the system, not a stopwatch.

% ══════════════════════════════════════════════
\section{Hybrid Zero Dynamics}
% ══════════════════════════════════════════════

When the virtual constraints are perfectly satisfied ($y = 0$), the entire complexity of the $n$-dimensional state space collapses down into a highly reduced dynamical system that evolves strictly along the manifold $\mathcal{Z}$. This is the \textbf{Zero Dynamics} of the system.

For systems that undergo discrete impacts (a foot hitting the ground, or a golf club striking a ball), the system evolves via continuous differential equations until an impact surface is triggered, followed by an instantaneous jump map (a reset) dictated by conservation of momentum.

\begin{principle}{Hybrid Zero Dynamics (HZD)}{hzd}
  A control architecture achieves Hybrid Zero Dynamics if the invariant manifold $\mathcal{Z}$ (the zero dynamics surface) is invariant not just under the continuous flow of the differential equations, but also invariant under the discrete impact map. That is, an impact starting on the manifold $\mathcal{Z}$ must immediately map to a new state that is also on $\mathcal{Z}$.
\end{principle}

By enforcing HZD, we can design rhythmic walking gaits that are astonishingly robust. If the robot trips, its phase slows down, but the virtual constraints remain active, keeping the limbs on the correct spatial path until the next footfall resets the sequence. The temporal sequence arises naturally from the interplay of gravity and momentum.

% ══════════════════════════════════════════════
\section{Application: The Golf Downswing}
% ══════════════════════════════════════════════

While the golf swing is not a periodic walking gait, the phase-variable approach perfectly encapsulates the athletic reality of the swing. The golfer's shoulder angle defines the phase $s$. The target manifold is $s = 1.0$ (the impact location). 

By structuring the control scheme computationally through Virtual Constraints, the golfer is impervious to slight hesitations at the top of the backswing. The transverse feedback LQR (Chapter 4) corrects spatial deviations from the unactuated club ($\theta_2$), while the phase variable ensures the arm actuation ($u_1$) applies the necessary parametric braking entirely in synchrony with the physical location of the arms, leading to perfectly timed passive impact dynamics regardless of chronological speed.

\chapter{Stochastic Trajectories \& Motor Variability}

\epigraph{%
  Perfection is inherently fragile. Biological optimization seeks the strongest balance between intent and reality.
}{}

\vspace{0.5cm}

% ══════════════════════════════════════════════
\section{Signal-Dependent Noise}
% ══════════════════════════════════════════════

When a roboticist specifies a torque of $50\text{ Nm}$, the motor delivers $50\text{ Nm}$ with a tiny, constant variance. Biological actuators do not work this way. Human muscles generate force by recruiting discrete motor units. As higher forces are required, larger, more fast-twitch oriented motor units are recruited. This physiological reality manifests as a profound control challenge: \textbf{motor noise is proportional to the command signal}.

\begin{definition}{Signal-Dependent Noise (SDN)}{sdn}
  A control input $\control(t)$ applied by a biological system is actually realized as a stochastic variable $\tilde{\control}(t)$:
  \begin{equation}
    \tilde{\control}(t) = \control(t) + \bm{W} \control(t) \cdot \bm{\epsilon}(t)
  \end{equation}
  where $\bm{W}$ is a scaling matrix representing the coefficient of variation, and $\bm{\epsilon}(t) \sim \mathcal{N}(0, I)$ is white noise. That is, the variance of the applied force grows squarely with the magnitude of the ordered force: $\text{Var}(\tilde{u}) = c u^2$.
\end{definition}

This simple reality crumbles deterministic optimal control. If you compute an optimal trajectory that requires maximum muscular exertion, the trajectory will mathematically possess minimum variance only if noise is constant. Under SDN, maximum exertion guarantees maximum chaotic noise.

% ══════════════════════════════════════════════
\section{Minimum Variance Theory of Human Movement}
% ══════════════════════════════════════════════

Why do humans move in smooth, bell-shaped velocity profiles rather than jerky "bang-bang" optimal control solutions when reaching for coffee?

Classical control attempts to explain this via "minimum jerk" theories, positing that the nervous system explicitly encodes an optimization functional to minimize $\int \dddot{x}^2 \dt$. The trajectory-first perspective offers a much deeper, physically grounded explanation championed by Harris and Wolpert: the \textbf{Minimum Variance Theory}.

\begin{keyidea}{Optimization for Repeatability}{min_variance}
  The human nervous system seeks to move in a way that minimizes the final endpoint variance evaluated across multiple trials.
\end{keyidea}

Under SDN, the only way to minimize endpoint variance (ensure you actually grasp the coffee cup smoothly) is to minimize sudden spikes in required actuator forces, because large forces inject huge amounts of uncertainty into the dynamics. Smooth trajectories inherently possess lower peak torque commands than rapid accelerations and decelerations. Smoothness is not the explicitly programmed objective; it is an emergent property of attempting to be consistently accurate under the unavoidable reality of signal-dependent noise.

% ══════════════════════════════════════════════
\section{The Speed-Accuracy Tradeoff in the Golf Swing}
% ══════════════════════════════════════════════

Fitts's Law states that the faster you try to move to a target, the less accurate you become. In golf, we see the supreme manifestation of this. The deterministic optimal control solution (Chapter 6) tells the solver to maximize input torque to achieve theoretically maximal clubhead speed. 

However, if we introduce the Stochastic OCP (via Covariance Steering) from Chapter 6 and apply the realistic premise of SDN, the optimizer faces a profound tradeoff.

Applying maximum $u_1$ torque during the downswing injects enormous variance into the kinematic state. The unactuated club ($\theta_2$) begins whipping through its zero dynamics, but its exact phase progression is now highly uncertain due to the noise flooding the arms. 

To ensure the clubface squares up at the moving goal, the stochastic optimizer will \emph{deliberately back off} the peak requested torque. It trades a deterministic $5$ mph gain in clubhead speed for a massive reduction in angular variance at the target manifold. The optimal biological swing is mathematically proven to be a sub-maximal physical exertion.

\chapter{Learning to Move}

\epigraph{%
  The first swing establishes the manifold. The ten-thousandth swing perfects the funnel.
}{}

\vspace{0.5cm}

% ══════════════════════════════════════════════
\section{Iterative Learning Control (ILC)}
% ══════════════════════════════════════════════

Trajectory optimization (Chapter 6) computes a reference motion strictly from a mathematical model. However, dynamic models of humans—and even robots—are invariably flawed. Frictional forces are unmodeled, inertias are estimated, and actuators possess hidden dynamics. If you execute a computed reference $\traj^*(t)$ open-loop, the physical system will likely fail to strike the moving target due to these discrepancies.

\textbf{Iterative Learning Control (ILC)} resolves this by leveraging the repetitive nature of tasks like a golf swing. ILC exploits the premise that while models are inaccurate, \emph{the unmodeled dynamics are highly deterministic across repeated trials}. 

If a robot under-swings by $5$ cm on the first attempt, the controller records this terminal error and updates the feedforward command $\control_{\text{ff}}$ for the \emph{second} attempt. Over repeated trials, the system learns the inverse dynamics of the true physical plant without ever explicitly identifying the model parameters. 

\begin{definition}{ILC Update Law}{ilc_law}
  Let $j$ denote the trial index. The feedforward command for trial $j+1$ is updated based on the command and the \emph{transverse error} $\bm{\xi}$ from trial $j$:
  \begin{equation}
    \control_j(t) = \control_{\text{ff}, j}(t) + \bm{K}_\perp(t) \bm{\xi}_j(t)
  \end{equation}
  \begin{equation}
    \control_{\text{ff}, j+1}(t) = \control_{\text{ff}, j}(t) + \bm{L}(t) \bm{\xi}_j(t)
  \end{equation}
  where $\bm{L}(t)$ is a learning filter. This updates the nominal trajectory to iteratively cancel out deterministic disturbances.
\end{definition}

% ══════════════════════════════════════════════
\section{Policy Search and Reinforcement Learning}
% ══════════════════════════════════════════════

While ILC updates the feedforward commands to zero-out deterministic tracking errors, \textbf{Policy Search} algorithms optimize the \emph{shape of the funnel} itself.

Instead of computing the transverse LQR gains $\bm{K}_\perp(t)$ from an analytical Riccati equation (Chapter 4), we parameterize the controller as a policy $\pi_{\bm{\theta}}(\state)$ and search the parameter space $\bm{\theta}$ to minimize the variance at the moving goal under realistic noise.

Reinforcement Learning (RL), specifically policy gradient methods, allows the agent to sample thousands of golf swings in simulation, continuously updating its policy parameters to maximize the expected reward (clubhead speed and squareness at impact). A heavily trained RL agent organically converges on the very principles we mathematically derived in earlier chapters:
\begin{enumerate}
    \item It abandons rigid time-tracking and autonomously discovers phase-variable mapping.
    \item It exploits the underactuated double-pendulum whipping motion to gain speed.
    \item It artificially widens its "funnel" during the transition phase, realizing that overly tight feedback control here merely amplifies motor noise and spoils the impact constraints.
\end{enumerate}

% ══════════════════════════════════════════════
\section{Trajectory Libraries}
% ══════════════════════════════════════════════

A singular optimal trajectory $\traj^*$ and its corresponding funnel $\mathcal{F}$ are highly specific to a single objective. What if the golfer wishes to fade the ball, hit a draw, or execute a punch shot out of the rough?

Instead of re-running a massive NLP optimization for every scenario, biological systems and advanced controllers maintain \textbf{Trajectory Libraries}. A library consists of several distinct, pre-computed optimal trajectories and their corresponding local funnels.

When faced with a new task (e.g., striking the ball on an uneven lie), the system searches its library for the closest matching funnel. By interpolating between established funnels using techniques like LQR-Trees~\cite{Tedrake2009} or simple spatial blending, the system can instantly generate a stable, feasible motion without any online optimization.

\chapter{Case Study: The Complete Golf Swing}

\epigraph{%
  The physics are invariant. What separates the professional is their mastery over the geometry of the resultant funnel.
}{}

\vspace{0.5cm}

% ══════════════════════════════════════════════
\section{Unifying the "Control is Motion" Framework}
% ══════════════════════════════════════════════

Throughout this text, we have systematically deconstructed the classical control paradigm and rebuilt it to handle tasks defined uniquely by motion. We conclude by applying our holistic framework to a 15-DOF musculoskeletal model of the human golf swing in order to formally synthesize the mathematical structure of athletic perfection.

The objective is clear: maximize negative clubhead velocity at the precise spatial intersection with the golf ball, while maintaining face angle orthogonality. The exact time interval $t_f$ in which this occurs is irrelevant; only the geometric state at impact matters. Furthermore, the commanded motion must not exceed the structural limits of biological torque, and the trajectory must be highly robust to physiological signal-dependent noise (SDN).

\subsection{Phase 1: Defining the Geometry and Optimization}

The system is heavily underactuated. The wrists are treated as passive pin joints, delegating the final propulsive phase entirely to inertial coupling. 

We formulate a \textbf{Stochastic Moving Target Optimal Control Problem} (Chapter 6). We deploy Direct Collocation to optimize the trajectory $\traj^*(s)$ parameterized over a spatial phase variable $s$ representing the shoulder angle (Chapter 8). The integral cost function strictly penalizes signal variance (Chapter 9) by incorporating local covariance propagation. 

The collocation solver successfully navigates the $\Bmat^\perp$ annihilator constraint (Chapter 5) and produces an optimal, smooth backswing. It discovers that a massive torque inversion (braking the torso) right before impact causes the zero dynamics of the passive wrist to wildly accelerate, perfectly squaring the face precisely as the clubhead spatial coordinate intersects the ball's location.

\subsection{Phase 2: Funnel Synthesis and Robustness}

With $\traj^*(s)$ computed, we lock down the trajectory using \textbf{Transverse Linearization} (Chapter 4). We ignore errors along the tangent of the path (since chronological timing is irrelevant) and construct a time-varying LQR feedback matrix to regulate only the transverse deviations. 

Finally, we compute the verified bounds of this stability region using \textbf{Sums-of-Squares Programming} (Chapter 7). The SDP calculates a rigorous boundary function $V(\state, t) = 1$ that proves the system is confined to an invariant funnel. 

\begin{keyidea}{The Architecture of the Perfect Swing}{perfect_swing}
  The final computed funnel reveals the mathematical secret of the professional golfer. During the backswing, the funnel is immensely wide; the motion is forgiving, and vast transverse deviations are easily tolerated without violating the invariant boundary.
  
  As the swing approaches the transition and the explosive downswing, the funnel aggressively tightens. The stochastic optimization has selected a trajectory whose passive unactuated dynamics are so profoundly stable in the final $0.15$ seconds that the funnel diameter at impact shrinks to virtually zero. Any small deviation injected by muscular noise is passively swallowed by the immense centrifugal geometry of the zero dynamics. The golfer does not aggressively steer the club to the ball; they simply set the initial conditions, and the physics irresistibly pull the clubface through the target manifold.
\end{keyidea}

Through this book, we hope we have proven that the thermostat has its place, but when the objective is a sweeping, dynamic motion cutting through the physical world, "Control is Motion" provides the only mathematically truthful language.


\bibliographystyle{plain}
\bibliography{../geometry_of_motion}


\backmatter
\chapter{Glossary of Important Concepts}
\markboth{Glossary of Important Concepts}{Glossary of Important Concepts}

These definitions distinguish geometric objects, physical models, and the guarantees obtained from a calculation. A definition does not establish that a particular golf model satisfies its assumptions.

\begin{description}[style=nextline,leftmargin=0pt,font=\bfseries,itemsep=0.5em]
\item[Articulated Body Algorithm]
A recursive forward-dynamics algorithm for a rigid-body tree. With bounded joint dimension, its arithmetic cost scales linearly with the number of bodies. Closed-loop constraints and contact require additional treatment; the tree algorithm alone does not solve every constrained contact problem.

\item[Configuration Space]
The set of configurations allowed by a model, including positions, orientations, and declared geometric constraints. It is a smooth manifold when suitable regularity conditions hold; redundant coordinates and constraint singularities require care. Velocities are additional state information.

\item[Contraction Analysis]
Analysis of convergence between trajectories under stated common inputs or closed-loop dynamics, in a specified metric and domain. A differential contraction inequality needs regularity, metric bounds, and domain/existence conditions to imply the corresponding finite-distance result. Contraction does not by itself prove input feasibility or a desired contact outcome.

\item[Control Contraction Metric (CCM)]
A positive-definite differential metric satisfying control-dependent contraction conditions that support feedback construction for feasible reference trajectories. The result depends on the theorem's domain, regularity, metric bounds, and controller construction. Actuator saturation and state constraints require explicit treatment.

\item[Degrees of Freedom (DOF)]
The number of locally independent coordinates needed to specify a configuration. Independent constraints can reduce it; redundant parameters can exceed it. The number of actuators and the dimension of the full state space are different quantities.

\item[Drift]
The evolution of the declared effective plant when the declared control channel is zero. It includes the configuration-velocity kinematics and all retained uncontrolled dynamics, such as gravity, inertial coupling, passive elasticity, and damping. Its meaning depends on the chosen state, control channel, and environment; zero control does not mean zero acceleration or zero velocity.

\item[Drift-Control Ratio (DCR)]
A model-based comparison of drift and bounded control authority in a declared acceleration or task-projected space:
\begin{equation*}
\operatorname{DCR}_{W,\mathcal U}(x)=
\frac{\|W a_{\mathrm{drift}}(x)\|_2}
{\sup_{u\in\mathcal U(x)}\|W B_a(x)u\|_2+\varepsilon} .
\end{equation*}
State the projection/scaling $W$, admissible input set $\mathcal U$, acceleration input map $B_a$, and regularizer $\varepsilon$, with compatible units. This scalar magnitude ratio omits directional alignment and does not by itself establish controllability, cancellation feasibility, or a universal coaching threshold.

\item[Euler Angles]
Three parameters describing an orientation using a declared sequence of rotations. Different sequences have different singular sets. Angle rates are related to angular velocity by a configuration-dependent map and are not generally its Cartesian components.

\item[Exponential Map]
For a Lie group, the map from a Lie-algebra element to the group element reached at parameter one along its generated one-parameter subgroup; for a matrix Lie group it is the matrix exponential. It need not be globally one-to-one. It is not general integration of an arbitrary time-varying velocity. The Riemannian exponential defined by geodesics is a related but distinct construction requiring a connection or metric.

\item[Flow]
The solution map of an ordinary differential equation where solutions exist uniquely. For an autonomous smooth vector field it forms a local one-parameter family with the composition property. Time-dependent systems use a two-time solution map; hybrid resets need not belong to a single smooth flow.

\item[Gimbal Lock]
A singularity of an Euler-angle rate-to-angular-velocity map. The parameterization loses local rank; the rigid body's physical rotational freedom does not disappear. Quaternions or overlapping coordinate charts handle orientation differently and have their own representation conventions.

\item[Jacobian]
The derivative of a map expressed in coordinates. For a task map $y=h(q)$, $\dot y=J(q)\dot q$. Acceleration also contains $\dot J\dot q$. Spatial and body Jacobians use different frame conventions, which must match the represented velocities and forces.

\item[Kinematics]
Description of position, orientation, velocity, and acceleration relationships without using forces to determine the motion. A kinematically allowed path is not automatically dynamically feasible under available actuation.

\item[Lagrangian Mechanics]
A formulation using a Lagrangian, often $L=T-U$, together with applied generalized forces and any constraints. For unconstrained coordinates, $d(\partial L/\partial\dot q)/dt-\partial L/\partial q=Q$. Dissipation, actuation, nonconservative loads, and contact must be specified; kinetic and potential energy alone do not determine a forced system's motion.

\item[Lie Algebra]
A vector space equipped with a bilinear antisymmetric bracket satisfying the Jacobi identity. For a Lie group, it is the tangent space at the identity with the induced bracket. Brackets of smooth vector fields are another important example; accessibility arguments require the appropriate system and theorem assumptions.

\item[Lie Group]
A smooth manifold that is also a group, with smooth multiplication and inversion. Examples include the rotation group $SO(3)$ and rigid-transformation group $SE(3)$. A choice of coordinates represents these objects without changing the group itself.

\item[Lyapunov Stability]
For an equilibrium, the property that sufficiently small initial perturbations remain small for all future times of the stated system. It does not require return to equilibrium. Asymptotic stability additionally requires attraction; orbital and finite-horizon stability statements refer to different targets and must be defined separately.

\item[Manifold]
A space locally described by Euclidean coordinates, with the usual topological regularity conditions. A smooth manifold also has smoothly compatible charts. Its dimension is local and does not depend on how many coordinates are used to embed or redundantly represent it.

\item[Screw Axis]
For a rigid-body twist with nonzero angular velocity, an instantaneous axis and pitch describing rotation together with translation along that axis. Pure translation is a limiting case, often represented as an axis at infinity; a finite unique rotational axis is then unavailable.

\item[State Space]
The set of variables sufficient to determine future evolution under specified inputs and a model. Mechanical models often include configuration and velocity or momentum, and may additionally include actuator states, flexible modes, estimators, or discrete contact modes.

\item[Tangent Bundle]
The union of tangent spaces $T_qQ$ over configurations $q\in Q$, with its bundle and smooth-manifold structure. For an $n$-dimensional smooth configuration manifold, $TQ$ has dimension $2n$. It represents configuration-velocity states for a basic mechanical model; momentum states instead lie in the cotangent bundle $T^*Q$.

\item[Tangent Space]
The vector space of instantaneous derivatives of smooth curves through a point on a manifold. Its vectors are local velocities, not finite displacements between arbitrary configurations. Applied differential constraints can restrict the admissible velocities further.

\item[Twist]
A six-component representation of an instantaneous rigid-body velocity field, combining angular and linear terms with a declared reference point and frame. In one frame, $v(r)=v_0+\omega\times(r-r_0)$. Spatial and body representations are related by the appropriate adjoint transformation; the linear term must not be confused with the velocity of an unspecified point.

\item[Wrench]
The dual force representation consisting of a moment about a declared reference point and a resultant force in a declared frame. With matching twist conventions, mechanical power is $\tau\cdot\omega+F\cdot v_0$. Changing point or frame requires the corresponding moment shift and dual transformation so that power is unchanged.

\item[Zero-Torque Counterfactual (ZTCF) Family]
Model interventions setting the declared applied generalized-control channel to zero while preserving the declared effective plant. Pointwise samples, stitched traces, forward trajectories, and branched trajectories answer different questions. A stitched pointwise acceleration trace is not automatically a feasible unforced trajectory.

\item[Zero-Velocity Counterfactual (ZVCF)]
The instantaneous acceleration after the declared velocity and control intervention:
\begin{equation*}
a_{\mathrm{ZVCF}}(q,z)=-M(q)^{-1}h(q,0,z_0;p) .
\end{equation*}
Specify the auxiliary-state reset $z_0$, retained loads, parameters, and contact mode. This separates the retained configuration-dependent loads under that intervention; it is not necessarily gravity alone, and it is not a forward simulation of a stationary body.
\end{description}

\chapter{Further Reading and References}
\markboth{Further Reading and References}{Further Reading and References}
\begingroup
\raggedbottom

The following books develop complementary parts of the framework used here. Choose a source for the particular claim or calculation being investigated, and consult its assumptions, notation, edition, and errata. A robotics or control theorem does not become evidence about human golf performance merely because it can be applied to a simplified swing model.

\section*{Geometry, Kinematics, and Mechanical Models}

\begin{itemize}
\item \textbf{A Mathematical Introduction to Robotic Manipulation}, Murray, Li, and Sastry \cite{Murray1994}, treats rigid-body motion, manipulator and hand kinematics, dynamics, and nonholonomic motion planning. The \href{https://www.cds.caltech.edu/~murray/mlswiki/index.php?title=Main_Page}{authors' companion site} provides chapter information and corrections. Use it to connect Lie groups and Jacobians to explicitly defined frames and mechanical systems.
\item \textbf{Geometric Control of Mechanical Systems}, Bullo and Lewis \cite{BulloLewis2004}, develops differential-geometric modeling and control of mechanical systems. Its \href{https://motion.me.ucsb.edu/book-gcms/}{companion page} includes contents, selected material, examples, and errata. This is a suitable route into connections, mechanical control structures, and the hypotheses behind geometric analysis.
\item \textbf{Robot Modeling and Control}, Spong, Hutchinson, and Vidyasagar \cite{Spong2006}, covers kinematics, Jacobians, trajectory planning, dynamics, joint and multivariable control, force control, and vision-based methods. The \href{https://bcs.wiley.com/he-bcs/Books?action=contents&bcsId=2888&itemId=0471649902}{publisher's chapter listing} identifies the scope of the cited edition. Its robot models still require adaptation and validation before supporting a biomechanical interpretation.
\end{itemize}

\section*{Computational Rigid-Body Dynamics}

\begin{itemize}
\item \textbf{Rigid Body Dynamics Algorithms}, Featherstone \cite{Featherstone2008}, develops spatial-vector conventions and recursive dynamics algorithms. The \href{https://royfeatherstone.org/teaching/index.html}{author's teaching materials} and \href{https://royfeatherstone.org/}{software resources} support implementation. Distinguish the linear-cost articulated-body calculation for a tree from the additional work needed for closed loops, constraints, and contacts.
\end{itemize}

\begin{samepage}
\section*{Nonlinear Stability and Feedback}

\begin{itemize}
\item \textbf{Applied Nonlinear Control}, Slotine and Li \cite{SlotineLi1991}, develops nonlinear stability and feedback methods. \href{https://www.mit.edu/~nsl/videos/lectures.html}{Slotine's accompanying course materials} identify the relevant book sections and separate the later contraction lectures and papers. Check the domain and model conditions of a stability argument before extending it to saturated or hybrid systems.
\item \textbf{Contraction Theory for Dynamical Systems}, Bullo \cite{BulloContraction}, develops contractivity and incremental stability using specified norms and metrics. The \href{https://fbullo.github.io/ctds}{author's book page} provides version information and materials. The bibliography cites the 2024 edition; later revisions may reorganize or extend the treatment. Contraction claims must retain their metric, input, domain, and existence assumptions.
\end{itemize}

The chapter-level primary references supply more specialized results for transverse dynamics, phase constraints, nonlinear funnels, stochastic models, and learning. Keep mathematical derivations, numerical verification, model calibration, and experimental validation identifiable when combining those sources.

\nocite{LohmillerSlotine1998, ManchesterSlotine2017, Khatib1987, Shiriaev2010}

\end{samepage}
\clearpage
\endgroup

\printindex
\end{document}
